\clearpage
\section{分解・補償・成分ノルム}
\label{sec:construction}
本付録では有界sourceの分解から予測可能補償と正準成分のノルム評価までを構成する。
\subsection{積分の初期値と停止}

elementary積分は区間 $(0,t]$ の増分の有限和で定めるため、時刻零で零となる。
完成マルチンゲール積分も時刻零で零であり、有限変動積分はStieltjes積分の一様極限として
時刻零で零となる。局所化した各座標の利得との区別不能性を時刻零で評価すると、
一般積分表示に属する利得 $X$ は $X_0=0$ a.s.を満たす。

分解 $X=M+A$ の各成分には定数シフトの自由度がある。
$M-M_0$ と $A+M_0$ に中心化すれば利得は変わらず、
$X_0=0$ の下で両成分の初期値を零にできる。以下のcostはこの零始点の成分を用いる。

有限停止時刻 $\tau$ に対し、閉停止を $X^\tau_t=X_{t\wedge\tau}$ と書く。
strict prefixは $t<\tau$ で $X_t$、$t\geq\tau$ で $X_{\tau-}$ とし、
$\tau=0$ では $X_{0-}=X_0$ と約束する。このとき
\[
 X^\tau=X^{\tau-}+1_{\{0<\tau\leq t\}}\Delta X_\tau.
\]
従って閉停止の全変動や成分costを評価するときは、strict prefixの評価に境界jumpを加える。
確率積分の停止に用いるのは閉停止であり、strict prefixは解析的なcostの比較に用いる。

\subsection{有界sourceの分解という基礎入力}
\label{sec:bounded-source}
以下で使う分解の入力は、Bichteler--Dellacherie定理の次の有界版である。
積分の値域閉性とは分けて、採用する基礎定理の範囲を明示する。

\begin{theorem}[有界good integratorの分解]\label{thm:bounded-source}
通常の条件の下で、$S$ は適合càdlàg過程であり、決定論的な $B<\infty$ により
$\sup_t|S_t|\leq B$ がほとんど確実に成り立つとする。
有界な予測可能elementary係数の一様収束に対し、各有限horizonでのelementary積分が
確率収束するという連続性を仮定する。このとき
\[
 S-S_0=M^S+A^S
\]
と分解できる。$M^S$ は零始点càdlàg local martingale、$A^S$ は零始点、
予測可能、càdlàg、局所有限変動の過程である。
中心化した成分は区別不能性を除いて一意である。測度は元の測度のままで、市場の仮定は不要である。
\end{theorem}
\begin{proof}
標準的な基礎定理として、Bichteler--Dellacherieの特徴付けの有界な場合を使う。
すなわち有限区間の適合càdlàg good integratorはlocal martingaleと適合càdlàg FV過程の和である。
連続性の仮定は\cite[Definition 1.1 and Theorem 1.2]{BSV}のelementary積分による条件である。
有限区間への制限と時間の再尺度化によってその定式化に帰着する。
elementary係数のほとんど確実な一様上界と全点での一様上界は、通常の条件の下で
時刻0の零集合上の係数を修正することにより、同じ判定条件を与える。
従って固定horizonで $S-S_0=N+D$ と零始点成分に分けられる。
この段階では $D$ は適合であるだけである。予測可能成分への置換には以下の議論が必要であり、
任意の分解の一意性だけでは得られない。

$N$ をこのhorizonまで真のマルチンゲールにする局所化を選び、さらに $|N|$ と
$\Var(D)$ が $r$ を初めて厳密に超える時刻で停止する。
得た閉停止を $\sigma_r$ とする。その直前までは両量とも $r$ 以下である。停止時には
\[
 |\Delta D_{\sigma_r}|
 \leq2B+|\Delta N_{\sigma_r}|
 \leq2B+r+|N_{\sigma_r}|,
 \qquad
 \Var(D^{\sigma_r})\leq2r+2B+|N_{\sigma_r}|
\]
となる。局所化した真のマルチンゲールの任意抽出により最後の上界は可積分である。
経路の有界性と有限変動性から、これらの停止は固定horizonを尽くす。
付録の\ref{sec:canonical-cost}節でノルム評価とともに明記する補償定理を
$D^{\sigma_r}$ に適用して
\[
 S^{\sigma_r}-S_0
 =\bigl(N^{\sigma_r}+D^{\sigma_r}-(D^{\sigma_r})^p\bigr)
   +(D^{\sigma_r})^p
\]
を得る。各座標でspecial分解が得られた。
共通停止上で二つの予測可能FV成分の差は零始点予測可能FV local martingaleであるから零である。
これが整合性と中心化した成分の一意性を与える。
停止マルチンゲール座標と予測可能FV座標を貼り合わせ、さらに正整数horizonで貼り合わせる。
可算個の一致と右連続性から、共通の区別不能性の事象を得る。
\end{proof}

Leanではこの有界分解を直接構成する。標本化したDoob分解の増分を
$\Delta A_j=E[\Delta S_j\mid\mathcal F_{t_j}]$、
$\Delta M_j=\Delta S_j-\Delta A_j$ とし、予測可能な符号testで変動を制御する。
共通停止、Hilbert凸化、正則化から整合的な正準極限を得る。
本稿の数学的説明では、明示した有界な特徴付けを基礎定理として採用する。
確率積分の値域の閉性は外部入力にしない。
特に後の一般good integratorの分解は、大ジャンプを除去して残差を停止した後、
この有界な入力だけを使うものであり、本証明へ循環して戻るものではない。

\subsection{Doléans--Dade--Yen 型分解}

一般 local martingaleは locally $M^2$ とは限らない。running supremumの $L^1$ boundだけでは、
first passage時の未観測 jumpが二乗可積分であるとも、決定論的に有界であるともいえない。
従って、任意の local martingaleに $M^2$ localizing sequenceがあると仮定する経路は取れない。

supplied special decompositionの局所 martingale成分を $N$ とする。最初に必要なのは、$N$ に対する
Doléans--Dade--Yen 型の分解
\[
 N=L+B
\]
である。$L$ は局所 martingaleで、その jumpは局所化後に決定論的に有界であり、
$B$ は適合局所有限変動過程でなければならない。証明は次を含む。
\begin{enumerate}
\item c\`adl\`ag pathの有限区間には、任意の閾値より大きい jumpが有限個しかないことを示す。
\item large jumpsを genuine stopping timesで時系列化し、適合 c\`adl\`ag 有限変動過程を作る。
\item その variationが可積分となる exhaustive localizationを選び、予測可能射影を取る。
\item 補償後の martingale jumpを announced-time estimateで制御し、停止との整合性を示す。
\end{enumerate}
非単調な dense skeletonの列挙順を jumpの発生順として使ってはならない。

決定論的な path core については、実数値 càdlàg path $x$、$c>0$、および
有限 horizon $T$ に対し
\[
 J_{c,T}=\{t\in(0,T]:c<|x_t-x_{t-}|\}
\]
が有限であることを用いる。$J_{c,T}$ の元を時刻順に保持する有限 step path
\[
 P_t=\sum_{s\in J_{c,T},\ s\leq t}(x_s-x_{s-})
\]
を構成すると、$P_0=0$、右連続性、左極限、各時刻での jump の exact identity
（$J_{c,T}$ の jump は保持し、それ以外は $0$）が得られる。またその累積絶対
jump 和は単調であり、任意の $0\leq u\leq v$ に対して、経路の実変動
$\operatorname{eVariationOn}$ と累積変動 $C$ の間に
\[
\operatorname{eVariationOn}(P,[u,v])
 = \operatorname{ofReal}\bigl(C(v)-C(u)\bigr)
\]
という厳密な等式が成り立つ。ここで $C(v)-C(u)$ は $(u,v]$ にある保持された
jump の絶対値の有限和である。この等式は各有限step経路について成り立つ。
さらに factorial skeleton の局所振動事象を用いると、固定閾値の最初の大きい左jumpが
時刻 $t$ より前に存在することを pathwise に可測事象として表せる。この表示から、
有限な jump 集合の infimum（空集合では horizon に clamped）は genuine stopping time
となり、非空時の attainment と earliest property が得られる。

確率過程への移行では、単なる各時刻の可測性ではなく、large-jump graph の
chronological debut を停止時刻として反復することが必要である。上の specialized
first-debut theorem は generic optional debut theoremを仮定せず、factorial skeletonの
局所振動と有限集合の infimum からこの最初の停止時刻を構成する。過去の停止時刻
後の部分を deterministic post-processing で切り出し、有理数にわたる可算合併
で pre-$T$ の successor を表し、$t>T$ では 結論 $T$ の jump を明示的に加えることで、
この構成を反復できる。得られた座標は genuine stopping times であり、active な座標は
時刻順に厳密に増加し、$T$ を含む exact jump set に属する。さらに全ての大きい left jump
時刻を coverage し、有限個を尽くした後は $\top$ となる。従って chronological stopping-time
enumeration は完了している。

時系列有限和については、停止左jump process の各座標を

\[
 J^{(n)}_t(\omega)=1_{\{\tau_n(\omega)\leq t\}}
 \bigl(N_{\tau_n(\omega)}-N_{\tau_n(\omega)-}\bigr)(\omega)
\]

（$\tau_n=\top$ では零）とし、$J=\sum_n J^{(n)}$ と定める。各座標は停止過程の
strong adaptation と停止時刻事象の indicator から strongly adapted であり、countable sum
の可測性を持つ。pathwise には enumeration が有限個の large jump を尽くした後 $\top$ に
安定するため、$J$ は前述の有限 step path と厳密に一致する。この一致から $J_0=0$、right
continuity、left limits、global/local bounded variation、および累積 jump 変動との
exact $\operatorname{eVariationOn}$ identity を得る。区間 $(0,T]$ で $|\Delta N|>c$ の jump は
コピーされ、その他の left jump は零である。また $T$ 後の constancy と、partial finite sums
の process-level eventual constancy も得られる。従って確率過程としての有限大jump processは得られた。右連続過程の level
hitting に対する停止時刻定理や、predictable interval-algebra に対する debut の構成を、
一般の optional graph の代用として扱ってはならない。

\subsection{有限大ジャンプ過程の可積分局所化}

零始点 c\`adl\`ag 局所 martingale $N$ の有限大ジャンプ過程を
\[
 J_t=\sum_{0<s\leq t\wedge T,\,|\Delta N_s|>c}\Delta N_s,
 \qquad
 V_t=\sum_{0<s\leq t\wedge T,\,|\Delta N_s|>c}|\Delta N_s|
\]
と書く。càdlàg経路は有限区間に大きさ $c>0$ を超えるjumpを有限個しか持たないため、
$J$ は有限変動であり、$T$ の後では一定である。決定論的 horizon
\[
 H_n=T+(n+1)
\]
を用い、$V$ の最初の strict passage によって
\[
 \rho_n=\min\bigl\{H_n,\inf\{s:V_s>n+1\}\bigr\}
\]
を定める。$H_n$ は cofinal であり、$\rho_n$ は variation level と horizon の双方の単調性から増加する停止時刻列である。有限変動性と $T$ 後の constancy により、各 sample path では有限段階の後に
\[
 \rho_n=H_n
\]
となる。従って $(\rho_n)$ は $\top$ に発散する genuine な localizing sequence である。

各段階では閉じた停止過程ではなく strict-prefix 過程
\[
 C^n_t=
 \begin{cases}
 J_t,&t<\rho_n,\\
 J_{\rho_n-},&t\geq\rho_n
 \end{cases}
\]
を source とする。この過程は適応 càdlàg 有限変動過程であり、停止時刻 $H_n$ までの累積変動は level $n+1$ 以下であるため、terminal cumulative variation は $L^1$ に属する。従って各 $C^n$ から、正規化された適応 càdlàg 有限変動 $L^1$ data を構成し、value truncation の定理によって predictable 有限変動極限と martingale residual を得る。source と strict-prefix の一致は $t<\rho_n$ の開いた prefix 上で厳密に保持される。

さらに $\rho_n$ を martingale localizer
$\lambda_n$ と absolute-passage localizer $\kappa_n$ で refineして
$\tau_n=\rho_n\wedge\lambda_n\wedge\kappa_n$ とする。これは有限停止時刻であり、
\[
 \operatorname{Var}(J^{\tau_n};[0,H_n])
 \le n+1+|\Delta J_{\tau_n}|.
\]
右辺の jump項は次節の評価で可積分となる。従って各 closed source $J^{\tau_n}$ からも
有限変動 $L^1$ data と value-truncation predictable projection が得られる。全ての $n$ に
対する選択を一つの family に固定する。$(\rho_n)$、$(\lambda_n)$、$(\kappa_n)$ はいずれも
a.e. increasing な localizing sequence なので、$(\tau_n)$ も a.e. increasing かつ exhaustive
である。従って $m\le n$ なら停止の吸収則により
\[
  (J^{\tau_n})^{\tau_m}\sim J^{\tau_m}
\]
を得る。この主張は全点での単調性を要求せず、process indistinguishability として述べる。
さらに $m\le n$ とする。上位の predictable projectionを $\tau_m$ で停止したものと下位の
projectionとの差は、上の closed-source identityによって local martingaleである。一方、この差は
strongly predictable、right-continuousかつ有限変動である。予測可能有限変動 local martingaleの
rigidityと、両 projectionの初期値が a.e. に零であることから、この差は indistinguishably 零である。
従って level 間の predictable projections は ordered overlap 上で compatible になった。

\subsection{有限大ジャンプを除いた residual}

元の局所 martingale $N$ から有限大ジャンプ過程を差し引いた residual を
\[
 R_t=N_t-J_t,
 \qquad
 J_t=\sum_{0<s\le t\wedge T,\,|\Delta N_s|>c}\Delta N_s
\]
と定める。left limit を用いると全時刻で
\[
 \Delta R_t=\Delta N_t-\Delta J_t
\]
が厳密に成り立つ。$0<t\le T$ では、$|\Delta N_t|>c$ のとき
$\Delta J_t=\Delta N_t$ であり、それ以外では $\Delta J_t=0$ である。従って
\[
 0<t\le T\quad\Longrightarrow\quad |\Delta R_t|\le c .
\]
さらに $t>T$ では $\Delta J_t=0$ なので $\Delta R_t=\Delta N_t$ である。
この境界は source-level の residual に限られる。$J$ の predictable projection や
compensator を差し引いた過程について同じ jump bound を主張するには、announced-time または
conditional-expectation による別の estimate が必要であり、ここでは仮定していない。closed source
と predictable projection の nested overlap に加え、貼り合わせたglobal predictable process $P$ も得た。
さらに停止列のexhaustionにより、同じ $P$ に対する $Q=J-P$ が局所martingaleであることを得た。
この $P$ のpredictable jump estimateも以下で得る。

一方、closed stop が追加する境界 jump については、局所化をさらに細かくすれば可積分性を
得られる。決定論的な $H<\infty$ に対して $\tau\leq H$ を満たす停止時刻 $\tau$ が、ある
local-martingale localizer と level $r$ の absolute strict passage の双方以下であるとする。
前者によって $N_\tau$ は true martingale の bounded optional sample と同一視できるので可積分であり、
後者と左極限の存在によって
\[
  |N_{\tau-}|\le r
\]
である。有限大ジャンプ過程 $J$ の jump は $\Delta N_\tau$ または $0$ なので、
\[
  |\Delta J_\tau|
  \le |N_\tau|+|N_{\tau-}|
  \le |N_\tau|+r .
\]
従って $\Delta J_\tau$ は可積分である。この議論には $\tau$ の announcement は不要であり、
境界 correction が predictable であるとも主張しない。この estimate は上の共通 refinement に
実際に用いられ、closed stopped slice の可積分全変動と予測可能射影を与える。

\subsection{予測可能補償過程の貼り合わせと jump bound}

各 $n$ について、$J^{\tau_n}$ の双対予測可能射影（補償過程）を $P^n$、martingale residualを
$Q^n=J^{\tau_n}-P^n$ と書く。局所構成で得た ordered overlap は
\begin{equation}
 m\leq n\quad\Longrightarrow\quad (P^n)^{\tau_m}\sim P^m
 \label{eq:compensator-ordered}
\end{equation}
である。一方、局所有限変動過程の貼り合わせが要求する形は
\begin{equation}
 (P^n)^{\tau_n\wedge\tau_m}\sim
 (P^m)^{\tau_n\wedge\tau_m}\qquad(n,m\in\mathbb N)
 \label{eq:compensator-minimum}
\end{equation}
である。式~\eqref{eq:compensator-ordered}で $m=n$ とすれば各 coordinate の自己停止が得られる。また一つの
full-measure set上で $(\tau_n)$ は増加するので、$m\leq n$ なら
$\tau_n\wedge\tau_m=\tau_m$ である。二重停止の吸収則と自然数の全順序を使えば
式~\eqref{eq:compensator-minimum}が従う。これにより ordered overlapを貼り合わせ定理のall-pairs入力へ移した。

式~\eqref{eq:compensator-minimum}から pointwise に eventually constantな局所極限を作ると、予測可能性は極限で保たれる。
compatibilityとscheduleのa.e. 増加性・exhaustionをそのまま用いることで、right continuityと
局所有限変動性がa.e. に得られる。そのpath-regularityの例外集合上で過程全体を零にすれば、
全sample pathでright-continuousかつ局所有限変動な予測可能version $P$ を得る。
compatibilityや $(\tau_n)$ 自体をpointwiseなdataへ置き換える必要はない。ただし各 $P^n_0=0$ は
a.e. の等式なので、貼り合わせだけでは $P_0=0$ も
a.e. にしか得られない。$\{P_0\neq0\}$ は時刻零で可測な零集合であり、そこでも過程全体を零にして
$P_0=0$ をpathwiseに正規化する。この二度目の正規化はpredictability、right continuity、
局所有限変動性、および各coordinateとのstopped indistinguishabilityを保つ。この構成により、同じ
versionの $P$ についてこれら全ての性質が得られた。

closed sourcesにも式~\eqref{eq:compensator-minimum}と同じminimum-stop compatibilityがあり、差を取ることで $(Q^n)$ にも
同じcompatibilityを得た。各 $Q^n$ は c\`adl\`ag true martingaleなので、martingale
貼り合わせ定理から、一つのstrongly adaptedな c\`adl\`ag local martingale $\widetilde Q$ を得る。この
$\widetilde Q$ は全ての $n$ について、$\tau_n$ で停止すると $Q^n$ とindistinguishableである。各
$\tau_n$ で
\[
 J^{\tau_n}\sim \widetilde Q^{\tau_n}+P^{\tau_n}
\]
が成り立つ。可算個のstopped agreementとexhaustionの共通full-measure set上では、任意の時刻 $t$ に
対し $t\leq\tau_n$ となる $n$ を選べる。従って $J\sim\widetilde Q+P$ である。正規化後の同じ
representative上で $Q=J-P$ と定めれば $Q\sim\widetilde Q$ である。従って $Q$ は
局所 martingaleであると同時に、適合 c\`adl\`ag 局所有限変動過程であり、$Q_0=0$ である。
この残差を局所 martingaleと示す際には、local propertyの時刻零indicatorを直接展開せず、
true-martingale coordinatesの貼り合わせとindistinguishabilityによる移送を用いる。

$L=N-Q=R+P$ と置けば $N=L+Q$ であり、$L$ も c\`adl\`ag 局所 martingaleである。しかし
$|\Delta R|\leq c$ に加え、次の補償過程の評価を用いる。

\begin{lemma}[補償過程の jump estimate]
\label{lem:compensator-jump-bound}
固定した $c>0$ と $T<\infty$ に対し、一つの full-measure set上で
\[
 0<t\leq T\quad\Longrightarrow\quad
 |\Delta P_t|\leq c,
 \qquad |\Delta L_t|\leq 2c .
\]
\end{lemma}

\begin{proof}
$R=L+(-P)$ を時刻 $T$ で停止する。予測可能 c\`adl\`ag 有限変動過程 $P^T$ の全ての非零jumpは、
累積変動の有理level crossingから作る可算個のannounced predictable stopping graphs
$(\sigma_k)$ で覆われる。完成した各時刻は $T+1$ 以下である。局所化と条件付き期待値の評価を
$L^T$ と $-P^T$ に適用すると、元の局所martingaleのjump全体の可積分性を仮定せずに
\[
 |\Delta (P^T)_{\sigma_k}|
 \leq
 \mathbb E[c\mid\mathcal F_{\sigma_k-}]=c
 \quad\text{a.s.}
\]
を得る。入力の評価は $|\Delta(R^T)_t|\leq c$ であり、$t>T$ では停止済み過程のjumpが零である。
可算個の零集合をまとめ、$(\sigma_k)$ のcoverageを用いれば、$P$ のjumpが零である時刻も含めて
最初の結論が従う。最後に $\Delta L=\Delta R+\Delta P$ と三角不等式を用いる。
$T=0$ は空な時間区間として分け、enumerationのdummy value $T+1$ では停止過程のjumpが零であることを使う。
\end{proof}

従って補償後に必要なestimateは、$Q=J-P$ 自身に対する $c$ boundではなく、まず
$|\Delta P|\leq c$、次にbounded-jump martingale $L=N-Q$ に対する $2c$ boundである。
これにより、元の零始点c\`adl\`ag local martingale $N$、$c>0$、有限horizon $T$、usual conditionsから
大jumpの局所化列と補償過程を構成し、同じ $L,Q$ に対する $N=L+Q$ を得る。
両成分は零始点c\`adl\`ag local martingaleであり、$Q$ は適合局所有限変動過程である。
$Q$ の予測可能性は主張しない。任意のpathwise stop $\rho\leq T$ についても、同じ零集合の外で
$|\Delta(L^\rho)_t|\leq2c$ が全時刻で成り立つ。これはjump評価の停止安定性である。
得られる分解は固定horizonに依存する。

\begin{lemma}[Horizon間の整合性]
各horizon $T$ で得た成分を $P^{(T)},Q^{(T)},L^{(T)}=N-Q^{(T)}$ と書く。
$U\leq T$ なら
\[
 (J^{c,T})^U=J^{c,U}
\]
が経路ごとに成り立ち、一つのfull-measure set上で全時刻について
\[
 (P^{(T)})^U=(P^{(U)})^U,\qquad
 (Q^{(T)})^U=(Q^{(U)})^U,\qquad
 (L^{(T)})^U=(L^{(U)})^U
\]
が成り立つ。各horizon内部のlocalizing sequenceを共通に選ぶ必要はない。
\end{lemma}

\begin{proof}
$(0,U]$ の大jump時刻集合は $(0,T]$ のものの部分集合である。
後者にのみ属する時刻は $U$ より大きく、時刻 $t\wedge U$ での有限和には寄与しない。
これにより最初の等式が従う。残りについては
\[
 A=(P^{(T)})^U-(P^{(U)})^U
   =(Q^{(U)})^U-(Q^{(T)})^U
\]
を考える。左辺は予測可能で、有限horizonで停止したため全時間軸で有限変動を持つ。
右辺はc\`adl\`ag local martingaleの差である。従ってpredictable FV local-martingale rigidityから
$A$ は初期値とindistinguishableである。初期値は零であるため、$P$ の一致が従い、同じ式から $Q$、
さらに $L=N-Q$ の一致を得る。$U=0$ は零初期値から直接従う。
\end{proof}

元の零始点local martingaleから各 $T_n=n+1$ の分解を構成し、上の補題をその同じ成分へ適用する。
可算交叉により、全 $m\leq n$ と全時刻でのoverlapを一つのfull-measure set上に揃えられる。
また同じ $P^{(T_n)}$ のjump boundも全 $n$ で共通のfull-measure set上に揃えられる。

\begin{lemma}[予測可能成分の全時刻への貼り合わせ]
元の零始点local martingaleから上記のcoherent familyと、
零始点の予測可能な右連続局所有限変動過程 $P$ を構成でき、
一つのfull-measure set上で全 $n,t$ に対して
$P_{t\wedge T_n}=P^{(T_n)}_{t\wedge T_n}$ が成り立つ。
さらに一つのfull-measure set上で全 $t$ に対し $|\Delta P_t|\leq c$ である。
\end{lemma}

\begin{proof}
各 $P^{(T_n)}$ を $T_n$ で停止して全時間軸で有限変動の座標とする。
二つのindexの大小に応じてoverlapを適用すれば、両座標は $T_n\wedge T_m$ で停止した後に一致する。
決定論的停止時刻列 $T_n$ は単調で無限大へ発散するため、予測可能な有限変動過程の貼り合わせを
適用できる。得られた過程の初期値はほとんど確実に零である。その例外集合はusual conditionsにより
時刻零で可測なので、その集合上で全pathを零に修正する。
可算交叉で全horizonの一致を同時に確保する。
任意の $t$ に対し $t\leq T_n$ となる $n$ を取り、停止前および停止時刻のleft jump一致から
$|\Delta P_t|=|\Delta P^{(T_n)}_t|\leq c$ を得る。
\end{proof}

\begin{lemma}[Local-martingale座標の共通細分]
零始点c\`adl\`ag local martingales $X^n$ が、ほとんど確実に単調でexhaustiveな停止時刻列 $(\tau_n)$ に対し、
$(X^n)^{\tau_n\wedge\tau_m}\sim(X^m)^{\tau_n\wedge\tau_m}$ を満たすとする。
このときc\`adl\`ag local martingale $Y$ が存在し、全 $n$ で
$Y^{\tau_n}\sim(X^n)^{\tau_n}$ である。
\end{lemma}

\begin{proof}
$X^n$ のlocalizing sequenceを $(\eta^n_j)_j$ とする。
可算な列の抽出により、$\sigma_n=\tau_n\wedge\eta^n_{k(n)}$ がほとんど確実に無限大へ発散するように
$k(n)$ を選べる。$\rho_n=\inf_{j\geq n}\sigma_j$ と置くと、filtrationの右連続性により停止時刻であり、
単調に無限大へ発散し、$\rho_n\leq\tau_n\wedge\eta^n_{k(n)}$ である。
零初期値はlocalizing stopに付く時刻零のindicatorを除去するため、optional stoppingから
$(X^n)^{\rho_n}$ はtrue martingaleである。
元のcompatible familyの点ごとの極限を $W$ とすると、$W^{\rho_n}\sim(X^n)^{\rho_n}$ である。
これにより細分した座標もcompatibleであり、true-martingale gluingを適用できる。
得られた $Y$ は全 $\rho_n$ で $W$ と一致するため、exhaustionから $Y\sim W$ である。
元の $\tau_n$ で停止し直すと所望の一致を得る。
\end{proof}

\begin{proposition}[全時刻のDoléans--Dade--Yen分解]
零始点c\`adl\`ag local martingale $N$ と $c>0$ に対し、零始点c\`adl\`ag local martingales
$L,Q$ が存在し、$N=L+Q$、$Q$ は局所有限変動であり、一つのfull-measure set上で全時刻について
$|\Delta L_t|\leq2c$ が成り立つ。
\end{proposition}

\begin{proof}
同じhorizon familyの $Q^{(T_n)}$ に共通細分とgluingを適用する。
各horizonでの一致からlocal BVと初期値零がほとんど確実に従う。
例外集合は時刻零で可測であり、その上でpathを零に修正して全pathのlocal BVと零初期値を得る。
同じfamilyから貼った $P,Q$ は、全horizonの共通零集合上で
$(J^{c,T_n})^{T_n}=(Q+P)^{T_n}$ を満たす。
$L=N-Q$ と置けば、$t\leq T_n$ に対して
$\Delta L_t=\Delta(N-J^{c,T_n})_t+\Delta P_t$ である。
任意の正の $t$ を含むhorizonを選び、それぞれのbound $c$ を加える。
時刻零のjumpは零である。$L$ の残りの性質はlocal martingalesの差から従う。
\end{proof}

この同じ大域 $L$ に次のquadratic constructionを適用する。
有限horizonの $L^T$ への適用も同じ構成の別の入力である。

\begin{lemma}[Bounded-jump local martingaleのquadratic局所化]
零始点c\`adl\`ag local martingale $M$ が、各整数horizon $n+1$ で一つのfull-measure set上の
jump bound $|\Delta M_t|\leq b_n$ を持つとする。ここで $b_n\geq0$ は定数である。
$M$ のlocal-martingale localizing sequenceを $(\eta_n)$ とし、
\[
 \theta_n=\inf\{t\geq0:|M_t|>n+1\},\qquad
 \rho_n=\eta_n\wedge\theta_n\wedge(n+1)
\]
と置く。この列はexhaustiveであり、各 $M^{\rho_n}$ はtrue martingaleで、全時刻共通に
\[
 |M^{\rho_n}_t|\leq n+1+b_n\quad\text{a.s.}
\]
が成り立つ。従って各座標の終端値は $L^2$ に属する。
\end{lemma}

\begin{proof}
absolute passageをhorizonで切った列と $(\eta_n)$ はともにlocalizing sequenceであり、その最小値も
同じ性質を持つ。零初期値により、$\eta_n=0$ のpathでもindicator付きlocalizing stopと通常のclosed stopが
一致する。従って $M^{\eta_n}$ はtrue martingaleである。$\rho_n\leq\eta_n$ でのbounded optional stoppingから
$M^{\rho_n}$ のtrue-martingale性を得る。passage直前の絶対値は $n+1$ 以下であり、境界で加わるjumpは
$b_n$ 以下であるため、表示した全時刻評価が従う。定数の $L^2$ majorantにより終端値の二乗可積分性を得る。
\end{proof}

各座標に有限horizon quadratic constructionを適用し、minimum-stop compatibilityとgluingを用いて
一つのglobal optional quadratic variation $V$ を得る。同じ $V$ のlocalizer rootと座標のterminal rootは
a.e. に一致し、その期待値costも一致する。各座標でDavis評価を適用すると
\[
 \mathbb E\sup_{t\leq n+1}|M^{\rho_n}_t|
 \leq 6\,\mathbb E\sqrt{V_{\rho_n}}
\]
を得る。ここでsupremumはc\`adl\`ag性により可算な時刻集合から検出する。

固定horizon DDY分解では $M=L^T$、$b_n=2c$ としてこの全構成を適用する。
元の零始点local martingaleからDDY成分を選び、その同じ $L^T$ に対するquadratic/Davis評価まで得られる。
この $V$ は $L^T$ のquadratic variationであり、元の $N$ のものではない。また各座標の評価から
unstopped $L^T$ のglobal $H^1$ 可積分性は主張しない。

\subsection{一般 quadratic variation と exact \texorpdfstring{$j^1$}{j1}}

bounded-jump martingaleに対しては、有限 horizonの $M^2$ quadratic process、停止整合性、
局所化列上の貼り合わせ、一意性、terminal root、および定数6の一方向 Davis estimateが得られている。
local martingaleへの適用では、horizonごとのjump boundを入力とする。
固定horizon DDYの $L^T$ ではその入力を分解自身から生成した。

\begin{lemma}[有限変動jump補正の絶対収束]
上で構成した同じDDY分解 $N=L+Q$ に対し、$V_T=\operatorname{Var}_{[0,T]}(Q)$ と置く。
一つのfull-measure set上で全 $T\geq0$ について
\[
 \sum_{0<t\leq T}|\Delta Q_t|\leq V_T,\qquad
 \sum_{0<t\leq T}|\Delta L_t\Delta Q_t|\leq2cV_T,\qquad
 \sum_{0<t\leq T}(\Delta Q_t)^2\leq V_T^2
\]
が成り立ち、各級数は絶対収束する。さらに
\[
 D_T:=\sum_{0<t\leq T}\bigl((\Delta N_t)^2-(\Delta L_t)^2\bigr)
   =2\sum_{0<t\leq T}\Delta L_t\Delta Q_t+
      \sum_{0<t\leq T}(\Delta Q_t)^2,\qquad
 |D_T|\leq4cV_T+V_T^2
\]
であり、$D_T$ を定める級数も絶対収束する。
\end{lemma}

\begin{proof}
まず右連続BV path $A$ のsigned Stieltjes measureを $\nu$ とする。
各singletonのJordan total variationは $|\nu\{t\}|=|\Delta A_t|$ である。
$(0,T]$ の任意の有限部分集合上の和は $|\nu|((0,T])=\operatorname{Var}_{[0,T]}(A)$ 以下であり、
非負項級数のsummabilityと表示したboundが従う。
locally BV pathには $T$ で停止してこの結果を適用する。停止時刻までのjumpと変動量は元のpathに一致する。
同じ $Q$ に適用すると第一の評価を得る。
$|\Delta L_t|\leq2c$ と $|\Delta Q_t|\leq V_T$ を掛ければ残り二つの評価が従う。
最後に $\Delta N=\Delta L+\Delta Q$ を二乗展開し、絶対収束する級数を加える。
三角不等式から $D_T$ のboundを得る。全時刻の $L$ のjump boundが成立する一つの集合上で
任意の $T$ を扱うため、$T$ ごとに例外集合を取り直さない。
\end{proof}

\begin{lemma}[同じ有限グリッド上の二乗和]
有限の単調な時刻列 $u_0\leq\cdots\leq u_m$ に対し
\[
 \begin{aligned}
 E_G(Y,t)&=\sum_{k<m}(Y_{t\wedge u_{k+1}}-Y_{t\wedge u_k})^2,\\
 C_G(Y,Z,t)&=\sum_{k<m}(Y_{t\wedge u_{k+1}}-Y_{t\wedge u_k})
                         (Z_{t\wedge u_{k+1}}-Z_{t\wedge u_k}).
 \end{aligned}
\]
と置く。同じDDY分解の全pathと全時刻で
\[
 E_G(N,t)=E_G(L,t)+2C_G(L,Q,t)+E_G(Q,t),
\]
\[
 \left|\sqrt{E_G(N,t)}-\sqrt{E_G(L,t)}\right|
   \leq\operatorname{Var}_{[0,t]}(Q)
\]
である。この評価はグリッドに依存する定数を含まない。
\end{lemma}

\begin{proof}
各cellの増分を二乗展開すると第一式を得る。Cauchy--Schwarzから
$C_G(Y,Z,t)\leq\sqrt{E_G(Y,t)}\sqrt{E_G(Z,t)}$ であり、二乗展開と合わせれば
$\sqrt{E_G(Y+Z,t)}\leq\sqrt{E_G(Y,t)}+\sqrt{E_G(Z,t)}$ である。
また単調なcell chainの絶対増分和はpath variation以下なので、
\[
 \sqrt{E_G(Q,t)}\leq\sum_{k<m}|Q_{t\wedge u_{k+1}}-Q_{t\wedge u_k}|
    \leq\operatorname{Var}_{[0,t]}(Q).
\]
$N=L+Q$ と $L=N-Q$ の両方へrootの三角不等式を適用する。
\end{proof}

\begin{lemma}[階乗分割の交差項の積分極限]
$f$ を実数値càdlàg path、$Q$ を右連続locally BV pathとする。
$Q^T_t=Q_{t\wedge T}$ のsigned Stieltjes measureを $\eta_T$ とし、
$G_r$ を時刻 $k/(r+1)!$、$0\leq k\leq(r+1)(r+1)!$ のグリッドとする。このとき
\[
 C_{G_r}(f,Q,T)\longrightarrow\int_{(0,T]}\Delta f_s\,d\eta_T(s).
\]
従って同じDDY分解に対して、全pathと全有限horizonで
\[
 E_{G_r}(N,T)-E_{G_r}(L,T)
 \longrightarrow
 2\int_{(0,T]}\Delta L_s\,d\eta_T(s)+\int_{(0,T]}\Delta Q_s\,d\eta_T(s).
\]
\end{lemma}

\begin{proof}
$s>0$ を含む半開cell $(a_r(s),b_r(s)]$ を取り、
$D_r f(s)=f(b_r(s)\wedge T)-f(a_r(s)\wedge T)$ と置く。
$a_r(s)<s\leq b_r(s)$ で両端は $s$ へ収束するため、
$s\leq T$ なら $D_r f(s)\to\Delta f_s$ である。
この主張は $s$ 自体がmesh上にある場合も成り立つ。
càdlàg pathのcompact boundから $|D_r f(s)|\leq2\sup_{u\leq T}|f(u)|$ である。
有限測度 $|\eta_T|$ の下で優収束を適用すると、$(0,T]$ 上の $L^1$ 誤差が零へ収束し、
符号付き積分も収束する。
各停止済みcell上の積分は
$(f(b\wedge T)-f(a\wedge T))(Q(b\wedge T)-Q(a\wedge T))$ に等しい。
$r$ が十分大きければこれらのcellは $(0,T]$ を覆うので、積分は交差増分和に一致する。
$f=L,Q$ の二回の適用と有限和の二乗展開から第二式を得る。
\end{proof}

\begin{lemma}[jump積分の級数表示]
$f$ を実数値càdlàg path、$\eta$ を有限signed measureとすると、
\[
 \int_{(0,T]}\Delta f_t\,d\eta(t)
   =\sum_{t\in(0,T]}\Delta f_t\,\eta(\{t\}).
\]
特に、上の $\eta_T$ と同じDDY分解に対し、全pathで
\[
 E_{G_r}(N,T)-E_{G_r}(L,T)\longrightarrow
 2\sum_{0<t\leq T}\Delta L_t\Delta Q_t+\sum_{0<t\leq T}(\Delta Q_t)^2.
\]
\end{lemma}

\begin{proof}
各 $n$ について $|\Delta f_t|>1/(n+1)$ となる $(0,T]$ 上の時刻は有限なので、
非零jumpのsupportは可算である。cell近似の一様boundを極限へ移すと
$\Delta f$ も $(0,T]$ 上で有界であり、近似の可測性と合わせて $|\eta|$ 可積分となる。
supportをsingletonの可算非交和に分けて積分を足し合わせれば、第一式を得る。
停止済みFV pathのStieltjes原子は $0<t\leq T$ で
$\eta_T(\{t\})=\Delta Q_t$ である。終端時刻のjumpも保持される。
$f=L,Q$ に第一式を適用し、前補題の積分極限を書き換える。
\end{proof}

\begin{lemma}[補正後の候補の非負性]
同じDDY分解の補正級数を $D_t$、$L$ のglued quadratic variationを $V_t$ とする。
一つのfull-measure set上で全時刻 $t$ に対して $V_t+D_t\geq0$ である。
\end{lemma}

\begin{proof}
有限horizon座標 $n$ の局所化時刻を $\tau_n$、horizonを $T_n$ とする。
$t\leq\tau_n\wedge T_n$ では、停止済みsourceの二乗増分和は
元のsourceの同じ階乗grid上の和に等しい。
$L^{\tau_n}$ のquadratic構成が保存するforward-convex weightsとcutoffを使う。
補正の全列収束はそのweightsの下でも保たれるため、$L^{\tau_n}$ のconvex二乗増分和の極限に加えると、
同じweightsを用いた $N^{\tau_n}$ の二乗増分和は座標variationと $D_t$ の和へ収束する。
各近似は非負なので極限も非負である。
座標variationと $V$ の停止後の一致、ならびに $\tau_n\to\infty$ を使って全時刻へ移す。
座標について可算交叉を取った後に $t$ を選ぶので、零集合は時刻に依存しない。
\end{proof}

\begin{lemma}[補正と候補の経路正則性]
補正過程 $D$ と候補 $V+D$ はstrongly adaptedであり、一つのfull-measure set上で
右連続かつ局所有限変動である。
\end{lemma}

\begin{proof}
固定horizon $T$ では補正の係数
$a_s=2\Delta L_s\Delta Q_s+(\Delta Q_s)^2$ は絶対summableである。
累積和 $\sum_{s\in(0,T]}a_s\mathbf1_{\{s\leq t\}}$ の各項は右連続であり、
$|a_s|$ による級数の優収束から累積和も右連続である。
係数を $(a_s+|a_s|)-|a_s|$ と書くと、両累積和は非負・単調・有界なので、差は有限変動である。
$t\leq T$ でこの累積和は $D_t$ に等しいことから局所的な結論を得る。
adaptednessについては、各有限gridの二乗増分和の差が時刻 $t$ で可測であり、
全pathで $D_t$ に収束することを使う。$V$ の正則性と加法性から候補にも同じ結論が従う。
\end{proof}

\begin{lemma}[補正後の候補の単調性]
一つのfull-measure set上で、候補 $V+D$ は全時刻に関して単調増加である。
\end{lemma}

\begin{proof}
前の非負性の証明で用いた同じweightsによる収束を使う。
$s\leq t$ を固定し、$t+1<\tau_n$ となる局所化座標を選ぶ。
$s,t$ をhorizonで切った右factorial近似 $u_r,v_r$ で近似する。
十分大きい $r$ では $u_r\leq v_r<\tau_n$ であり、両時刻は同じ粗いgridに属する。
weightsのtail条件から、十分先のconvex近似が使う全てのgridはこの粗いgridを含む。
二乗増分和はgrid時刻間では単調なので、同じ非負weightsで平均した値も順序を保つ。
まずconvex近似の極限を取り、停止一致で $V+D$ へ戻す。
次に $r\to\infty$ とし、右連続性から $(V+D)_s\leq(V+D)_t$ を得る。
座標の収束と停止一致に関する可算交叉を先に取るため、同じevent上で任意の $s,t$ を選べる。
\end{proof}

\begin{lemma}[全経路で正則な代表元]
usual conditionsの下で、同じ候補 $V+D$ にindistinguishableなstrongly adapted過程 $A$ が存在し、
全経路で右連続・単調増加・局所有限変動であり、左極限を持ち、$A_0=0$ である。
\end{lemma}

\begin{proof}
候補の右連続性または単調性が破れる標本の集合を $B$ とする。前補題から $B$ は零集合であり、
usual conditionsにより時刻 $0$ のsigma algebraに属する。
$B$ 上で全時刻の値を零に置き換え、それ以外では候補をそのまま保つ。
この変更はstrong adaptednessとindistinguishabilityを保つ。
候補の初期値はglued variationの初期値と空のjump和から零であり、変更後も零である。
全経路の単調性から局所有限変動性が従い、局所有限変動性から左極限が得られる。
\end{proof}

\begin{lemma}[補正後のjump-square同定]
上の正則代表元 $A$ は一つの確率1の集合上で、全時刻について
$\Delta A_t=(\Delta N_t)^2$ を満たす。
\end{lemma}

\begin{proof}
各有限horizonで絶対summableな係数 $a_s$ に対し、累積和
$F(u)=\sum_{0<s\le u}a_s$ を考える。
$u\uparrow t$ のとき各項のindicatorは $\mathbf1_{\{s<t\}}$ に収束し、
$|a_s|$ で支配される。級数の優収束により
$F(t-)=\sum_{0<s<t}a_s$ となり、$t>0$ で $\Delta F_t=a_t$ を得る。
これを $a_s=2\Delta L_s\Delta Q_s+(\Delta Q_s)^2$ に適用する。
glued variationについて $\Delta V_t=(\Delta L_t)^2$ であるから、
\[
 \Delta(V+D)_t=(\Delta L_t)^2+2\Delta L_t\Delta Q_t+(\Delta Q_t)^2
 = (\Delta N_t)^2.
\]
時刻0では左jumpの規約から両辺とも零である。
級数の局所絶対収束とglued variationのjump等式は共通event上で全時刻に成立する。
最後に全時刻の経路一致を用いて正則代表元 $A$ に同じ等式を移す。
\end{proof}

\begin{lemma}[補正の閉停止と同じquadratic座標との一致]
有限値のcutoff $\tau$ に対して、閉停止した $L,Q$ のjumpから計算した補正は $D^\tau$ に等しい。
特に、scheduleの局所化時刻を $\tau_n$、そのquadratic座標を $V^n$ とすると、
同じ正則代表元 $A$ は
\[
 A^{\tau_n}=(V^n)^{\tau_n}
   +2\sum_{0<s\le\cdot}\Delta L^{\tau_n}_s\Delta Q^{\tau_n}_s
   +\sum_{0<s\le\cdot}(\Delta Q^{\tau_n}_s)^2
\]
をindistinguishabilityの意味で満たす。
\end{lemma}

\begin{proof}
閉停止した過程のjumpは $s\le\tau$ で元のjumpと一致し、$s>\tau$ で零となる。
従って閉停止成分の補正和を $0<s\le t$ 上で取ることは、元の成分の和を
$0<s\le t\wedge\tau$ 上で取ることに等しい。
この等式は経路ごとに成立し、停止時刻でのjumpも含む。
glued variationの座標一致 $V^{\tau_n}=(V^n)^{\tau_n}$ と
$A\sim V+D$ を組み合わせれば結論を得る。
\end{proof}

\begin{lemma}[有限変動成分の二乗残差]
右連続で左極限を持つ局所有限変動過程 $Q$ について、各経路と各 $T$ で
\[
 Q_T^2-Q_0^2-\sum_{0<s\le T}(\Delta Q_s)^2
   =2\int_{(0,T]}Q_{s-}\,dQ_s.
\]
右辺は $Q$ を $T$ で停止して得られるsigned Stieltjes measureに関する積分である。
\end{lemma}

\begin{proof}
factorial gridの各セル $(u,v]$ で左端点値 $Q_u$ を使う。
正の時刻を常に厳密に左から近似するため、この階段関数は $Q_{s-}$ に収束する。
有限horizon上のcadlag pathは有界であり、integratorの全変動measureは有限である。
従って優収束で左端点和がStieltjes積分へ収束する。
各有限gridで
\[
 2\sum Q_u(Q_v-Q_u)+\sum(Q_v-Q_u)^2=Q_T^2-Q_0^2
\]
が成立する。右辺の端点はhorizonで切り、十分細かいgridでは $T$ に達する。
二乗増分和の極限は有限変動cross-increment極限とjump積分の級数表示から
$\sum_{0<s\le T}(\Delta Q_s)^2$ である。両辺の極限を取って結論を得る。
\end{proof}

同じDDY成分にこの公式を適用し、$I_Q(t)=\int_{(0,t]}Q_{s-}\,dQ_s$ と置くと、
正則代表元 $A$ について共通event上の全時刻で
\[
 N_t^2-A_t=(L_t^2-V_t)+
   2\left(L_tQ_t-\sum_{0<s\le t}\Delta L_s\Delta Q_s+I_Q(t)\right)
\]
を得る。ここでは $Q_0=0$ を用いた。
このpathwiseな等式に、以下の共通局所化と可積分支配極限を適用する。

\begin{lemma}[全変動が可積分となる共通局所化]
同じDDY成分 $L,Q$ に対して有限値の局所化列 $\rho_n\uparrow\infty$ を選べ、
$\rho_n\le n+1$、$L^{\rho_n},Q^{\rho_n}$ はtrue martingalesであり、
\[
 |L^{\rho_n}_t|\le n+1+2c,\qquad
 |Q_s|\le n+1\quad(s<\rho_n)
\]
が成立する。前者は共通の確率1の集合上で全時刻の評価である。
さらに $\operatorname{Var}_{[0,\rho_n]}(Q)$ は可積分であり、
$Q^{\rho_n}$ の全horizonの変動量を支配する。
\end{lemma}

\begin{proof}
$A_t=\operatorname{Var}_{[0,t]}(Q)$ と置く。有限gridによる変動量の表示から $A$ は適合的であり、
局所有限変動と右連続性から右連続・単調で左極限を持つ。
両local martingalesのtrue localizers、$|L|$ と $A$ のlevel $n+1$ のpassages、
決定論的horizon $n+1$ の最小値を $\rho_n$ とする。
passage localizerとlocal martingale localizerの最小値なので、これは局所化列である。
零初期値によりlocalizerの時刻0のindicatorを除け、追加停止後もtrue martingale性が保たれる。
$L$ のjump boundからclosed-stop評価を得る。
$s<\rho_n$ では $|Q_s|\le A_s\le n+1$ である。
変動量の左jumpは $|\Delta Q|$ なので、$a=n+1$ とすると
\[
 A_{\rho_n}\le a+|\Delta Q_{\rho_n}|
 \le 2a+|Q_{\rho_n}|.
\]
最後の停止値はtrue martingaleの有界停止値であり可積分である。
停止した累積変動の可測性とこの支配から $A_{\rho_n}$ の可積分性を得る。
停止pathの変動量は $A_{t\wedge\rho_n}$ であり、$A_{\rho_n}$ 以下である。
\end{proof}

この同じ局所化上で、全ての有限gridについて
$I_G(Q^{\rho_n},Q^{\rho_n})$、$I_G(Q^{\rho_n},L^{\rho_n})$、
$I_G(L^{\rho_n},Q^{\rho_n})$ はtrue martingalesとなる。
最初の二つでは係数を $Q_s\mathbf1_{\{s<\rho_n\}}$ に置き換える。
これは適合的で有界であり、停止後のintegrator増分が零なので元の有限和と等しい。
最後の和ではclosed $L$ の有界性を使う。従って $Q_{\rho_n}$ の有界性も二乗可積分性も不要である。

\begin{lemma}[可積分支配下のマルチンゲール極限]
マルチンゲール列 $M^k$ が各時刻でほとんど確実に強適合的過程 $Y$ へ収束し、
各 $t$ に対して可積分な $B_t$ が存在して $|M^k_t|\le B_t$ がほとんど確実に全ての $k$ で成立するとする。
このとき $Y$ はマルチンゲールである。
\end{lemma}
\begin{proof}
極限も $B_t$ に支配されるので可積分である。
$s\le t$ と $E\in\mathcal F_s$ に対して、$E$ 上の積分に支配収束を適用すると
$\int_E Y_s\,d\mu=\int_E Y_t\,d\mu$ を得る。条件付き期待値の特徴付けから結論が従う。
\end{proof}

\begin{lemma}[自己積分と交差項の局所マルチンゲール性]
同じDDY成分について
\[
 I_t=\int_{(0,t]}Q_{s-}\,dQ_s,\qquad
 K_t=L_tQ_t-\sum_{0<s\le t}\Delta L_s\Delta Q_s
\]
は局所マルチンゲールである。
\end{lemma}
\begin{proof}
共通局所化を一つ固定し、$\bar L=L^{\rho_n}$、$\bar Q=Q^{\rho_n}$、
$a=n+1$、$b=n+1+2c$、$A=\operatorname{Var}_{[0,\rho_n]}(Q)$ と書く。
有限grid $G$ 上の自己積分和は、strict-prefix係数への置換により
\[
 |I_G(\bar Q,\bar Q)_t|
 \le a\operatorname{Var}_{[0,t]}(\bar Q)\le aA
\]
を満たす。その階乗grid極限は停止pathに対するStieltjes積分である。
停止前後の二乗公式と
$\sum_{s\le t}(\Delta Q^{\rho_n}_s)^2=\sum_{s\le t\wedge\rho_n}(\Delta Q_s)^2$
を比較すれば、この極限は元の $I_{t\wedge\rho_n}$ と等しい。

交差項には、零から始まる有限gridの終端を $u_G$ として
\[
 I_G(\bar L,\bar Q)_t+I_G(\bar Q,\bar L)_t
 =\bar L_{t\wedge u_G}\bar Q_{t\wedge u_G}-C^G_t
\]
という積の公式を用いる。ここで $C^G$ は同じcell上の交差増分和である。
$|\bar Q_s|\le\operatorname{Var}_{[0,s]}(\bar Q)$ と $|\bar L_s|\le b$ から積は $bA$ 以下、
また各 $\bar L$ の増分は絶対値 $2b$ 以下なので $|C^G_t|\le2bA$ である。
従って二つの積分和の和は $3bA$ に支配される。
交差増分和のStieltjes極限とjump級数表示、jumpの停止整合性から、
和全体の階乗grid極限は $K_{t\wedge\rho_n}$ となる。

両極限は全pathの時刻ごとの極限なので、有限和の強適合性も継承する。
$A$ の可積分性と前補題によって、両closed stopsはtrue martingalesである。
両過程の初期値は零なので、局所化の時刻0のindicator規約とも一致し、結論を得る。
この証明は $\int Q_-\,dL$ を別個のcompleted integralとして構成する必要がない。
\end{proof}

\begin{theorem}[元の局所マルチンゲールの二次変分]
usual conditionsを満たすfiltered probability space上の零始点・強適合的c\`adl\`ag
局所マルチンゲール $N$ に対して、適合的c\`adl\`ag増加過程 $V$ が存在し、
\[
 V_0=0,\qquad \Delta V=(\Delta N)^2,\qquad N^2-V\text{ は局所マルチンゲール}
\]
を満たす。jump等式は一つの共通の確率1の集合上で全時刻に成立する。
この $V$ は区別不能性を除いて一意である。
\end{theorem}
\begin{proof}
存在には正則化済み補正過程 $V$ をそのまま使う。
二乗公式は $N^2-V=(L^2-[L])+2(K+I)$ を与える。
$K+I$ は同じ共通局所化で局所マルチンゲールである。
その倍数を区別不能な正則過程 $N^2-V-(L^2-[L])$ に移してから、
 $L^2-[L]$ と局所化列をそろえて加えると、同じ $V$ の二乗残差の局所マルチンゲール性を得る。

一意性では、より一般に、適合的c\`adl\`ag・局所FV候補 $V,W$ が同じ初期値とjumpを持ち、
$N^2-V,N^2-W$ が局所マルチンゲールであるとする。候補の単調性はここでは不要である。
差 $D=V-W$ は局所FVの局所マルチンゲールであり、$\Delta D=0$ からほとんど確実に連続である。
各決定論的整数horizonで停止すると経路の全変動は有限になる。
停止過程の左極限はpredictableであり、連続なpath上では元の過程と等しい。
これを共通零集合上で正則化して、predictable FV局所マルチンゲールの剛性を適用すると、
各horizonで差は初期値の零に等しい。可算個のhorizonの共通eventを取れば全時刻で $V=W$ となる。
\end{proof}

従って、同じ $L$ のquadratic/Davis構成から得た補正は
\[
 [N]_T=[L]_T+D_T
\]
という二次変分の等式を与える。DDY閾値やscheduleが異なっても、同じsourceのcertificateは区別不能である。
この結果だけで非convexな二乗増分和の全列収束は主張しない。
同じscheduleの座標との停止整合性は上記の補題で得られる。
さらに任意の停止時刻との整合性も次のとおり得られる。

\begin{lemma}[二次変分の停止整合性]
上の仮定の下で、$\tau$ を無限大を取り得る任意の停止時刻とする。このとき
\[
 [N^\tau]\sim[N]^\tau.
\]
左辺は停止後のsourceから独立に構成してよい。
\end{lemma}
\begin{proof}
零始点の右連続局所マルチンゲールは、任意のclosed stoppingで局所マルチンゲールのままである。
実際、元のlocalizerで停止したtrue martingaleをさらに $\tau$ で止め、
二つの停止の交換と時刻0のindicatorの除去を用いればよい。
従って $N^\tau$ にも存在定理を適用できる。

$V=[N]$ とすると、$V^\tau$ の強適合性、c\`adl\`ag性、単調性と零初期値は停止で保たれる。
単調性から局所有限変動性も得る。
$t\le\tau$ では停止前後のjumpが同じであり、$\tau<t$ では両jumpとも零だから
$\Delta(V^\tau)=(\Delta(N^\tau))^2$ である。
二乗残差は
\[
 (N^\tau)^2-V^\tau=(N^2-V)^\tau
\]
であり、右辺は零始点の局所マルチンゲールを停止した過程である。
これで $V^\tau$ は $N^\tau$ の二次変分の全条件を満たす。
停止後のsourceに対する一意性から、独立に構成したどの二次変分とも区別不能となる。
\end{proof}

\begin{lemma}[二次変分の平方根の三角不等式]
零始点・強適合的c\`adl\`ag局所マルチンゲール $X,Y$ に対して、共通の確率1の集合上で
\[
 \sqrt{[X+Y]_t}\le\sqrt{[X]_t}+\sqrt{[Y]_t}\qquad(t\ge0)
\]
が成立する。
\end{lemma}
\begin{proof}
$V=[X]$、$W=[Y]$、$S=[X+Y]$ と書く。実数 $a$ に対して
\[
 A^a=aS+(1-a)V+(a^2-a)W
\]
は零始点の適合的c\`adl\`ag・局所FV過程である。
二乗残差には
\[
 (X+aY)^2-A^a
 =a\bigl((X+Y)^2-S\bigr)+(1-a)(X^2-V)+(a^2-a)(Y^2-W)
\]
が成立するので、これは局所マルチンゲールである。
また $\Delta A^a=(\Delta X+a\Delta Y)^2$ である。
従って、単調性を要求しない一意性の比較を用いると、$A^a\sim[X+aY]$ を得る。

各有理数 $a$ について $X+aY$ に二次変分の存在定理を適用し、可算交叉を取る。
すると一つのevent上で全時刻 $t$ と全有理数 $a$ に対して
\[
 0\le aS_t+(1-a)V_t+(a^2-a)W_t
   =W_ta^2+(S_t-V_t-W_t)a+V_t
\]
となる。時刻とpathを固定した二次式の連続性から全実数 $a$ でも非負である。
判別式を評価すると $(S_t-V_t-W_t)^2\le4V_tW_t$ を得る。
$V_t,W_t\ge0$ なので $S_t\le(\sqrt{V_t}+\sqrt{W_t})^2$ となり、結論が従う。
\end{proof}

\begin{lemma}[有限グリッド上の $L^1$ Davis 評価]
離散時間の実数値マルチンゲール $Y$ に対して
\[
 Y_n^*=\max_{0\le k\le n}|Y_k|,\qquad
 R_n=\sqrt{Y_0^2+\sum_{k<n}(Y_{k+1}-Y_k)^2}
\]
とおく。この二つは可積分であり、
\[
 E[Y_n^*]\le6E[R_n],\qquad E[R_n]\le7E[Y_n^*]
\]
が成立する。
各時刻の二乗可積分性は要求しない。
\end{lemma}
\begin{proof}
それぞれ $\sum_{k\le n}|Y_k|$ と
$|Y_0|+\sum_{k<n}|Y_{k+1}-Y_k|$ で抑えられるので可積分である。
有界な適合的係数 $H_k$ に対し、$H_k(Y_{k+1}-Y_k)$ は可積分であり、
$\mathcal F_k$ に関する条件付き期待値は零である。
従って左端係数による有限積分和は零始点のマルチンゲールとなる。
経路ごとのDavis不等式に現れる係数は絶対値が $1$ 以下なので、
その不等式を積分すると補正積分の期待値が消え、結論を得る。

逆向きには、$m\ge|x|$, $m,q\ge0$ に対して
\[
 F(x,m,q)=-2m+\sqrt{m^2+q}+\frac{m^2-x^2}{2\sqrt{m^2+q}}
\]
を使う。零除算は零とする。$M=\max\{m,|x+d|\}$ とおくと
\[
 F(x+d,M,q+d^2)-F(x,m,q)
 \le-\frac{xd}{\sqrt{m^2+q}}+5(M-m).
\]
この一段階評価は次のように得る。$r=\sqrt{m^2+q}>0$,
$u=\sqrt{M^2+q+d^2}$ とする。
$M\le2m$ の場合、平方根の接線評価と $u\ge r$ により、左辺に $xd/r$ を加えたものは
$(M^2-m^2)/r-2(M-m)\le M-m$ 以下である。
$M>2m$ の場合は $M=|x+d|$, $|d|\le M+m\le3(M-m)$ であり、
$u-r\le4(M-m)$ と $xd/r\le|d|$ を使う。
新しいpotentialの分子 $M^2-(x+d)^2$ は零、古い分子は非負なので、目的の上界を得る。
$r=0$ なら $x=m=q=0$ となり、直接計算で評価できる。

$x=Y_k$, $m=Y_k^*$, $q=R_k^2$, $d=Y_{k+1}-Y_k$ として和を取る。
$F(Y_0,|Y_0|,Y_0^2)\le0$ と $F(Y_n,Y_n^*,R_n^2)\ge R_n-2Y_n^*$ から
\[
 R_n\le7Y_n^*-\sum_{k<n}
 \frac{Y_k}{\sqrt{R_k^2+(Y_k^*)^2}}(Y_{k+1}-Y_k)
\]
を得る。同じ絶対値1以下の適合係数を使っているので、この積分和の期待値も零である。
\end{proof}

DDY分解の共通局所化 $\rho_j$ では、$L^{\rho_j}$ と $Q^{\rho_j}$ がともに
マルチンゲールである。従って元の過程 $N^{\rho_j}=L^{\rho_j}+Q^{\rho_j}$ を
任意の有限時刻グリッドで標本化した列にも、この評価が適用できる。
ここで停止列は、閉じて停止した $Q$ の全変動が可積分となるものと同じである。
平方根は凹関数なので、個々のグリッドの逆評価を凸結合したenergyの平方根へそのまま移せない。
そこで元の二乗増分和の収束を示す。

\begin{lemma}[有界マルチンゲールの元の二乗増分和]
usual conditionsの下で、$M$ を強適合的c\`adl\`agマルチンゲールとし、
有限horizon $T$ まで $|M_t|\le C$ が全pathで成立するとする。
階乗グリッドを $\pi_i$ とし、
\[
 I_i=2\sum_{[u,v]\in\pi_i}M_u(M_v-M_u),\qquad
 E_i=\sum_{[u,v]\in\pi_i}(M_v-M_u)^2
\]
とおく。保存済みの二次変分構成のmartingale成分を $Y$、増加成分を $V$ とすると、
$I_i\to Y_T$ および $E_i\to V_T$ が $L^2$ で成立する。
有界性が全時刻共通にほとんど確実に成立する場合にも、$E_i\to V_T$ は確率収束する。
\end{lemma}
\begin{proof}
同じ $V$ の予測可能エネルギー測度を $\nu$ とする。
グリッド左端係数 $H_i$ は有界予測可能初等過程であり、$(0,T]$ 上で $M_-$ に各点収束する。
$\nu$ は有限であり、時刻零と $T$ より後を測度零にする。
従って支配収束により $H_i\to M_-$ が $L^2(\nu)$ で成立する。
格子によらない等長性
\[
 \|(H_i-H_j)\mathbin{\cdot}M_T\|_{L^2(P)}
 =\|H_i-H_j\|_{L^2(\nu)}
\]
から $I_i$ は $L^2(P)$ でCauchyである。その極限を $I$ とする。
保存済みのforward convex weightsを $I_i$ に適用しても極限は $I$ のままである。
一方、その同じ重みによる積分和はほとんど確実に $Y_T$ へ収束する。
確率極限の一意性から $I=Y_T$ である。
$I_i+E_i=M_T^2-M_0^2$ と保存済みの二乗分解から $E_i\to V_T$ を得る。

有界性の例外が共通零集合に含まれる場合、usual conditionsによりその集合は
$\mathcal F_0$ に属する。そこで $M$ を零に置き換えると、全pathで有界な
区別不能のc\`adl\`agマルチンゲールを得る。
同じ重みと同じ二つの成分はこのsourceにも二次変分dataを与える。
元の二乗増分和も区別不能なsource間でほとんど確実に等しいので、確率収束を移せる。
\end{proof}

\begin{theorem}[一般局所マルチンゲールの逆Davis評価]
usual conditionsの下で、$N$ を零始点・強適合的c\`adl\`ag局所マルチンゲールとする。
任意の内在的二次変分と有限 $T$ に対して
\[
 E\sqrt{[N]_T}\le7E[N_T^*]
\]
が成立する。期待値は無限でもよい。
\end{theorem}
\begin{proof}
DDY共通停止列の一つ $\rho_j$ を固定する。$L^{\rho_j}$ はほとんど確実に一様有界であり、
前補題から元の二乗増分和がその二次変分へ確率収束する。
$N^{\rho_j}$ と $L^{\rho_j}$ の二乗増分和の差は、有限変動補正へpathwiseに収束する。
和を取り、停止整合性を使うと、元の $N^{\rho_j}$ の二乗増分和が
$[N]_{T\wedge\rho_j}$ へ確率収束する。
ほとんど確実に収束する部分列を選び、平方根を取る。
各グリッドの最大値は $(N^{\rho_j})_T^*$ 以下なので、有限グリッドの逆評価とFatouの補題から
\[
 E\sqrt{[N]_{T\wedge\rho_j}}
 \le7E[(N^{\rho_j})_T^*]\le7E[N_T^*]
\]
を得る。ほとんど確実に $\rho_j\to\infty$ なので、同じ確率1の集合上で左辺の被積分関数は最終的に
$\sqrt{[N]_T}$ と等しい。再びFatouを使い、二次変分の一意性で任意のcertificateへ移す。
\end{proof}

同じDDY局所化は、次の平方根の一様可積分性にも使える。

\begin{lemma}[凸結合した二乗増分和の平方根の一様可積分性]
上の共通停止列の一つと有限horizonを固定する。
元の停止過程とその $L$ 成分のfactorial grid上の二乗増分和を、それぞれ $E_i,B_i$ とする。
任意のforward convex weights $W$ に対し、$\sqrt{W(E)_k}$ は一様可積分である。
\end{lemma}
\begin{proof}
$A=\operatorname{Var}(Q^{\rho_j})_\infty$ とおくと、$A\ge0$ は可積分である。
同じグリッドの平方根の三角不等式と有限変動評価から
$\sqrt{E_i}\le\sqrt{B_i}+A$、従って $E_i\le2B_i+2A^2$ となる。
重みの非負性と和が $1$ であることから
\[
 W(E)_k\le2W(B)_k+2A^2,\qquad
 \sqrt{W(E)_k}\le2W(B)_k+2(1+A).
\]
停止した $L$ は有界なマルチンゲールなので、そのterminal valueは二乗可積分である。
二次変分近似の評価により $B_i$ は一様可積分であり、任意のforward convexification
$W(B)_k$ も一様可積分である。右辺の残りは一つの可積分関数なので、結論が従う。
この証明では $A^2$ の可積分性は使わない。
\end{proof}

\begin{lemma}[閉じた停止後の終端平方根の $L^1$ 極限]
同じ共通停止列の一つ $\rho$ と有限horizon $T$ を固定する。
$L^\rho$ の二次変分構成から得た有限horizonの候補を $V^L$、その構成の重みと部分列による
元の停止sourceの凸結合energyを $E_k$ とする。DDYのjump補正を $C$ とすると、
\[
 \sqrt{E_k(T)}\longrightarrow
 \sqrt{V^L_T+C_{T\wedge\rho}}\quad\text{in }L^1,
\]
かつ右辺は可積分である。
\end{lemma}
\begin{proof}
停止した過程のgrid上の増分和は、元の過程を $t\wedge\rho$ で評価した増分和に等しい。
従ってjump補正のgrid極限を、$t\le\rho$ に制限することなく適用できる。
同じ重みで $L^\rho$ 側のenergyが $V^L$ に収束し、補正項が $C_{t\wedge\rho}$ に収束するので、
$E_k(T)\to V^L_T+C_{T\wedge\rho}$ がほとんど確実に成立する。
平方根の連続性と上の一様可積分性から、Vitaliの定理によって結論を得る。
$L^\rho$ の終端は二乗可積分なので、ここで使う有限horizonの二次変分データは実際に構成できる。
\end{proof}

\begin{lemma}[元の二次変分の停止値との同定]
上の記号の下で、元の過程 $N$ の正則化済み二次変分を $V$ とすると、
\[
 V^L_T+C_{T\wedge\rho}=V_{T\wedge\rho}\quad\text{a.s.}
\]
従って同じ凸結合energyについて
$\sqrt{E_k(T)}\to\sqrt{V_{T\wedge\rho}}$ が $L^1$ で成立し、極限は可積分である。
\end{lemma}
\begin{proof}
有限horizonの候補 $V^L$ は $(L^\rho)^T$ の二次変分の条件を満たす。
実際、その構成のマルチンゲールは二乗残差に等しく、jumpはsourceのjumpの二乗である。
一方、大域的に構成した $[L]$ を $\rho$ と $T$ で順に停止した過程も、
停止整合性により同じsourceの二次変分となる。
一意性から $V^L_T=[L]_{T\wedge\rho}$ を得る。
正則化済みの $V$ は $[L]+C$ と区別不能なので、ランダム時刻 $T\wedge\rho$ でも
この等式を使える。最後に上の $L^1$ 極限の終端値をこの等式で置き換える。
\end{proof}

\begin{lemma}[一般の局所マルチンゲールのDavis評価]
usual conditionsの下で、$N$ を零始点・強適合的c\`adl\`ag局所マルチンゲールとする。
任意の有限horizon $T$ に対し、非負拡張実数の期待値として
\[
 E\sup_{t\le T}|N_t|\le6E\sqrt{[N]_T}
\]
が成立する。右辺が有限なら、左辺の最大値も可積分である。
\end{lemma}
\begin{proof}
まず共通局所化の一つ $\rho$ を固定する。
一つのcoarse factorial gridの最大値は、それ以後の各gridの最大値以下である。
従ってtail convex weightsによる平均を取り、有限グリッドのDavis評価と
平方根の凹性を使うと、その最大値の期待値は $6E\sqrt{E_k(T)}$ 以下となる。
上の $L^1$ 収束で $k\to\infty$ として、右辺を
$6E\sqrt{[N]_{T\wedge\rho}}$ に置き換えられる。
coarse gridを細かくすると、右連続性と単調収束により、全時刻の最大値に対して
\[
 E\sup_{t\le T}|N^\rho_t|\le6E\sqrt{[N]_{T\wedge\rho}}
\]
を得る。この時点で左辺の可積分性も従う。
次に共通停止列 $\rho_n\to\infty$ を使う。
ほとんど全ての経路では最終的に $\rho_n\ge T$ となるため、停止過程の最大値は元の最大値に等しくなる。
右辺は二次変分の単調性から $6E\sqrt{[N]_T}$ 以下である。
Fatouの補題により停止を除去できる。二次変分の一意性から、構成に用いた候補を
同じsourceの任意の二次変分に置き換えても結論が成立する。
\end{proof}

\begin{lemma}[無限horizonと全分解上の最大値評価]
同じ仮定の下で $R_N:=\sup_{t\ge0}\sqrt{[N]_t}$ とおくと、
\[
 E\sup_{t\ge0}|N_t|\le6E R_N
\]
が非負拡張実数の期待値として成立する。
零始点の過程 $Y$ に対し、$Y=N+A$ が区別不能の意味で成立する全ての零始点分解を考える。
ここで $N$ は強適合的c\`adl\`ag局所マルチンゲール、$A$ は強適合的c\`adl\`ag局所有限変動過程である。
\[
 \mathcal J(Y):=\inf_{Y=N+A}E[R_N+\Var(A)_\infty]
\]
とおくと、$E\sup_{t\ge0}|Y_t|\le6\mathcal J(Y)$ が成立する。
下限では有限costを仮定せず、空の分解族の下限は $\infty$ とする。
\end{lemma}
\begin{proof}
整数horizonでの最大値は単調に全時刻の最大値へ増加する。
有限horizonのDavis評価の右辺を $6E R_N$ で抑え、単調収束により最初の主張を得る。
二次変分の単調性により $R_N$ は整数時刻で検出できるので可測である。
各分解について、$A_0=0$ と全変動の定義から
\[
 \sup_t|Y_t|\le\sup_t|N_t|+\Var(A)_\infty.
\]
期待値を取りDavis評価を使うと、右辺は
$6E R_N+E\Var(A)_\infty\le6E[R_N+\Var(A)_\infty]$ 以下となる。
全ての分解について下限を取ればよい。
各 $N$ の二次変分は存在し、全時刻共通の零集合を除いて一意であるため、
二次変分の構成上の選択はcostに影響しない。
\end{proof}

\begin{lemma}[閉じた停止と停止直前の延長の下限]
任意の停止時刻 $\tau$ に対し、$\mathcal J(Y^\tau)\le\mathcal J(Y)$ である。
停止時刻には $\infty$ を許す。さらに
\[
 \mathcal J_{\tau-}(Y):=
 \inf\{\mathcal J(L):L_t=Y_t\text{ for all }t<\tau\text{ a.s.}\}
\]
と置く。ただし一致は全時刻共通の零集合を除いて要求する。このとき
\[
 \mathcal J_{\tau-}(Y)\le\mathcal J(Y^\tau)\le\mathcal J(Y),\qquad
 E\sup_{t<\tau}|Y_t|\le6\mathcal J_{\tau-}(Y).
\]
$\tau=0$ で空となるsupremumは $0$ とする。
\end{lemma}
\begin{proof}
分解 $Y=N+A$ を同じ $\tau$ で停止すると、
$Y^\tau=N^\tau+A^\tau$ は再び許容される零始点分解となる。
停止した二次変分は $[N]^\tau$ なので、その全時刻平方根は $R_N$ 以下である。
時刻を $t\wedge\tau$ に制限しても全変動は増えない。
従って停止した分解のcostは元のcost以下であり、全分解について下限を取れる。
また $Y^\tau$ は停止直前まで $Y$ と一致する大域延長なので、外側の下限の上界を与える。
最後に任意の延長 $L$ について $\sup_{t<\tau}|Y_t|\le\sup_t|L_t|$ がa.s.で成立する。
既に得た最大値評価を適用し、延長について下限を取れば最後の不等式を得る。
\end{proof}

\begin{lemma}[二つの下限の三角不等式]
同じ仮定の下で
\[
 \mathcal J(X+Y)\le\mathcal J(X)+\mathcal J(Y),\qquad
 \mathcal J_{\tau-}(X+Y)\le\mathcal J_{\tau-}(X)+\mathcal J_{\tau-}(Y)
\]
が成立する。後者の $\tau$ は同じものであり、その停止時刻性は必要ない。
\end{lemma}
\begin{proof}
二つの分解 $X=N+A$、$Y=M+B$ の成分を足すと、$X+Y=(N+M)+(A+B)$ の許容分解を得る。
二次変分の三角不等式から $R_{N+M}\le R_N+R_M$ がa.s.で成立する。
全変動の劣加法性から $\Var(A+B)_\infty\le\Var(A)_\infty+\Var(B)_\infty$ でもある。
全時刻variationは整数horizonのvariationのsupremumであり、右連続な適合過程について可測である。
従って期待値を足し合わせることができ、和の分解のcostは二つのcostの和以下となる。
二つの分解族について独立に下限を取ると最初の主張を得る。
同じstrict prefixに一致する二つの大域延長を足せば、和のsourceの大域延長になる。
最初の不等式を適用して二つの延長族について下限を取れば後半を得る。
この議論は無限costや空の分解族も含む。
\end{proof}

\begin{lemma}[符号反転とsuccessive differences]
同じ仮定の下で、$\mathcal J(-X)=\mathcal J(X)$ および
$\mathcal J_{\tau-}(-X)=\mathcal J_{\tau-}(X)$ である。
従って一つの固定した $\tau$ について
\[
 \mathcal J_{\tau-}(Z_k-X)\longrightarrow0
 \quad\Longrightarrow\quad
 \mathcal J_{\tau-}(Z_{k+1}-Z_k)\longrightarrow0.
\]
\end{lemma}
\begin{proof}
$N$ の二次変分はそのまま $-N$ の二次変分の条件を満たす。
二乗残差は変わらず、jumpの符号反転は二乗で消えるからである。
また全変動も符号反転で不変である。従って分解 $(N,A)$ を $(-N,-A)$ に移すとcostは保存される。
下限を取り、符号反転を二度適用すれば $\mathcal J$ の等式を得る。
同じprefix上で一致する延長も符号反転できるため、外側の下限にも同じ議論を適用できる。
三角不等式とこの対称性から
\[
 \mathcal J_{\tau-}(Z_{k+1}-Z_k)
 \le\mathcal J_{\tau-}(Z_{k+1}-X)+\mathcal J_{\tau-}(Z_k-X)
\]
となり、非負性と合わせて結論を得る。無限costも許すが、共通停止時刻の存在をこの補題からは主張しない。
\end{proof}

\begin{lemma}[有限costからの局所成分の構成]
usual conditionsの下で、$\tau\le T$ を有限停止時刻とする。
$\mathcal J_{\tau-}(X)<r$ なら、$X$ と区別不能な過程 $X'$ と、
$t<\tau$ で点ごとに $X'_t=N_t+A_t$ を満たす局所マルチンゲール $N$、
大域的有限変動過程 $A$ が存在し、
\[
 E\sup_{t\le T}|N_{t\wedge\tau}|
 +E\Var_{[0,T]}(A^{\tau-})<6r
\]
となる。両成分は適合的で右連続かつ左極限を持ち、左辺の二項は有限である。
\end{lemma}
\begin{proof}
下限の定義から延長 $Y$ と分解 $Y=N+B$ を選び、そのcostを $r$ 未満にする。
このcostは有限である。$A=B^T$ と置くと経路の大域的有限変動性を得る。
停止前の全変動は $B$ の全変動以下であり、停止した $N$ の最大値には
全時間のDavis評価を適用できる。従って左辺は選んだcostの6倍以下となる。
$t<\tau$ では $X'_t=N_t+A_t$、それ以外では $X'_t=X_t$ と定義する。
延長と分解の一致が成立する共通の確率1の集合上で、全時刻 $X'=X$ となる。
これにより停止前の点ごとの等式と、元の過程との区別不能性が同時に得られる。
\end{proof}

\begin{lemma}[和可能なcostからの成分列]
usual conditionsの下で、$\tau\le T$ を一つの有限停止時刻とする。
$\sum_k\mathcal J_{\tau-}(X_k)<\infty$ なら、前の補題のexact代表元 $X'_k$ と成分 $(N_k,A_k)$ を選び、
\[
 \sum_k\left(E\sup_{t\le T}|N_{k,t\wedge\tau}|
       +E\Var_{[0,T]}(A_k^{\tau-})\right)
 \le6\left(\sum_k\mathcal J_{\tau-}(X_k)+1\right)<\infty
\]
とできる。さらに一つの確率1の集合上で、全 $k,t$ について $X'_{k,t}=X_{k,t}$ であり、
\[
 \sum_k\left(\sup_{t\le T}|N_{k,t\wedge\tau}|
       +\Var_{[0,T]}(A_k^{\tau-})\right)<\infty
\]
となる。
\end{lemma}
\begin{proof}
$c_k=\mathcal J_{\tau-}(X_k)$、$\varepsilon_k=2^{-k-1}$ と置く。
各 $c_k$ は有限なので $c_k<c_k+\varepsilon_k$ であり、前の補題からcostが
$6(c_k+\varepsilon_k)$ 未満の代表元と成分を選べる。
$\sum_k\varepsilon_k=1$ より期待値の総和の評価を得る。
各martingale最大値とstrict-prefix全変動は可測であり、Tonelliにより経路ごとの総和はa.s.有限である。
全 $k$ の区別不能性のeventとの可算交叉を取れば、同じ集合上で元の過程との全時刻一致も得る。
\end{proof}
ここで $\tau$ と和可能なcostは入力であり、共通局所化の構成は含めない。
また成分は一般の分解から選ばれるので、original-sourceの実現された積分への所属は別に証明する必要がある。

\begin{lemma}[停止直前の比較と境界jump]
usual conditionsの下で、$\tau\le T$ を有限停止時刻とする。このとき
\[
 \mathcal J(X^{\tau-})\le2\mathcal J(X).
\]
さらに $X_0=0$ なら
\[
 \mathcal J_{\tau-}(X)\le\mathcal J(X^{\tau-})\le2\mathcal J_{\tau-}(X).
\]
$X$ が零始点の適合的c\`adl\`ag過程なら
\[
 \mathcal J(X^\tau)\le2\mathcal J_{\tau-}(X)+E|\Delta X_\tau|.
\]
いずれも無限costを許す。
\end{lemma}
\begin{proof}
分解 $X=N+A$ を固定する。$A$ を決定論的な $T$ で停止してもstrict prefixは変わらない。
従って $X^{\tau-}=N^\tau+(A^{\tau-}-B^\tau(N))$ を許容分解にできる。
二次変分の単調性、零初期値とjump条件より
$|\Delta N_\tau|\le\sqrt{[N]_\tau}$ である。
停止したマルチンゲールのroot costは元のroot cost以下であり、
有限変動成分の全変動は元の全変動に $|\Delta N_\tau|$ を加えたもの以下となる。
従ってこの分解のcostは元のcostの2倍以下であり、下限を取れば第一式を得る。
次に停止前で一致する延長の分解を選ぶ。その成分の和と $X$ の左極限は、
$\tau>0$ では停止前の一致から、$\tau=0$ では零初期値から一致する。
第一式の分解ごとの評価を適用して延長と分解の下限を取れば第二式の上側評価を得る。
下側評価は $X^{\tau-}$ 自身が延長であることによる。
最後に境界過程 $B^\tau(X)$ は適合的な零始点FV過程であり、
零マルチンゲールとの分解から $\mathcal J(B^\tau(X))\le E|\Delta X_\tau|$ となる。
$X^\tau=X^{\tau-}+B^\tau(X)$ に三角不等式を適用すれば第三式を得る。
\end{proof}

これらは零始点coreと外側の下限の評価、および具体的な局所成分の構成である。
主定理の零始点の代表列に必要な代数・共通prelocal制御・入力生成は、次節の位相比較と共通局所化で供給する。
一般の非零初期値は別座標として扱い、この零始点costに隠さない。

有限 $M^2$ quadratic rootは exact $j^1$ そのものではない。
stopping-time jump seminormは、二次変分と全変動の和からなるcostとは別に扱う。
quadratic variationは定数 martingaleを検出しないので、$N-N_0$ を用い、一般過程の距離には
初期値を別座標として持たせる。

\subsection{任意の分解から正準成分への定量評価}
\label{sec:canonical-cost}
補償定理を次の形で用いる。適合的・零始点・càdlàg な過程 $D$ が可積分な全変動を持ち、
有限の決定論的時刻以後一定なら、予測可能・零始点・càdlàg で可積分変動を持つ $D^p$ が存在し、
$D-D^p$ は一様可積分マルチンゲールとなる。
これは\emph{双対予測可能射影}であり、
\[
 E\int K\,dD^p=E\int K\,dD
 \quad\text{（任意の有界予測可能な }K\text{）}
\]
で特徴付けられる。これは $D^p_t=E[D_t\mid\mathcal F_{t-}]$ という主張ではない。
存在は標準的な双対射影定理による。その特徴付けと jump 公式は
\cite[Section 4.3]{Nikeghbali}を参照されたい。必要な変動評価は
\begin{equation}\label{eq:compensator-variation}
 E\Var(D^p)\leq E\Var(D)
\end{equation}
である。実際、$dD^p$ の予測可能 Hahn 符号を $K$ に選ぶと、左辺は期待全変動となり、
右辺は $E\Var(D)$ 以下である。二つの増加 Jordan 成分を別々に補償し、正値性を使ってもよい。
$K=1_F1_{(s,t]}$、$F\in\mathcal F_s$ に対する同じ等式が $D-D^p$ のマルチンゲール性を与える。
Lean では切断と予測可能有限変動近似の極限からこの補償過程を構成している。
本稿では存在定理自体を標準的な基礎定理として明示的に使用する。

\begin{lemma}[有限 horizon 後一定な利得の正準成分 cost]
\label{lem:canonical-cost}
零始点・càdlàg な $X$ は $T<\infty$ 以後一定で、
正準分解 $X=M^{\rm can}+A^{\rm can}$ を持つとする。通常条件の下で
\[
 \begin{aligned}
 E\Var_{[0,T]}(A^{\rm can})&\leq12\mathcal J(X),\\
 E\sup_{t\leq T}|M^{\rm can}_t-M^{\rm can}_0|
       +E\Var_{[0,T]}(A^{\rm can})&\leq30\mathcal J(X)
 \end{aligned}
\]
が拡張値の不等式として成立する。$\mathcal J$ の下限で許す FV 成分は、
予測可能とは限らず適合的でよい。
\end{lemma}
\begin{proof}
下限を定める一つの有限 cost $j$ の分解を固定する。
strict-prefix witness の構成により、決定論的停止 $\sigma=T+1$ の前で $X$ を表す
成分 $N,B$ を
\[
 c_N=E\sup_t|N^\sigma_t|,\qquad
 c_B=E\Var(B^{\sigma-}),\qquad c_N+c_B\leq6j
\]
となるように選べる。初期値をそのまま保持し、
$J^N_t=\Delta N_\sigma1_{[\sigma,\infty)}(t)$ として
\[
 D=B^{\sigma-}-B_0-J^N
\]
と置く。これは適合的であり、
\[
 E\Var(D)\leq c_B+E|\Delta N_\sigma|\leq c_B+2c_N
\]
より可積分変動を持つ。この具体的な残差を補償する。
\[
 \widehat M=N^\sigma+D-D^p,\qquad \widehat A=B_0+D^p
\]
は $(N+B)^{\sigma-}=X$ の special 分解である。
最後の等式には、開区間上の一致だけでなく、$X$ が $T$ 以後一定であることを用いる。
式\eqref{eq:compensator-variation}より
\[
 E\Var(\widehat A)\leq c_B+2c_N\leq2(c_N+c_B)\leq12j
\]
となる。ここで初めて正準分解の一意性を適用し、構成した分解の中心化成分を
$X$ の正準成分に同定する。$j$ の下限を取れば第一の評価を得る。
零始点性より
\[
 \sup_{t\leq T}|M^{\rm can}_t-M^{\rm can}_0|
 \leq\sup_t|X_t|+\Var_{[0,T]}(A^{\rm can})
\]
である。既証の最大値評価 $E\sup_t|X_t|\leq6\mathcal J(X)$ と合わせて、
$6+12+12=30$ を得る。無限 cost の場合も拡張値のまま扱える。
\end{proof}

正則な零始点 $X$ と $\tau\leq T$ に対する prelocal witness の
M 最大値 cost と変動 cost の和を $c$ とする。
後の $r^1$ 局所化小節で示す witness と strict-prefix infimum の比較より $\mathcal J_{\tau-}(X)\leq16c$ である。
二次変分 cost の節で示した境界 jump 補正より
\[
 \mathcal J(X^\tau)\leq2\mathcal J_{\tau-}(X)+E|\Delta X_\tau|
 \leq32c+E|\Delta X_\tau|
\]
となる。独立に構成した actual integral の利得が $X^\tau$ なら、
この同じ $T$ 以後一定な利得に補題\ref{lem:canonical-cost}を適用して
\begin{equation}\label{eq:actual-closed-cost}
 E\sup_{t\leq T}|M^{\rm actual}_t-M^{\rm actual}_0|
 +E\Var_{[0,T]}(A^{\rm actual})
 \leq30\bigl(32c+E|\Delta X_\tau|\bigr)
\end{equation}
を得る。actual integral の存在自体は別途供給する。
この議論はその正準成分に定量評価を渡すものであり、一意性だけで評価が得られるわけではない。


\clearpage
\section{共通局所化と成分級数}
\label{sec:localization}
本付録では有限horizonの評価を共通停止列へまとめ、成分級数の極限を構成する。
\subsection{Elementary testsに一様な極限への接続}

強progressiveかつ右連続なsource $S$ のelementary積分列を $X^n$、その指定された利得を $X$ とする。
有限horizon $T$ におけるcapped絶対envelopeを $e_T$ と書き、
unit-bounded predictable elementary tests $K$ に対して
\[
 \Phi_{K,T}(Y)=E[e_T(K\mathbin{\cdot}Y)],\qquad
 g_{n,T}=\sup\bigl(\{0\}\cup\{\Phi_{K,T}(X^n-X):K\}\bigr)
\]
と置く。ここで積分はelementary testの有限和であり、一般predictable積分の存在は用いない。

\begin{lemma}[Cauchy性と固定test極限から一様極限へ]
realized strategyの代表列について、各 $T$ で $g_{n,T}\to0$ であり、さらに
\[
 \sum_{r\ge0}2^{-r-1}g_{n,r+1}\longrightarrow0.
\]
\end{lemma}

\begin{proof}
$\varepsilon>0$ と $T$ を固定する。代表列のÉmery-Cauchy性から、
$m,n\ge N$ なら全ての $K$ で
$\Phi_{K,T}(X^m-X^n)<\varepsilon$ となる共通の $N$ を選べる。
ここで初めて $n\ge N$ と $K$ を固定する。固定testでの誤差は確率収束で零へ行くので、
補助添字 $m$ の部分列を選ぶとcapped誤差がa.s.で零へ行く。
その絶対値は1以下であり、優収束により期待値も零へ行く。
envelopeの三角不等式から
\[
 \Phi_{K,T}(X^n-X)
 \le\Phi_{K,T}(X^m-X^n)+\Phi_{K,T}(X^m-X)
 \le\varepsilon+\Phi_{K,T}(X^m-X).
\]
補助部分列に沿って極限を取ると右辺は $\varepsilon$ となる。
閾値 $N$ は $K$ に依存しないため、その後にsupremumを取れる。
各 $g_{n,T}$ は $[0,1]$ に属するので、幾何重みを支配列とする級数の優収束から後半を得る。
\end{proof}

この結論は同じ代表列に対するelementary testsの一様収束である。
本稿の位相はこの elementary test 族で定義するため、一般の predictable test への拡張を必要としない。

\subsection{補助位相と共通局所化}

\begin{lemma}[capped級数からの経路的有限性]
具体的な零始点分解 $X_k=N_k+A_k$ に対して
\[
 B_k(t)=\sqrt{[N_k]_t}+\Var_{[0,t]}(A_k)
\]
と置く。全整数 $n\ge0$ で $\sum_k E[1\wedge B_k(n)]<\infty$ なら、
一つの確率1の集合上で全時刻 $t\ge0$ について $\sum_k B_k(t)<\infty$ である。
\end{lemma}
\begin{proof}
各 $B_k$ は有限非負値で増加し、固定時刻で可測である。
整数 $n$ を固定すると、Tonelliから $\sum_k(1\wedge B_k(n))<\infty$ a.s.を得る。
この級数の項は零へ収束するので、経路ごとに十分大きな $k$ では $B_k(n)<1$ となる。
その尾部ではcapが外れ、有限個の初期項も有限だから $\sum_k B_k(n)<\infty$ である。
整数 $n$ の可算交叉を取り、任意の $t$ には $n=\lceil t\rceil$ と単調性を使えばよい。
\end{proof}
この補題は総和の期待値の有限性を与えるものではない。

\begin{lemma}[総和の単調性と右連続性]
同じ仮定の下で、一つの確率1の集合上で
$U_t=\sum_k B_k(t)$ は有限実数値の増加・右連続関数である。
\end{lemma}
\begin{proof}
局所有限変動な右連続関数の累積全変動は右連続である。
実際、累積全変動の右極限とその値との差は、元の関数の右jumpの絶対値であり、これは零である。
従って各 $B_k$ は右連続である。前の補題のevent上では、各時刻の級数は有限であり、
項ごとの単調性から総和の単調性も従う。
$t$ を固定し、右側近傍 $t\le s<t+1$ を考えると $0\le B_k(s)\le B_k(t+1)$ である。
右辺は $k$ について和可能なので、級数の支配収束により $s\downarrow t$ と総和を交換できる。
これが同じevent上の全時刻で右連続性を与える。
\end{proof}

\begin{lemma}[具体的な分解列の共通passage]
usual conditionsの下で、上のcapped級数の仮定を満たす分解列に対して、
総和 $U$ の強適合的c\`adl\`ag代表元を選べる。
その絶対値が $r+1$ を初めて厳密に超える時刻を $r+1$ で切ったものを $\tau_r$ とすると、
$(\tau_r)$ は一つの局所化列であり、$\tau_r\le r+1$ である。
一つの確率1の集合上で、全 $r$ と全 $t<\tau_r$ に対して
\[
 \sum_k B_k(t)\le r+1
\]
が成り立つ。
\end{lemma}
\begin{proof}
各時刻 $t$ の累積全変動は、$t$ 以下の有限グリッド変動の可算上限から復元できるため
$\mathcal F_t$-可測である。二次変分の適合性と合わせて各 $B_k(t)$ は $\mathcal F_t$-可測となり、
実数値級数の総和も強適合的である。
前の補題による共通event上の単調性から左極限が存在し、右連続性と合わせてc\`adl\`ag性を得る。
usual conditionsにより例外の零集合は $\mathcal F_0$ に属するので、そこで零へ修正しても適合性を保つ。
修正した同じ過程にc\`adl\`ag絶対値passageの局所化定理を適用する。
得られる族は単調かつ無限大へ発散し、各有限代表は対応する決定論的horizon以下である。
passageより厳密に前では $|U_t|\le r+1$ であり、元の総和との全時刻一致から結論を得る。
\end{proof}
ここでは具体的な分解列のcapped評価を入力に共通局所化を構成した。
Émery収束からの入力生成は本節後半の閉グラフによる位相比較が供給し、
停止時の境界jumpを含むprelocal cost総和の評価は以下で証明する。
そのうち、closed stoppingによるquadratic rootの境界補正は次のとおりである。

\begin{lemma}[停止したquadratic rootの境界補正]
零始点の二次変分 $[N]$ について、一つの確率1の集合上で全時刻 $t$ に対し
\[
 \sqrt{[N]_t}\le\sqrt{[N]_{t-}}+|\Delta N_t|
\]
である。零始点のc\`adl\`ag過程 $N$ の有限値停止時刻 $\tau$ に対して、
同じ二次変分を停止したcertificateの全時間rootは $\sqrt{[N]_\tau}$ に等しく、
\[
 E\sup_{t\ge0}\sqrt{[N^\tau]_t}
 \le E\bigl[\sqrt{[N]_{\tau-}}+|\Delta N_\tau|\bigr].
\]
期待値は無限でもよく、$\tau$ の決定論的有界性は要求しない。
\end{lemma}
\begin{proof}
二次変分の単調性と零初期値から $[N]_{t-}\ge0$ である。
jumpの同定により $[N]_t=[N]_{t-}+(\Delta N_t)^2$ なので、
非負な平方根を比較すれば最初の不等式を得る。
$[N]^\tau_t=[N]_{t\wedge\tau}$ は増加し、$t=\tau$ で終端値に達するため全時間supremumを同定できる。
共通event上で $t=\tau$ を代入し、非負積分の単調性を使えば期待値の評価を得る。
時刻零では左極限を初期値とする規約と $[N]_0=0$ を用いる。
\end{proof}

\begin{lemma}[停止直前成分とjumpによるprelocal cost評価]
usual conditionsの下で、零始点分解 $X=N+A$ と有限停止時刻 $\tau\le T$ に対して
\[
 \mathcal J_{\tau-}(X)
 \le E\bigl[\sqrt{[N]_{\tau-}}+|\Delta N_\tau|\bigr]
       +E[\Var(A)_{\tau-}]
\]
が成立する。$A$ は局所有限変動でよく、期待値は無限でもよい。
\end{lemma}
\begin{proof}
まず $A$ を決定論的時刻 $T$ で停止する。この過程は大域的有限変動であり、
$[0,\tau)$ で元の $A$ と一致するため、strict prefixも元のものと等しい。
累積全変動の単調性から、全 $t<\tau$ で $\Var(A)_t\le\Var(A)_{\tau-}$ である。
strict-prefix変動評価を適用すると
$\Var(A^{\tau-})_\infty\le\Var(A)_{\tau-}$ を得る。
従って $Y=N^\tau+A^{\tau-}$ は $X$ のstrict prefix上の延長であり、
右辺の二成分はその許容分解となる。マルチンゲール成分のcostには前の境界root評価を使い、
FV成分には上の全変動評価を積分して用いる。
最後に延長と分解についての下限の定義を適用すれば結論を得る。
\end{proof}
\begin{lemma}[共通passageにおける左極限の総和]
上のcapped仮定から構成した同じ共通passageについて、一つの確率1の集合上で全 $r$ に対し
\[
 \sum_k\left(\sqrt{[N_k]_{\tau_r-}}+\Var(A_k)_{\tau_r-}\right)\le r+1
\]
が成立する。
\end{lemma}
\begin{proof}
各clockの左極限は、二次変分の左極限に平方根を適用したものと累積全変動の左極限との和である。
有限部分和を固定する。$t<\tau_r$ ではこの部分和は全級数以下であり、共通passageの評価から
$r+1$ 以下となる。$\tau_r>0$ なら左からの極限を取り、有限部分和の左極限にも同じ上界を得る。
最後に有限部分集合について上限を取ると非負級数全体の評価となる。
$\tau_r=0$ では各左極限が零初期値に一致するため同じ結論となる。
\end{proof}
\begin{lemma}[共通局所化によるcostとwitnessの和可能性]
usual conditionsの下で、上のcapped仮定を満たす零始点分解列について、
同じ共通停止時刻族 $\tau_r\le r+1$ は各 $r$ で
\[
 \sum_k\mathcal J_{\tau_r-}(X_k)
 \le r+1+\sum_k E|\Delta N_{k,\tau_r}|
\]
を満たす。ここで
\[
 s(N):=\sup_{\tau\text{：有界停止時刻}}E|\Delta N_\tau|
\]
と定める。さらに $\sum_k s(N_k)<\infty$ なら、各 $r$ で形式の
exact-representation witness列を選べ、そのcost総和は
\[
 6\left(r+1+\sum_k s(N_k)+1\right)<\infty
\]
で抑えられる。停止時刻族は列全体に共通である。
\end{lemma}
\begin{proof}
停止時jumpは進行的可測過程の停止値なので可測である。
従って一分解のcost評価からその期待値を分離できる。
残る左極限clockについては、有限個の期待値の和は和の非負積分以下であり、
経路的総和上界を積分すると $r+1$ 以下となる。
有限部分集合について上限を取り、jumpの期待値を加えれば最初の不等式を得る。
各 $\tau_r$ は有界停止時刻なので、jump項は $\sum_k s(N_k)$ 以下である。
よってprelocal costが和可能となり、近似的最小分解の選択定理を適用すると、
総和が1の正の誤差列とDavis定数6からwitnessの上界を得る。
\end{proof}
このwitnessは一般分解のものであり、元の価格過程に対するactual integralの所属は意味しない。
\begin{definition}[零始点の補助cost]
零始点分解 $D=(N,A)$ とその二次変分に対し
\[
 c_{r^1}(D):=s(N)+\sum_{n=0}^\infty 2^{-n-1}
 E\left[1\wedge\left(\sqrt{[N]_n}+\Var(A)_n\right)\right],
 \qquad R(X):=\inf_{X\sim N+A}c_{r^1}(N,A)
\]
と定める。値は無限でもよい。二つの項は同じ分解から計算し、その和の下限を取る。
\end{definition}
\begin{lemma}[許容分解を持つ全過程での補助costの有限性]
usual conditionsの下で、零始点の許容分解 $X\sim N+A$ が存在するなら $R(X)\le3$ である。
逆に $R(X)<\infty$ なら許容分解が存在する。
この主張に $\mathcal J(X)<\infty$ は不要である。
\end{lemma}
\begin{proof}
与えられた $N$ に閾値1の大域DDY分解を適用し、$N=L+Q$ とする。
$L$ は零始点局所マルチンゲールで、その全時刻のjumpは共通の確率1の集合上で2以下である。
$Q$ は零始点の適合的càdlàg局所有限変動過程であるから、
$X\sim L+(Q+A)$ は同じtargetの許容分解である。
同じ $L$ の二次変分を構成する。全有界停止時刻でjumpを評価して期待値を取ると $s(L)\le2$ である。
一方、capped clockの各期待値は1以下であり、重みの総和も1なので、その級数は1以下である。
従ってこの分解のjoint costは3以下であり、下限 $R(X)$ も3以下となる。
逆方向では、許容分解が存在しなければ外側の下限が空であり $R(X)=\infty$ となる。
\end{proof}
ここで有限性を得た空間は許容分解を持つ零始点過程全体である。
主定理で必要な零始点・正則なelementary good-integratorからこの空間への所属は、
後述の分解生成とelementary-test完備性で供給する。
\begin{lemma}[大域的costから補助costへの評価]
usual conditionsの下で、拡張値のまま $R(X)\le2\mathcal J(X)$ が成立する。
従って $\mathcal J(X)<\infty$ なら $R(X)<\infty$ であり、$\mathcal J(X_k)\to0$ なら $R(X_k)\to0$ である。
\end{lemma}
\begin{proof}
分解 $D=(N,A)$ を固定し、その大域的costを
\[
 j(D):=E\left[\sup_t\sqrt{[N]_t}+\Var(A)_\infty\right]
\]
と書く。二次変分のjump条件と単調性から
$|\Delta N_\sigma|\le\sqrt{[N]_\sigma}\le\sup_t\sqrt{[N]_t}$ が
全時刻共通の確率1の集合上で成立する。期待値を取り、全有界停止時刻について上限を取ると
$s(N)\le j(D)$ を得る。一方、各時刻のclockは全時間rootと全変動の和以下なので
各capped期待値も $j(D)$ 以下である。重み $2^{-n-1}$ の総和は1であり、
$c_{r^1}(D)\le2j(D)$ となる。全分解について下限を取れば主張を得る。
無限値から実数への変換は不要である。零収束の主張は $0\le R(X_k)\le2\mathcal J(X_k)$ から従う。
\end{proof}
\begin{lemma}[補助costの三角不等式と差]
usual conditionsの下で
\[
 R(X+Y)\le R(X)+R(Y),\qquad R(-X)=R(X),\qquad
 R(X-Y)\le R(X)+R(Y)
\]
が成立する。従って $R(Z_k-X)\to0$ なら $R(Z_{k+1}-Z_k)\to0$ である。
\end{lemma}
\begin{proof}
二つの許容分解を固定して足す。局所マルチンゲールの和の二次変分を構成し、
root三角不等式と全変動の劣加法性を使うと、和のclockは二つのclockの和以下である。
$1\wedge(a+b)\le(1\wedge a)+(1\wedge b)$ を積分し、正の重みを掛けて足すと
capped項も劣加法的となる。一方、全有界停止時刻でjumpの三角不等式を積分してから
上限を取れば、jump項の劣加法性を得る。停止時jumpの可測性を用いて非負積分の和を分離する。
最後に、二項を同じ分解に保ったまま全分解と二次変分について下限を取る。
負号は二次変分と全変動を変えず、jumpの絶対値も保存するので $R(-X)=R(X)$ である。
差の評価は和と負号の結果から従う。
$Z_{k+1}-Z_k=(Z_{k+1}-X)-(Z_k-X)$ に適用して零収束を得る。
\end{proof}
\begin{lemma}[補助costとclosed stop]
任意の停止時刻 $\tau$ に対して $R(X^\tau)\le R(X)$ が成立する。
また $\sigma\le\tau$ なら $R(X^\sigma)\le R(X^\tau)$ である。
停止時刻は無限値を取ってよい。
従って $R(X_k)\to0$ なら、各項で任意に選んだ停止時刻 $\tau_k$ に対しても
$R(X_k^{\tau_k})\to0$ である。
\end{lemma}
\begin{proof}
一つの分解 $X\sim N+A$ とその二次変分を固定し、両成分を同じ $\tau$ で停止する。
任意の有限時刻 $s$ において、$s\le\tau$ なら $\Delta N^\tau_s=\Delta N_s$ であり、
$\tau<s$ なら $\Delta N^\tau_s=0$ である。これを各有界停止時刻で評価して積分すると
$s(N^\tau)\le s(N)$ を得る。
停止した二次変分は $[N]^\tau$ であり、その単調性と停止後の全変動の縮小から
各時刻のclockも元のclock以下である。従ってcapped積分と重み付き級数も縮小する。
同じ分解の二座標を合わせ、全分解について下限を取れば第一の主張を得る。
第二の主張は $(X^\tau)^\sigma=X^\sigma$ に第一の主張を適用して従う。
列については $0\le R(X_k^{\tau_k})\le R(X_k)$ を用いる。
\end{proof}
\begin{lemma}[補助costの零収束から共通局所化へ]
usual conditionsの下で $R(X_k)\to0$ なら、一つの狭義単調な部分列 $f$ と
一つのlocalizing sequence $\tau_r\le r+1$ が存在し、各 $r$ について
$X_{f(k)}$ のexact-representation witness列を選べる。そのcost総和は $6(r+4)$ 以下であり、
costは零へ収束する。
\end{lemma}
\begin{proof}
まず $\sum_k R(X_k)<\infty$ の場合を考える。総和が1の正の誤差列を使って
各下限を近似すると、一つの分解列 $D_k$ を選べて
$\sum_k c_r(D_k)\le\sum_k R(X_k)+1<\infty$ となる。
jump項の総和はこのcost総和以下である。また整数 $n$ を固定すると、
重み $2^{-n-1}>0$ を掛けたcapped期待値の総和もこのcost総和以下である。
重みが正なので各整数horizonのcapped期待値の総和も有限となる。
前の共通局所化定理をこの同じ分解列へ適用すると、各 $r$ でwitness cost総和は
$6(r+1+\sum_k R(X_k)+2)$ 以下となる。
一般の零収束列については $R(X_{f(k)})<2^{-k-1}$ となる狭義単調な部分列を選ぶ。
その補助cost総和は1以下なので上界 $6(r+4)$ を得る。
非負級数の有限性から各項の零収束も従う。
\end{proof}
Émery収束から補助costの零収束への接続は、後述の同じ分解可能過程の商上の閉グラフ定理で証明する。
大域的costの零収束 $\mathcal J(X_k)\to0$ が与えられる場合には、上の比較をこの構成へ適用できる。
同じ結論、すなわち一つの部分列、一つの共通停止時刻族、および各段階でcost総和が
$6(r+4)$ 以下かつ零収束するexact-representation witness列を得る。
\begin{lemma}[補助costの分離性]
区別不能な過程は同じ補助costを持つ。usual conditionsの下で
\[
 R(X)=0\ \Longleftrightarrow\ X\sim0,
 \qquad R(X-Y)=0\ \Longleftrightarrow\ X\sim Y
\]
が成立する。
\end{lemma}
\begin{proof}
区別不能なtargetへ同じ分解と二次変分を移してもcostは変わらず、下限も等しい。
$R(X)=0$ なら定数列 $X_k=X$ の補助cost総和は零である。
近似的最小分解を選び共通局所化へ渡すと、各 $r$ で
$\sum_k\mathcal J_{\tau_r-}(X)<\infty$ を得る。定数非負級数なので
$\mathcal J_{\tau_r-}(X)=0$ である。
まず $X$ が適合的で右連続な場合、各固定時刻 $t$ について
$|X_t|\mathbf1_{\{t<\tau_r\}}$ は可測であり、停止前の最大値評価からその期待値は零となる。
可算個の $r$ に共通の確率1の集合を取り、$\tau_r\to\infty$ を使うと $X_t=0$ を得る。
さらに可算right-dense skeleton上の等式を右連続性で全時刻へ延ばす。
一般のraw targetでは、$R(X)<1$ から許容分解を一つ取り、その和を正則な代表元として使う。
逆向きは零分解と零二次変分のcostが零であることと、区別不能性による不変性から従う。
差に適用すれば二つ目の同値も得る。
\end{proof}
\begin{lemma}[補助costの和可能性と経路的一様収束]
usual conditionsの下で $\sum_k R(X_k)<\infty$ なら、一つの確率1の集合上で全有限 $T$ に対し
\[
 \sum_k\sup_{t\le T}|X_k(t)|<\infty
\]
である。従って各有限horizonで一様に零へ収束する。
特に $R(X_k)\to0$ なら、この結論を満たす一つの狭義単調な部分列を選べる。
\end{lemma}
\begin{proof}
まず各過程が適合的で右連続である場合を考える。
共通停止時刻 $\tau_r$ を固定し、可算right-dense skeletonのうち $t<\tau_r$ の点における
$|X_k(t)|$ の上限を $G_{k,r}$ とする。これは可測であり、停止前の最大値評価から
$E G_{k,r}\le6\mathcal J_{\tau_r-}(X_k)$ である。
右辺を $k$ について足すと有限なので、Tonelliにより $\sum_kG_{k,r}<\infty$ が確率1で成立する。
可算個の $r$ に共通の集合を取り、任意の有限 $T$ について $T<\tau_r$ となる $r$ を選ぶ。
各 $t\le T$ を停止前のskeleton点で右から近似し、右連続性を使うと
$\sup_{t\le T}|X_k(t)|\le G_{k,r}$ を得る。これを足せば結論となる。
一般のraw target列では、近似的最小分解の和を正則な代表元として使い、
全rowと全時刻に共通の区別不能性により同じ結論を移す。
最後に零収束列から $\sum_kR(X_{f(k)})\le1$ となる部分列を選ぶ。
\end{proof}
\begin{lemma}[補助costの収束はucp収束を含意する]
usual conditionsの下で $R(X_k-X)\to0$ なら、全有限 $T$ と全 $\varepsilon>0$ に対し
\[
 P\{\exists t\le T:\ |X_k(t)-X(t)|>\varepsilon\}\longrightarrow0
\]
である。raw targetの可測性や経路正則性は追加仮定しない。
\end{lemma}
\begin{proof}
まず補助costが和可能な列を考える。各項は許容分解を持つので、その正則な和の
factorial envelopeを使い、元のtargetの有限horizon最大値がa.e.可測であることを得る。
前の補題から最大値はa.e.有限で零へ収束するため、その実数値版は確率収束する。
最大値の有限性を使って実数値版の確率尾部を拡張非負実数値の最大値の確率尾部と同定する。
一般の零収束列については、任意の部分列から補助costが和可能な部分列を再抽出できる。
従って確率尾部の数列にも零へ収束する再抽出があり、部分列判定から元の数列が零へ収束する。
最後に $X_k-X$ に適用する。ある時刻で誤差が $\varepsilon$ を越える事象は、
その最大値が $\varepsilon$ 以上である事象に含まれる。
\end{proof}
従って $R(X_k)\to0$ なら、項ごとに任意の停止時刻 $\tau_k$ を選んでも
$X_k^{\tau_k}\to0$ がucpで成立する。これはclosed stopによる補助costの縮小を
上の確率尾部の評価へ適用した帰結である。

\begin{proposition}[分解可能過程の商上の距離とCauchy列]
usual conditionsの下で、許容分解を持つ零始点過程全体を $\mathcal D_0$ とする。
$d(X,Y):=R(X-Y)$ は有限値の擬距離であり、その分離商は区別不能性による商と一致する。
商上の $X_k\to X$ は $R(X_k-X)\to0$ と同値であり、ucp収束を含意する。
さらに商上で $(X_k)$ がCauchyなら、一つの狭義単調な $f$ が存在し、
\[
 \sum_k R(X_{f(k+1)}-X_{f(k)})\le1
\]
となる。同じ差分列に対して、一つのlocalizing sequence $\tau_r\le r+1$ と、
各 $r$ でcost総和が $6(r+4)$ 以下かつ有限なexact-representation witness列を選べる。
\end{proposition}
\begin{proof}
二つの許容分解の差も許容分解であるから、DDYによる有限性の結果を距離に適用できる。
対称性と三角不等式は補助costの代数から従い、距離零と区別不能性の同値性から商の同定を得る。
商の距離は代表元の $R(X-Y)$ に等しいので、距離収束の同定とucpへの含意も従う。
Cauchy性から $b_n\to0$ を満たす拡張非負実数列を選び、$i,j\ge n$ なら
$d(X_i,X_j)\le b_n$ とできる。$b_{f(k)}<2^{-k-1}$ となる狭義単調な $f$ を選べば
$f(k+1)\ge f(k)$ により隣接差のcost総和は1以下となる。
この同じ差分列に補助costの和可能性からの共通局所化を適用する。
\end{proof}
\begin{lemma}[同じ部分列の局所一様極限]
隣接差について $\sum_k R(Y_{k+1}-Y_k)<\infty$ なら、一つの過程 $Y$ が存在し、
一つの確率1の集合上で全有限 $T$ に対し $Y_k\to Y$ が $[0,T]$ 上一様に成立する。
従って前の命題のCauchy列には、共通局所化witnessを持つ同じ部分列に対して、この極限を選べる。
\end{lemma}
\begin{proof}
補助costの和可能性から、共通の確率1の集合上で全有限 $T$ に対して
\[
 \sum_k q_k(T)<\infty,\qquad
 q_k(T):=\sup_{t\le T}|Y_{k+1}(t)-Y_k(t)|
\]
を得る。各項は有限であり、各固定時刻の隣接差の距離級数はこの級数で抑えられる。
実数の完備性から各時刻の極限を得る。この点ごとの極限を $Y_t$ と定めると、
和可能な一様step上界 $q_k(T)$ によって $[0,T]$ 上の一様収束へ上げられる。
$Y_t$ の定義はhorizonに依存しない。前の命題が選んだ部分列の差分costは和可能なので、
その同じ部分列にこの結論を適用する。
\end{proof}
この極限はここではraw processとして得られる。以下で大域成分級数を使って許容分解を構成し、
三座標のtail評価から商上の補助距離での収束と完備性を示す。
\begin{lemma}[固定停止段階のマルチンゲール成分の和]
true martingale列 $(M_k)$ と非負可測関数 $(g_k)$ が
$|M_k(t)|\le g_k$ を全時刻で満たし、$\sum_k E g_k<\infty$ とする。
このとき $M_t:=\sum_k M_k(t)$ はstrongly adaptedなtrue martingaleである。
従って上で得た共通局所化の各段階 $r$ で、同じwitnessの成分を零始点に正規化した後の
\[
 M^{(r)}_t:=\sum_k (N_k-N_k(0))_{t\wedge\tau_r}
\]
もtrue martingaleとなる。
\end{lemma}
\begin{proof}
Tonelliにより $G:=\sum_k g_k$ は可積分であり、確率1で有限である。
その同じ集合上で全時刻の $\sum_k |M_k(t)|$ が有限となる。
各固定時刻で級数の和は可測であるから、$M$ はstrongly adaptedである。
有限部分和はtrue martingaleであり、その絶対値は可積分な $G$ 以下である。
支配収束によるmartingale極限定理を適用すれば第一の主張を得る。
witnessへの適用では、零始点正規化したclosed-stopped成分がtrue martingaleであることを使う。
停止時刻がhorizon以下なので、有限horizonのenvelopeは停止成分を全時間で支配する。
正規化のcostは元のcostの2倍以下であり、witness costの和可能性から必要なenvelope期待値の
和可能性が従う。
\end{proof}
\begin{lemma}[中心化した有限変動成分の級数]
零初期値の実数値経路列 $(a_k)$ が $\sum_k\Var(a_k)<\infty$ を満たすなら、
部分和は $a:=\sum_k a_k$ へ全時間で一様収束し、$a$ は有限変動であり、
\[
 \Var\left(a-\sum_{k<n}a_k\right)\longrightarrow0
\]
となる。従って固定段階の和可能なwitness列に対し
$a_k(t):=A_k^{\tau-}(t)-A_k^{\tau-}(0)$ とすれば、同じ三性質が共通の確率1の集合上で成立する。
\end{lemma}
\begin{proof}
零初期値から $|a_k(t)|\le\Var(a_k)$ であり、和可能な上界による級数の一様収束定理を使える。
部分和の隣接差の全変動は $\Var(a_k)$ であるから、有限変動極限定理と
全変動距離での収束定理を適用して残る二性質を得る。
witnessへの適用では、定数を引いても変動は変わらず、strict prefixは停止時刻以後一定である。
停止時刻がhorizon以下なので、全時間の変動はwitnessの有限horizon変動以下となる。
cost総和の有限性とTonelliから、これらの変動の総和は確率1で有限である。
中心化は、変動から制御できない初期定数の級数を取り除くために必要である。
\end{proof}
この結果は、同じCauchy列の部分列と同じ共通停止時刻族の各段階へ適用済みである。
以下では大域分解列を保持する経路を使い、段階ごとの成分の貼り合わせを避けて極限の分解を構成する。

\begin{lemma}[大域分解を保持した共通成分評価]
上の商空間のCauchy列に対して、一つの狭義増加列 $f$ と、実際の隣接差
$X_{f(k+1)}-X_{f(k)}=N_k+A_k$ の大域分解列を選べる。
同じ分解の補助costの総和は2以下であり、一つの局所化列 $\tau_r\le r+1$ に対して
\[
 \sum_k\left(E\sqrt{[N_k]_{\tau_r}}+
 E\Var(A_k^{\tau_r-})\right)\le r+3
\]
が成立する。全変動は全時間で取る。
\end{lemma}
\begin{proof}
隣接差の補助距離の総和を1以下にする部分列を選び、下限の近似によって一つの分解列を
選ぶと、そのcost総和は2以下になる。この同じ列の共通clockを用いた停止では、
停止直前のrootと累積変動の総和の期待値が $r+1$ 以下となる。
各項について
\[
 \sqrt{[N_k]_{\tau_r}}\le
 \sqrt{[N_k]_{\tau_r-}}+|\Delta N_{k,\tau_r}|,
 \qquad
 \Var(A_k^{\tau_r-})\le\Var_{[0,\tau_r)}(A_k).
\]
境界jumpの期待値は同じ分解のbounded-stopping-jump cost以下であり、
その総和は2以下である。非負積分の有限和の評価を上限へ移して、目的の評価を得る。
停止した二次変分は増加し、停止時刻以後一定なので、停止時のrootは全時間のroot上限に等しい。
\end{proof}
\begin{lemma}[同じ大域マルチンゲール成分の和]
前の補題で選んだ成分について、$M_t:=\sum_kN_k(t)$ と定めると、各 $M^{\tau_r}$ は
true martingaleであり、$M$ 自身は零始点のlocal martingaleである。
\end{lemma}
\begin{proof}
一つの $r$ を固定する。停止二次変分の全時間root上限は $\sqrt{[N_k]_{\tau_r}}$ に等しい。
Davis評価から
\[
 \sum_k E\sup_t|N_k(t\wedge\tau_r)|
 \le 6\sum_k E\sqrt{[N_k]_{\tau_r}}<\infty
\]
を得る。$\tau_r\le r+1$ なので、各停止成分の有限horizonの可測envelopeが全時間を支配する。
その可積分性によって各停止成分はtrue martingaleとなり、上の支配されたmartingale級数の
補題を適用できる。停止と点ごとの級数は同じ時点の値を評価するので、得られる和は
$M^{\tau_r}$ そのものである。各 $N_k(0)=0$ から $M_0=0$ であり、共通局所化列を使えば
局所化の零時刻indicatorを除去して $M$ のlocal martingale性を得る。
\end{proof}
\begin{lemma}[マルチンゲール級数の正則な代表元]
同じ $M=\sum_kN_k$ は、零始点でstrongly adaptedかつ全経路càdlàgな
local martingaleと区別不能である。
\end{lemma}
\begin{proof}
各共通停止段階で用いたenvelopeの期待値総和は有限なので、Tonelliにより経路ごとの
envelope総和も確率1で有限となる。この同じ上界によって停止成分の有限部分和は全時間で
一様収束する。各部分和はcàdlàgであり、一様極限も右連続で左極限を持つ。
全停止段階の共通の確率1の集合を取る。停止時刻族は無限大へ達するため、各時刻 $t$ に対し
ある段階の停止時刻が $t+1$ より大きくなる。大域和とその停止は $[0,t+1)$ で同じ値を持ち、
大域和の右連続性と左極限をこの停止段階から得る。
各固定時刻の級数和は可測である。càdlàgでない経路の零集合を初期時刻可測な集合として
その上で零に置き換えると、適合性・零初期値・区別不能性を保存する。
区別不能性による移送定理を使い、この同じ代表元へlocal martingale性を移す。
\end{proof}
\begin{lemma}[同じ大域有限変動成分の和]
同じ分解列について $A_t:=\sum_k A_k(t)$ と定めると、一つの確率1の集合上で
$A$ は局所有限変動であり、全ての有限 $T$ に対して
\[
 \sum_{k<n}A_k\longrightarrow A\quad\text{一様に }[0,T]\text{ 上で},
 \qquad
 \Var_{[0,T]}\left(A-\sum_{k<n}A_k\right)\longrightarrow0
\]
が成立する。
\end{lemma}
\begin{proof}
capped clockの評価から、一つの確率1の集合上で、全ての $T$ に対し
$\sum_k\Var_{[0,T]}(A_k)<\infty$ となる。
$T$ を固定し、$a_k(t):=A_k(t\wedge T)$ と置く。
その全時間の変動は $\Var_{[0,T]}(A_k)$ 以下であり、$a_k(0)=0$ である。
零始点の変動級数定理により、$\sum_k a_k$ は全時間で有限変動であり、
有限部分和は一様収束し、全時間の変動距離でも収束する。
$[0,T]$ 上では $a_k=A_k$ なので、一様収束と変動距離の収束を元の成分へ移せる。
任意のコンパクト区間はある $[0,T]$ に含まれるため、$A$ の局所有限変動性も従う。
\end{proof}
\begin{lemma}[有限変動級数の正則な代表元]
前の有限変動級数は、零始点でstrongly adapted、全経路càdlàgかつ局所有限変動の過程と
区別不能である。
\end{lemma}
\begin{proof}
各時刻で可測な関数の級数を取るため、和はstrongly adaptedである。
一つの確率1の集合上で、有限部分和は全ての有限区間で一様収束する。
各有限部分和は右連続であり、時刻 $t$ の右側の近傍を $[0,t+1]$ に制限して
一様極限の定理を使えば、和の右連続性を得る。
局所有限変動性と右連続性が同時に成立しない経路を一つの零集合にまとめる。
usual conditionsによりこの集合は初期時刻で可測なので、その上で和を零に置き換えても
適合性と区別不能性を保存する。零初期値も保たれる。
修正した全経路は右連続かつ局所有限変動であり、後者から左極限の存在が従う。
\end{proof}
この構成では停止段階ごとに分解を選び直さない。adaptedな有限変動成分を持つ分解では、
和の一致から各成分の一致は従わないためである。
\begin{lemma}[分解可能過程空間内の極限]
補助距離についてCauchyな列 $(X_k)$ には、一つの狭義増加列 $f$ と許容分解を持つ過程 $Y$ が存在し、
一つの確率1の集合上の全有限区間で $X_{f(n)}\to Y$ が一様に成立する。
同じ部分列の点ごとの極限も $Y$ と区別不能であり、許容分解を持つ。
\end{lemma}
\begin{proof}
上で構成した同じ大域成分の正則な和を $M,A$ とし、$Y:=X_{f(0)}+M+A$ と置く。
基準項の許容分解と $(M,A)$ の分解を加えると $Y$ の許容分解を得る。
全ての隣接差の分解、両成分の代表元との一致、停止後の一様収束、FV部分和の局所一様収束を
一つの確率1の集合上で用いる。有限区間 $[0,T]$ に対して $\tau_r>T$ となる段階を選べば、
停止後のmartingale部分和はその区間で元の部分和と一致する。従って両成分の部分和は
それぞれ $M,A$ へ一様収束する。一方、各有限和で
\[
 X_{f(0)}+\sum_{k<n}(N_k+A_k)=X_{f(n)}
\]
であるから、$Y$ への一様収束が従う。
各時刻の極限の一意性をこの同じ集合上で使うと、点ごとの極限は $Y$ と全時刻で一致し、
区別不能性によって許容分解を移せる。
\end{proof}
\begin{lemma}[級数極限の停止jump評価]
許容分解のmartingale成分列 $U_n$ が左極限を持つ過程 $U$ へa.e.に局所一様収束するなら、
bounded-stopping-jump cost $s$ に対して $s(U)\le\liminf_n s(U_n)$ となる。
従って上で構成した同じ正則なmartingale和は
\[
 s(M)\le\sum_k s(N_k)
\]
を満たす。
\end{lemma}
\begin{proof}
有界停止時刻 $\sigma\le T$ を固定する。$T$ で停止した経路は全時間で一様収束し、
左極限の収束定理によって $\Delta U_n(\sigma)\to\Delta U(\sigma)$ となる。
各項のjumpは可測であり、Fatouから
$E|\Delta U(\sigma)|\le\liminf_n E|\Delta U_n(\sigma)|\le\liminf_n s(U_n)$ を得る。
最後の上界は $\sigma$ に依存しないので、有界停止時刻全体の上限を取れる。
級数には、元の分解を有限個加えた分解を使う。有限和のjump costは各項のcostの和以下であるから、
下半連続性を適用して目的の評価を得る。構成済みの $M$ への適用では、共通停止後の一様収束と
停止時刻族のexhaustivityにより、元の部分和の局所一様収束を使える。
\end{proof}
\begin{lemma}[成分級数のtail評価]
上と同じ正則な和 $M,A$ と元の分解列に対して、任意の $n$ で
\[
 s\left(M-\sum_{k<n}N_k\right)\le\sum_{k=0}^{\infty}s(N_{k+n})
\]
である。従って、この左辺は $n\to\infty$ で零に収束する。また各有限horizon $T$ で
\[
 E\left[\min\left\{\operatorname{Var}_{[0,T]}
       \left(A-\sum_{k<n}A_k\right),1\right\}\right]\longrightarrow0
\]
である。
\end{lemma}
\begin{proof}
$n$ を固定し、元の分解列を $n$ だけずらす。そのmartingale成分の部分和は
$M-\sum_{k<n}N_k$ へ同じ確率1の集合上で局所一様収束する。
前補題をこの列に適用すると最初の不等式を得る。
$\sum_k s(N_k)<\infty$ なので、右辺は零へ収束する。
FV成分については、既に得た経路ごとの全変動の零収束を用いる。
正則な和と各部分和は強適合的で右連続であるため、差の有限区間全変動は可測である。
同じ代表元への区別不能性を使い、定数1による支配収束で期待値の収束を得る。
\end{proof}
\begin{lemma}[同じマルチンゲール級数のquadratic-root tail]
上の正則な和を $M$ とし、$U_n=M-\sum_{k<n}N_k$ とおく。
元のroot期待値総和が有限である共通有界停止時刻 $\tau$ について
\[
 E\sqrt{[U_n]_\tau}
 \le42\sum_{k\ge n}E\sqrt{[N_k]_\tau}\longrightarrow0
\]
が成立する。$[U_n]$ は各tailの内在的二次変分として構成でき、停止時刻の選択には依存しない。
\end{lemma}
\begin{proof}
まず有限 $T$ 上で部分和が局所マルチンゲール $U$ へ一様収束する場合を考える。
非負のenvelope級数が無限となるpathも許して、部分和の極限から
\[
 U_T^*\le\sum_k (N_k)_T^*
\]
を得る。実際、各時刻で絶対値の有限和評価を極限へ移し、全時刻の上限を取ればよい。
逆Davis評価、Tonelli、順方向のDavis評価から
\[
 E\sqrt{[U]_T}\le7E U_T^*
 \le7\sum_k E(N_k)_T^*\le42\sum_k E\sqrt{[N_k]_T}
\]
となる。
次に $N_k^\tau$ の和へ適用する。共通停止後のenvelope期待値総和は有限なので、
その部分和は全時間でほとんど確実に一様収束する。
元の和から最初の $n$ 項を引くと、同じ確率1の集合上で
残りの部分和は $U_n^\tau$ へ一様収束する。
$\tau\le T$ を満たす決定論的horizonを取り、停止整合性を使うと主張の評価を得る。
右辺は非負の可算和のtailであるから零へ収束する。
\end{proof}

\begin{lemma}[capped rootの共通局所化の除去]
$\tau_r\to\infty$ がほとんど確実に成立し、各 $r$ で
$E\sqrt{[U_n]_{\tau_r}}\to0$ とする。このとき各有限 $T$ で
\[
 E\bigl[\min\{\sqrt{[U_n]_T},1\}\bigr]\longrightarrow0.
\]
\end{lemma}
\begin{proof}
二次変分の単調性から
\[
 \min\{\sqrt{[U_n]_T},1\}
 \le\sqrt{[U_n]_{\tau_r}}+\boldsymbol1_{\{\tau_r\le T\}}.
\]
右辺第二項の期待値はexhaustionと支配収束により $r\to\infty$ で零へ行く。
まずこれを小さくする $r$ を固定し、その後 $n\to\infty$ とすればよい。
\end{proof}

\begin{theorem}[零始点分解可能過程の補助距離の完備性]
usual conditionsの下で、許容分解を持つ全零始点過程の区別不能性商は、
距離 $d(X,Y)=R(X-Y)$ に関して完備である。
\end{theorem}
\begin{proof}
商のCauchy列から速い部分列 $X_{f(k)}$ を選び、その実際の隣接差の分解 $(N_k,A_k)$ を保持する。
上で構成した正則な成分和を $M,A$ とし、$X=X_{f(0)}+M+A$ とおく。
この $X$ は許容分解を持つ。同じ分解から最初の $n$ 項を引けば、
$X-X_{f(n)}$ のmartingale成分とFV成分はそれぞれ
\[
 U_n=M-\sum_{k<n}N_k,\qquad B_n=A-\sum_{k<n}A_k
\]
となる。sourceの同定には元の隣接差のtelescopingを用いる。
jump tail評価により $s(U_n)\to0$ であり、各有限 $T$ で
\[
 E\min\{\sqrt{[U_n]_T},1\}\to0,\qquad
 E\min\{\operatorname{Var}_{[0,T]}(B_n),1\}\to0
\]
も得た。$\min\{a+b,1\}\le\min\{a,1\}+\min\{b,1\}$ を使うと、同じ分解の
capped clock期待値が各整数horizonで零へ行く。
各期待値は1以下なので、幾何重み $2^{-j-1}$ による可算和にも支配収束を適用できる。
従って同じ分解の $r^1$ costが零へ収束し、その下限である $R(X-X_{f(n)})$ も零へ行く。
負号不変性から元の部分列は商上で $X$ へ収束する。
最後にCauchy列に収束部分列があれば列全体も同じ極限へ収束することを使う。
商の任意の列を代表元へ持ち上げてこの議論を適用できるので、完備性を得る。
\end{proof}

\begin{lemma}[スカラー倍の連続性]
許容分解を持つ零始点過程について
\[
 R(cX)\le\max\{1,|c|\}\,R(X)
\]
である。さらに、$a_k\to a$ かつ $R(X_k-X)\to0$ なら
$R(a_kX_k-aX)\to0$ であり、同じ分離商上のスカラー倍は連続である。
\end{lemma}
\begin{proof}
分解 $X=N+A$ と二次変分 $[N]$ を固定する。
$cX=cN+cA$ は許容分解で、$[cN]=c^2[N]$ である。
後者はjumpの二乗の恒等式と、$c^2(N^2-[N])$ の局所マルチンゲール性から従う。
変動の評価とjumpの斉次性から、clockとjump costはそれぞれ
$|c|B(t)$ と $|c|s(N)$ 以下である。
$1\wedge(|c|b)\le\max\{1,|c|\}(1\wedge b)$ を適用し、
同じ分解で評価してから下限を取ると最初の不等式を得る。

零へのスカラー収束にはこの不等式だけでは足りない。
DDY分解を用いて、固定した $X$ の有限jump costを持つ分解 $E=N+A$ を選ぶ。
$a_k\to0$ なら $|a_k|s(N)\to0$ であり、各有限horizonでは
有限のclockに対して $1\wedge(|a_k|B(t))\to0$ a.s.である。
1による支配収束と幾何重みの可算和への支配収束から
同じ分解のcostが零へ行き、$R(a_kX)\to0$ を得る。
一般のjoint limitは
\[
 a_kX_k-aX=a_k(X_k-X)+(a_k-a)X
\]
と三角不等式、最初の評価、固定過程に対する零収束から従う。
距離空間では逐次連続性は連続性と同値であり、連続なスカラー倍は分離商へ降りる。
\end{proof}


\clearpage
\section{Émery位相と同じ係数による積分実現}
\label{sec:emery-realization}
本付録ではgood-integrator性とÉmery位相の完備性を結び、M/FV両積分を一つの予測可能係数で同時に実現する。
\subsection{Good-integrator性と位相比較}
\label{sec:emery-completeness}
\begin{lemma}[局所有限変動過程のgood-integrator性]
強適合的で右連続な実数値過程 $A$ が全経路で局所有限変動を持つとする。
有限測度の下で $A$ はelementary good integratorである。
\end{lemma}
\begin{proof}
予測可能elementary integrands $H_n$ が全時間・全標本上で一様に零へ行くとする。
固定した $T$ で $A$ を停止すると、全時間軸上で有限変動の経路を得る。
時刻 $T$ のelementary gainは変わらず、各経路で
\[
 (H_n\mathbin{\cdot}A)_T
 =\int 1_{(0,T]}(t)H_n(t)\,dA^T_t
\]
というStieltjes積分で表せる。被積分関数は各時刻で零へ行き、十分大きな $n$ では
絶対値が1以下である。経路ごとの全変動測度は有限なので、ベクトル測度の支配収束により
gainは各標本で零へ行く。
強適合性と右連続性からsourceはprogressiveであり、elementary gainは可測である。
従って有限測度下で確率収束を得る。
strategyのblock数の上界や変動の期待値の有限性は使わない。
\end{proof}

零始点分解 $X=N+A$ にこれを適用すると、分解からgood-integrator性への方向は
$N$ のelementary good-integrator性へ還元される。
elementary gainは価格に線形であり、同じ全時刻の零集合の外で分解に沿って加算されるため、
raw targetが例外集合上で非正則でも結論は保たれる。
平方可積分な真のマルチンゲールには次の評価を適用し、DDY 分解と停止で一般の場合へ戻す。

\begin{lemma}[elementary martingale積分の表現に依存しない評価]
右連続な真のマルチンゲール $M$ と有限horizon $T$ に対して $M_T\in L^2$ とする。
予測可能elementary integrand $H$ が $|H|\le B$ を全時間・全標本で満たし、$B\ge0$ なら
\[
 \|(H\mathbin{\cdot}M)_T\|_2
 \le B\sqrt{E(M_T-M_0)^2}.
\]
従って各時刻で平方可積分な右連続マルチンゲールはelementary good integratorである。
\end{lemma}
\begin{proof}
$M$ を $T$ で停止し、chronological factorial grid上で左端値による離散積分を作る。
標本化されたsourceは平方可積分マルチンゲールであり、離散積分の等長性から
その二乗期待値は $B^2E(M_T-M_0)^2$ 以下である。
gridの終点は $T$ を覆い、停止sourceの終値は $M_T$、始点は $M_0$ である。
右連続性からgrid積分は各経路でelementary gainへ収束する。
$L^2$ seminormのFatou評価で極限へ上界を渡す。

$H_n$ が一様に零へ行く場合、固定した $T$ について
$C=\sqrt{E(M_T-M_0)^2}$ とおく。
任意の $\varepsilon>0$ に対して、十分大きな $n$ では
$|H_n|\le\varepsilon/(C+1)$ なのでgainの $L^2$ normは $\varepsilon$ 以下である。
gainの可測性と $L^2$ 収束から確率収束を得る。
この評価にはstrategyのblock数も係数の絶対値和も現れない。
\end{proof}

\begin{lemma}[good-integrator性の局所化解除]
有限測度の下で $S$ は強適合かつ右連続とする。
停止時刻列 $\tau_r\to\infty$ がほとんど確実に成り立ち、各 $S^{\tau_r}$ が
elementary good integratorなら、$S$ もそうである。
停止時刻列の単調性は不要である。
\end{lemma}
\begin{proof}
固定した $T$ に対し $P(\tau_r\le T)\to0$ である。
これは事象の指示関数の概収束と有限測度から従う。
$\tau_r>T$ 上では、elementary gainに現れる全ての評価時刻が $T$ 以下なので、
停止前後のgainは等しい。従って
\[
 P\bigl(|(H_n\mathbin{\cdot}S)_T|\ge\varepsilon\bigr)
 \le P(\tau_r\le T)
 +P\bigl(|(H_n\mathbin{\cdot}S^{\tau_r})_T|\ge\varepsilon\bigr).
\]
先に $r$ を固定し、次に $n\to\infty$ とすれば結論を得る。
\end{proof}

\begin{proposition}[分解可能過程のgood-integrator性]
通常の条件を満たす確率空間上で、零始点の許容分解 $X=N+A$ を持つ過程は
elementary good integratorである。
\end{proposition}
\begin{proof}
DDY 分解を $N=L+Q$ と書く。$L$ は零始点の局所マルチンゲールで
ジャンプが $2$ 以下、$Q$ は強適合な局所有限変動過程である。
共通局所化は exhaustive な停止列 $\tau_r$ を与え、$L^{\tau_r}$ は
真のマルチンゲールかつ決定論的定数で有界となる。
従って各時刻で平方可積分であり、上の積分評価と局所化解除から $L$ はgood integratorである。
有限変動成分 $Q+A$ を加え、区別不能なraw targetへ移送して $X$ の結論を得る。
一般の $N$ 自体が局所平方可積分であることは使わない。
\end{proof}

\begin{lemma}[good-integrator性と有限停止]
有限値停止時刻 $\tau$ に対し、$S$ がelementary good integratorなら $S^\tau$ もそうである。
\end{lemma}
\begin{proof}
予測可能elementary strategyの停止を、元のstrategyから予測可能な停止後tailを引いて定める。
そのintegrandは $1_{\{t\le\tau\}}H_t$ であるから、一様零収束は保たれる。
有限和の各評価時刻について最小値を入れ替えると
$(H\mathbin{\cdot}S^\tau)_T=((1_{\{t\le\tau\}}H)\mathbin{\cdot}S)_T$ となる。
この同一性に元の可測性とgood-integrator性を適用する。
\end{proof}

\begin{proposition}[有界ジャンプgood integratorの停止分解]
通常の条件の下で、$X$ は零始点、強適合、càdlàgなelementary good integratorとする。
$c\ge0$ に対し $|\Delta X_t|\le c$ が全時刻共通の零集合の外で成り立つなら、
\[
 \sigma_n=\inf\{t:|X_t|>n+1\}\wedge(n+1)
\]
はexhaustiveな停止列であり、各 $X^{\sigma_n}$ は上界 $n+1+c$ の有界sourceとなる。
従って有界sourceの分解定理（定理\ref{thm:bounded-source}）から、各座標のglobal special decompositionを得る。
\end{proposition}
\begin{proof}
exhaustivenessはcàdlàg経路のコンパクト区間上の有界性から従う。
初期値は閾値以下であり、停止直前までの上界と停止時のjump boundから
$|X^{\sigma_n}_t|\le n+1+c$ となる。強適合性、右連続性、左極限は停止で保たれ、
good-integrator性は前補題から従う。これらが定理\ref{thm:bounded-source}の全入力を供給する。
\end{proof}

\begin{proposition}[一般good integratorの大ジャンプ除去]
通常の条件の下で、零始点・強適合・càdlàgなgood integrator $X$ と $c>0$ に対し、
\[
 J_t=\sum_{0<s\le t}\Delta X_s1_{\{|\Delta X_s|>c\}},\qquad Y=X-J
\]
を定める。$J$ は零始点の強適合càdlàg局所有限変動過程であり、$Y$ は零始点の
強適合càdlàg good integratorである。全ての経路・時刻で $|\Delta Y_t|\le c$ が成り立つ。
従って前命題は $Y$ に適用でき、各停止段階のglobal special decompositionを与える。
\end{proposition}
\begin{proof}
有限horizon $T$ の有限大ジャンプ和を $J^T$ とする。
horizonの整合性から、$t\le T$ なら $J^t_t=J^T_t$ である。
そこで $J_t=J^t_t$ と置く。固定時刻での強適合性は $J^t$ から従い、
右連続性は $t$ より大きなhorizon上の一致から、左極限と局所有限変動は
各コンパクト区間上の一致から従う。零始点性も有限和から従う。
$t>0$ では $Y$ と有限horizon $t$ の残差が時刻 $t$ まで一致するので、
左極限を含む一致により残差のジャンプ上界を使える。時刻 $0$ のジャンプは零である。
局所有限変動過程のgood-integrator性と、elementary積分の線形性・確率収束の減法安定性から
$X-J$ のgood-integrator性を得る。前命題の全入力がこれで揃う。
\end{proof}
$J$ の補償過程はこの還元では必要なく、最終的に求める分解の有限変動成分は適合でよい。

\begin{theorem}[正則零始点good integratorの許容分解]\label{thm:regular-gi-decomposition}
通常の条件を満たす確率空間上で、零始点・強適合・càdlàgな過程 $X$ が
elementary good integratorであることと、許容分解 $X=N+B$ を持つことは同値である。
ここで $N$ は零始点のcàdlàg local martingale、$B$ は零始点の強適合càdlàg局所有限変動過程である。
\end{theorem}
\begin{proof}
分解からgood integratorへの向きは既に示した。逆向きでは、前命題で $c=1$ とし、
$X=Y+J$ と分ける。有界な停止source $Y^{\sigma_n}$ のspecial分解を $M^n+A^n$ とする。
この分解は初期値が正規化されているとは限らないので、
\[
 \widehat M^n=(M^n-M^n_0)^{\sigma_n},\qquad
 \widehat A^n=(A^n-A^n_0)^{\sigma_n}
\]
とする。初期値の大域的可積分性を仮定せずにlocal martingaleは中心化できる。
予測可能な初期値の時間一定延長と有限停止は予測可能性を保つ。
各経路で有限時刻に停止するため、$\widehat A^n$ は全時間上で有限変動である。
停止sourceの零始点性と停止の冪等性から、なお
$Y^{\sigma_n}=\widehat M^n+\widehat A^n$ が区別不能性の意味で成立する。

$\rho=\sigma_n\wedge\sigma_m$ で二分解を再停止する。その有限変動成分の差は、
二つの利得が同じであることからマルチンゲール成分の差と区別不能であり、local martingaleとなる。
零始点・予測可能・右連続・全時間有限変動でもあるため、rigidityから零である。
従って有限変動成分もマルチンゲール成分も共通停止上で区別不能となる。

local martingaleのgluingと、予測可能有限変動過程のgluingをこの同じ列に適用する。
得られる $M,A$ は各 $\sigma_n$ 上で対応する成分と一致し、停止列のexhaustivenessにより
$Y=M+A$ が全時刻共通の零集合の外で成立する。$M$ は強適合càdlàg local martingaleであり、
$A$ は予測可能・右連続・局所有限変動である。後者は左極限も持つ。
gluing後の初期値を再中心化し、
\[
 N=M-M_0,\qquad B=A-A_0+J
\]
とする。$Y_0=0$ より $M_0+A_0=0$ が概至る所成立し、$X=N+B$ を得る。
$J$ は適合局所有限変動過程なので $B$ もそうである。$B$ の予測可能性は結論しない。
\end{proof}

\begin{proposition}[一様elementary-test Cauchy評価の極限への移送]
強発展可測で右連続な過程列 $X_n$ と過程 $Y$ について、$X_n\to Y$ がucpで成り立つとする。
unit-bounded predictable elementary test $J$ に対し
\[
 \bar e_T^J(U,V)=\mathbb E\!\left[1\wedge\sup_{t\le T}|(J\cdot U)_t-(J\cdot V)_t|\right]
\]
と置く。各 $T,\varepsilon>0$ について、ある $N$ が存在し、全ての $m,n\ge N$ と全ての $J$ で
$\bar e_T^J(X_m,X_n)\le\varepsilon$ なら、$\sup_J \bar e_T^J(X_n,Y)\to0$ である。
逆に、この一様収束は前述のCauchy条件を含意する。
\end{proposition}
\begin{proof}
固定した $J$ について、有限和評価は、test積分のcapped最大誤差を
$\max(8\operatorname{length}(J),1)$ 倍の元の過程のcapped最大誤差で抑える。
この定数は固定testにのみ使用する。ucp収束と $[0,1]$ 内の有界性から、
$\bar e_T^J(X_m,Y)\to0$ を得る。

$T,\varepsilon$ に対しCauchy条件の $N$ を先に選び、$n\ge N$ と $J$ を固定する。
三角不等式から、全ての $m\ge N$ に対し
\[
 \bar e_T^J(X_n,Y)\le \bar e_T^J(X_m,X_n)+\bar e_T^J(X_m,Y)
 \le\varepsilon+\bar e_T^J(X_m,Y).
\]
$m\to\infty$ として $\bar e_T^J(X_n,Y)\le\varepsilon$ を得る。
$N$ は $J$ に依存しないので、ここで初めて全testの上限を取れる。
逆向きは $Y$ を介する三角不等式と $\varepsilon/2$ の誤差評価による。
\end{proof}
この命題は正則なucp極限を入力とする。その存在は以下で構成する。
極限のgood-integrator性と、同じ商上のelementary-test位相・完備性も以下で構成する。
固定sourceの積分範囲の閉性は命題の証明に使用していない。

\begin{proposition}[増大時間区間上の速い部分列]
強発展可測で右連続な過程列 $(X_n)$ が上の一様elementary-test Cauchy条件を満たし、
全ての $n$ で $X_n(0)=X_0(0)$ がほとんど確実に成り立つとする。
このとき狭義増加列 $f$ が存在し、ほとんど全ての経路について、十分大きい全ての $k$ で
\[
 \sup_{t\le k+1}|X_{f(k+1)}(t)-X_{f(k)}(t)|\le 2^{-k-1}
\]
が成り立つ。特に、各有限時間区間では、十分後の隣接差の上限が経路ごとに可算和可能である。
\end{proposition}
\begin{proof}
$J=1_{(0,T]}$ による積分は $t\le T$ で $X(t)-X(0)$ に等しい。
従って共通初期値の仮定により、一様test Cauchy評価は
\[
 \mathbb E\!\left[1\wedge\sup_{t\le T}|X_m(t)-X_n(t)|\right]\le\varepsilon
\]
を十分大きい全ての $m,n$ に対して与える。
$q_k=2^{-k-1}$ とし、$T=k+1$, $\varepsilon=q_k^2$ の閾値を $N_k$ とする。
$f(0)=N_0$, $f(k+1)=\max(f(k)+1,N_{k+1})$ と選ぶ。
隣接差のcapped最大値を $E_k$ とすれば、$\mathbb E E_k\le q_k^2$ である。
Markov評価から $\mathbb P(E_k\ge q_k)\le q_k$ を得る。
$\sum_k q_k<\infty$ なのでBorel--Cantelliにより、共通零集合の外で最終的に $E_k<q_k$ となる。
$q_k<1$ であるからcapを外せる。停止factorial gridは右連続性により
区間全体の最大値を検出するため、結論は全ての $t\le k+1$ に同時に成り立つ。
\end{proof}
この抽出を次の極限構成に使用する。

\begin{proposition}[一様elementary-test Cauchy列の正則極限]
フィルトレーションは通常の条件を満たすとする。
強発展可測でcàdlàgな過程列 $(X_n)$ が一様elementary-test Cauchy条件を満たし、
全ての $n$ で $X_n(0)=X_0(0)$ がほとんど確実に成り立つなら、
強発展可測でcàdlàgな過程 $Y$ が存在し、
\[
 Y(0)=X_0(0)\quad\text{a.s.},\qquad
 \sup_J \bar e_T^J(X_n,Y)\longrightarrow0\quad(T<\infty)
\]
が成り立つ。列全体は $Y$ にucp収束する。
\end{proposition}
\begin{proof}
前の命題の部分列を選ぶ。共通零集合の外で、各有限時間区間上の十分後の隣接差は
可算和可能な数列に支配される。三角不等式と級数のtail評価から局所一様Cauchy性を得る。
実数の完備性により各時刻で極限 $Z(t)$ が存在し、局所一様収束となる。
非収束点を含む全標本点に定義した可測な極限選択を使えば、各 $Z(t)$ は
時刻 $t$ の情報に関して可測であり、$Z$ は強適合である。

各経路を有限時刻で停止すると全時間上一様収束に帰着する。
右連続性と左極限はこの一様極限に保存されるため、$Z$ はほとんど確実にcàdlàgである。
通常の条件により例外集合は時刻零の情報に含まれる。
その集合上で $Z$ を零に置き換え、強適合で全経路がcàdlàgな $Y$ を得る。
右連続な強適合過程は強発展可測である。局所一様収束は共通零集合の外で保存され、
時刻零の極限から $Y(0)=X_0(0)$ を得る。

部分列の局所一様収束は各capped最大誤差の確率収束を与える。
先の一様Cauchy評価の移送証明では、補助添字 $m$ を $f(m)$ に取り替えてもよい。
$f(m)\to\infty$ なので同じCauchy閾値を使用でき、
固定したtestについて $m\to\infty$ とした後、全test共通の誤差評価を得る。
従って元の列全体がelementary-test収束する。共通初期値とunit testからucp収束も従う。
\end{proof}
この構成は近似過程のgood-integrator性や固定sourceの積分範囲の閉性を使用しない。

\begin{proposition}[good-integrator性の極限安定性]\label{prop:gi-limit}
強発展可測で右連続な過程 $X_n,Y$ について、$\sup_J \bar e_T^J(X_n,Y)\to0$ が
各有限 $T$ で成り立ち、各 $X_n$ がgood integratorなら、$Y$ もgood integratorである。
\end{proposition}
\begin{proof}
一様に零へ向かうpredictable elementary integrand列 $H_k$ と時刻 $T$ を固定する。
任意の $\varepsilon,\delta>0$ に対して $a=\min(\varepsilon/2,1)>0$ と置く。
一様test収束から、全unit-bounded test $J$ に対して
$\bar e_T^J(X_m,Y)\le a\delta/4$ を満たす一つの $m$ を選ぶ。
この $m$ を固定したまま、$X_m$ のGI性から、十分大きい $k$ で
\[
 \mathbb P\bigl(|(H_k\cdot X_m)_T|\ge\varepsilon/2\bigr)<\delta/2
\]
を得る。同じtailで $|H_k|\le1$ であるから、積分差のcapped最大値を $E_k$ とすれば
$\mathbb E E_k\le a\delta/4$ であり、Markov評価より
$\mathbb P(E_k\ge a)\le\delta/4$ である。
右連続性により、時刻 $T$ のcapped積分差は $E_k$ 以下である。
三角不等式から
\[
 \{|(H_k\cdot Y)_T|\ge\varepsilon\}
 \subseteq \{|(H_k\cdot X_m)_T|\ge\varepsilon/2\}\cup\{E_k\ge a\}.
\]
従って左辺の確率は $\delta$ 未満となる。$\delta$ は任意なので確率収束を得る。
gainの可測性は $Y$ の強発展可測性とelementary積分の有限和表示による。
\end{proof}

\begin{corollary}[分解可能な領域に属するCauchy極限]
通常の条件と既述のフィルトレーションの$\sigma$有限性の下で、正則零始点のGI過程列が
一様elementary-test Cauchy条件を満たすなら、正則零始点の極限 $Z$ が存在し、
全列が $Z$ にelementary-test収束し、$Z$ は零始点局所マルチンゲールと適合局所有限変動過程の
和に分解できる。
\end{corollary}
\begin{proof}
前の極限生成から $Y$ を得て $Z_t=Y_t-Y_0$ と置く。
中心化はelementary積分を変えず、正則性と強適合性を保つので、全列のtest収束も保存される。
極限安定性から $Z$ はGIであり、零始点càdlàg GIからの分解定理を適用する。
有限変動成分には適合性を要求し、predictabilityは要求しない。
\end{proof}
以上で過程列の極限とその分解を生成した。

\begin{proposition}[同じ区別不能性商上のelementary gauge族]
許容分解を持つ過程の区別不能性商を $\mathcal S$ とする。
各有限 $T$ に対して
\[
 e_T([X],[Y])=\sup_J \bar e_T^J(X,Y)
\]
は代表元によらず定まる。全ての $T,\varepsilon>0$ に対し、十分大きい全ての $m,n$ で
$e_T(x_m,x_n)\le\varepsilon$ となる列 $(x_n)$ には、$y\in\mathcal S$ が存在して
$e_T(x_n,y)\to0$ が全ての $T$ で成り立つ。この極限は一意である。
\end{proposition}
\begin{proof}
区別不能な価格過程は全時刻共通の零集合の外で一致するため、有限和である各elementary積分も
全時刻で一致する。capped grid envelopeとその積分も保存され、testの上限へ降ろせる。

元の代表元自身の適合性・経路正則性は要求しない。
各許容分解 $X=N+A$ の成分和を代表元に選ぶと、強発展可測でcàdlàgな零始点過程が得られる。
一様test Cauchy条件はこの置き換えで保存される。分解からのGI性を使い、前の系により
同じ分解可能領域の極限を得る。区別不能な元の代表列へ戻してから商を取れば存在が従う。

二つの極限があれば、両者と近似列を正則零始点代表元へ置き換える。
固定testの極限の一意性とunit testによって両極限の増分が一致し、共通初期値から
過程全体の区別不能性を得る。これは $\mathcal S$ の等号である。
\end{proof}
この極限生成では、補助距離の完備性・連続性を使用していない。
同じ商を使用するため、補助距離の零距離と区別不能性の同定のみを利用している。
\begin{proposition}[時間重み付きelementary距離]
同じ商 $\mathcal S$ 上で
\[
 d_{\mathcal E}(x,y)=\sum_{n\ge0}2^{-n-1}e_{n+1}(x,y)
\]
は $1$ 以下の距離である。
\end{proposition}
\begin{proof}
代表元のtest誤差も、正則代表元の誤差とのa.e.一致により可積分である。
capped誤差は $[0,1]$ に値を取り、確率測度で積分するから $0\le e_T\le1$ である。
誤差の対称性・三角不等式を積分して全testの上限を取れば、gaugeにもそれらが従う。
時間区間の拡大に関する単調性は正則代表元を選んで全時刻supの単調性から得る。

重みの和は $1$ なので級数は有限であり、$d_{\mathcal E}\le1$ となる。
対称性・三角不等式・対角上の零は各項から従う。
$d_{\mathcal E}(x,y)=0$ なら、各重みが正であるため、全ての $n$ で $e_{n+1}(x,y)=0$ である。
horizon単調性から全ての有限 $T$ でも $e_T(x,y)=0$ となり、
定数列 $x$ のgauge極限の一意性を $x,y$ に適用して $x=y$ を得る。
\end{proof}
\begin{proposition}[elementary距離の独立完備性]
距離空間 $(\mathcal S,d_{\mathcal E})$ は完備である。
\end{proposition}
\begin{proof}
$d_{\mathcal E}$-Cauchy列 $(x_j)$ を取る。有限 $T$ に対して $T\le k+1$ となる整数 $k$ を選ぶと、
\[
 2^{-k-1}e_T(x_m,x_n)\le d_{\mathcal E}(x_m,x_n)
\]
である。重みは正なので、この列は全ての有限horizonでgaugeのCauchy条件を満たす。
既に構成したgauge族の極限 $y\in\mathcal S$ を取る。
固定 $k$ では $e_{k+1}(x_n,y)\to0$ であり、
$0\le2^{-k-1}e_{k+1}(x_n,y)\le2^{-k-1}$ である。
可算和の支配収束から $d_{\mathcal E}(x_n,y)\to0$ を得る。
この証明は補助距離の完備性にも、固定sourceの積分範囲の閉性にも依存しない。
\end{proof}
\begin{proposition}[elementary距離と加法構造]
商 $\mathcal S$ の自然な実線形構造について、$d_{\mathcal E}$ は平行移動不変であり、
差の写像は積の最大距離に関して $2$-Lipschitzである。
従って、この距離の一様構造は加法群構造と両立する。
\end{proposition}
\begin{proof}
有限和積分の価格過程に関する線形性から、同じ過程を二つの引数に加えても
test誤差は変わらない。積分と上限を取り、時間重みで足せば距離の平行移動不変性を得る。
対称性と合わせて $d_{\mathcal E}(-x,-y)=d_{\mathcal E}(x,y)$ も従う。
三角不等式と平行移動不変性により
\[
 d_{\mathcal E}(x-y,x'-y')
 \le d_{\mathcal E}(x,x')+d_{\mathcal E}(y,y')
 \le2\max\{d_{\mathcal E}(x,x'),d_{\mathcal E}(y,y')\}
\]
となる。これは差の一様連続性を与え、加法と負号の連続性も従う。
\end{proof}
\begin{proposition}[固定係数のスカラー倍]
任意の $a\in\mathbb R$ に対して
\[
 d_{\mathcal E}(ax,ay)\le\max\{|a|,1\}\,d_{\mathcal E}(x,y)
\]
であり、固定係数のスカラー倍は連続である。
\end{proposition}
\begin{proof}
$a>0$ の場合、有限和積分の線形性とcapped包絡の評価
$1\wedge(au)\le\max\{a,1\}(1\wedge u)$ を使う。
可積分なtest誤差を積分し、testの上限を取り、時間重みで足せば評価を得る。
$a=0$ は直ちに従い、$a<0$ は負号が等距離写像であることから正の場合に帰着する。
\end{proof}
この上界の係数は $a\to0$ で零へ行かないため、次のGI評価を別途使う。

\begin{lemma}[unit積分の最大値の一様確率有界性]
強適合・右連続なGI過程 $S$ と有限 $T$ に対して、unit-bounded predictable elementary
strategy $H$ 全体の $\sup_{t\le T}|(H\cdot S)_t|$ は一様に確率有界である。
ここではフィルトレーションの右連続性を仮定する。
\end{lemma}
\begin{proof}
GI性から、固定時刻 $T$ のunit積分全体は一様に確率有界である。
尾確率の許容値に対してその上界を $r\ge0$ とし、$R=r+1$ と置く。
$X=H\cdot S$ の絶対値が $R$ を初めて厳密に超える時刻を $\tau$ とする。
右連続性と適合性により $\tau$ は停止時刻であり、有限なら $|X_\tau|\ge R$ である。
$H$ を $\tau\wedge T$ で止めたstrategy $K$ は依然としてunit-boundedである。
$\sup_{t\le T}|X_t|>R$ なら $\tau\le T$ であり、
$|(K\cdot S)_T|\ge R>r$ となる。従って元の最大値の尾確率は
停止strategyの固定時刻の尾確率以下である。
\end{proof}

\begin{proposition}[elementary距離のスカラー同時連続性]
$(\mathcal S,d_{\mathcal E})$ は自然な実線形構造に関して完備実位相ベクトル空間である。
\end{proposition}
\begin{proof}
まず固定 $x$ に対して $a\to0$ の連続性を示す。正則な代表元 $S$ を選ぶと、その分解から
GI性が従う。有限 $T$ と $\varepsilon>0$ を固定し、前補題で最大値の尾確率が
$\varepsilon/2$ 以下となる共通の $R>0$ を選ぶ。
$|a|\le\varepsilon/(4R)$ なら、最大値が $R$ 以下の経路で
\[
 f=1\wedge\sup_{t\le T}|a(H\cdot S)_t|\le\varepsilon/4.
\]
従って可測な事象 $E=\{f\ge\varepsilon/2\}$ の確率は $\varepsilon/2$ 以下であり、
$f\le\varepsilon/2+1_E$ を積分すれば $\mathbb E f\le\varepsilon$ となる。
係数の閾値は $H$ に依存しないので、$e_T(ax,0)\to0$ である。
gauge収束から距離収束への橋により $d_{\mathcal E}(ax,0)\to0$ を得る。

$a_n\to a$, $x_n\to x$ に対して、平行移動不変性と三角不等式から
\[
 d_{\mathcal E}(a_nx_n,ax)
 \le\max\{|a_n|,1\}d_{\mathcal E}(x_n,x)
   +d_{\mathcal E}((a_n-a)x,0)\longrightarrow0.
\]
距離空間なので、この列による連続性が同時連続性を与える。
加法構造と完備性は既に得られている。
\end{proof}

\begin{proposition}[補助距離に関する完備実位相ベクトル空間]
許容分解を持つ過程の区別不能性商は、補助距離 $R$ と自然な加法・実スカラー倍について
完備な距離付け可能実位相ベクトル空間である。
\end{proposition}
\begin{proof}
分解の和、負号、スカラー倍から、許容過程は実ベクトル空間をなす。
加法の連続性は三角不等式から従う。さらに積距離を最大値で定めると
\[
 R\bigl((X-Y)-(X'-Y')\bigr)
 \le R(X-X')+R(Y-Y')
 \le2\max\{R(X-X'),R(Y-Y')\}
\]
であり、差は一様連続である。
スカラー倍の連続性は前補題による。
これらの演算は分離商に降り、商写像は連続線形である。
商における等号は区別不能性と一致し、完備性は既に証明した同じ距離の完備性である。
\end{proof}

\begin{proposition}[二つの位相の間の閉グラフ]
同じ区別不能性商上の恒等写像
$(\mathcal S,d_{\mathcal E})\to(\mathcal S,R)$ のグラフは閉集合である。
\end{proposition}
\begin{proof}
elementary距離で $x_n\to y$、補助距離で $x_n\to z$ とする。
各過程を分解成分の和による正則零始点代表元に置き換える。
商の等号は区別不能性なので、両方の収束は保たれる。
elementary距離の各項の重みは正であり、horizonについてgaugeは単調である。
従って距離収束から任意の有限horizonでの全test共通gauge収束が従う。
unit testと共通初期値から、$x_n-y$ のcapped factorial envelopeは測度収束で零となる。
一方、補助距離の最大値尾確率評価と、同じenvelopeが有限horizonの最大値以下であることから、
$x_n-z$ についても同じ収束を得る。
右連続ucp極限の一意性により $y,z$ は区別不能である。
積空間は距離付け可能なので、列によるこの閉性はグラフの閉性を与える。
\end{proof}

\begin{lemma}[平行移動不変な完備距離を持つ実TVSの開写像と閉グラフ]
$E,G$ を平行移動不変な距離を持つ実位相ベクトル空間とする。
$E$ が完備、$G$ がBaire空間なら、連続な全射線形写像 $f:E\to G$ は開写像である。
両空間が完備なら、閉グラフを持つ線形写像は連続である。
\end{lemma}
\begin{proof}
零近傍 $V\subset E$ は吸収的であり、$G=\bigcup_{n\ge0}(n+1)\overline{f(V)}$ である。
Baireの定理と非零スカラー倍の同相性から $\overline{f(V)}$ の内部は非空である。
任意の零近傍 $U$ に対し対称な $V$ を $V+V\subset U$ と選ぶと、内部の一点を引くことで
$\overline{f(U)}$ が零近傍になる。

$\varepsilon>0$ とし、各 $n$ で
\[
 0<\delta_n\le2^{-n},\qquad
 B_G(0,\delta_n)\subset
 \overline{f\bigl(B_E(0,(\varepsilon/4)2^{-n})\bigr)}
\]
を選ぶ。$d_G(y,0)<\delta_0$ なら、$x_0=0$ から逐次補正して
$d_G(f(x_n),y)<\delta_n$、
$d_E(x_n,x_{n+1})\le(\varepsilon/4)2^{-n}$ とできる。
完備性から $x_n\to x$、幾何級数評価から $d_E(x,0)\le\varepsilon/2$ であり、
連続性と残差の消失から $f(x)=y$ となる。従って $f(B_E(0,\varepsilon))$ は零近傍である。
平行移動により開写像性を得る。
閉グラフは積空間の閉線形部分空間として完備であり、その第一射影に開写像性を適用すれば、
逆写像と第二射影の合成である元の線形写像の連続性が従う。
\end{proof}

\begin{proposition}[elementary収束から補助cost収束と共通局所化へ]\label{prop:topology-localization}
分解可能過程について $X_n\to X$ がelementary-test収束なら $R(X_n-X)\to0$ である。
さらに、強発展可測なc\`adl\`ag GI過程 $S$ に対する実現されたelementary代表列のgainを $G_n$、
指定された利得を $G$ とすると、一つの部分列 $(n_k)$ と一つのexhaustiveな停止列 $(\tau_r)$ を選べる。
各 $\tau_r\le r+1$ で、$G_{n_k}-G$ のprelocal exact representationのcost $c_{r,k}$ は
\[
 \sum_k c_{r,k}\le6(r+4)<\infty,\qquad c_{r,k}\longrightarrow0
\]
を満たす。別の同時選択により、連続する $G_{n_{k+1}}-G_{n_k}$ のcostについても
総和上界 $6(r+5)$ を得る。通常の条件とsigma-finite filtrationを仮定する。
\end{proposition}
\begin{proof}
両距離は平行移動不変であり、恒等写像の閉グラフ性と前補題から
$(\mathcal S,d_{\mathcal E})\to(\mathcal S,R)$ の連続性が従う。
代表列については、各elementary gainのGI性と正則零始点性からJ1分解を生成する。
独立に証明したelementary完備性により正則な分解可能極限を得て、共通ucp極限の一意性で
指定された利得と区別不能であることを示す。従って $R(G_n-G)\to0$ である。
補助costの部分列抽出と共通局所化を適用すると、主張の停止列とcost評価を得る。
連続差については、まず $\sum_k R(G_{n_k}-G)\le1$ となる部分列を選ぶ。
三角不等式から $\sum_k R(G_{n_{k+1}}-G_{n_k})\le2$ なので、
再度の部分列抽出を行わずに総和可能列用の共通局所化を適用できる。
この表示は極限gainの積分としての実現を前提にしない。
\end{proof}

この比較と局所化は固定sourceの積分範囲の閉性に依存しない。
さらに有界sourceについて、次の同時選択を得た。

\begin{proposition}[共通測度上のprelocal costと境界jumpの総和可能性]\label{prop:joint-costs}
前命題の条件に加え、ある定数 $B$ に対して全 $t,\omega$ で
$|S_t(\omega)|\le B$ とする。同値確率測度 $Q$、部分列 $(n_k)$、
exhaustiveな停止列 $\tau_r\le r+1$ を一度に選べる。
$Z_k=G_{n_{k+1}}-G_{n_k}$ のprelocal exact representationのwitness costは
\[
 \sum_k c_{r,k}\le6(r+5),\qquad
 \sum_k\mathbb E_Q|\Delta Z_k(\tau_r)|<\infty
\]
を満たす。同時に非負可測な $\eta\in L^2(Q)$ を保持でき、$Q$ は $\eta$ による
exponential tiltであり、その密度は上に有界である。
各 $k$ について $[0,k+1]$ 上の $|G_{n_k}-G|$ は共通の確率1の集合上で
$\eta$ 以下である。従って $\xi\in L^2(\mu)$ なら $\xi+\eta\in L^2(Q)$ である。
\end{proposition}
\begin{proof}
まずgrowing-horizon誤差包絡から $Q$ を固定する。その後に補助cost収束を適用する。
補助costとgrowing-horizon誤差の $L^1(Q)$ ノルムがともに零へ収束するため、
二つの和を幾何級数以下にする部分列を選ぶ。部分列の添字が増えるので、
誤差を測るhorizonを $k+1$ に縮めても包絡と総和上界1を保持する。
三角不等式から連続差の補助cost総和は2以下であり、共通局所化を適用する。
固定した $r$ について、$k+1\ge r+1$ のtailでは、誤差包絡を $a_k$ と書くと
\[
 |\Delta Z_k(\tau_r)|\le2(a_{k+1}+a_k)
\]
であるから期待値の総和は有限である。残る有限個の項は、有界sourceと各elementary
strategyの有限な係数上界からgain自体を全時間で有界にでき、jumpの期待値も有限となる。
追加の部分列抽出は行わず、同じ停止列で両方の総和可能性が成立する。
\end{proof}

残るのはbounded closed-stopごとのactual成分データと積分同定への接続である。
その比較に必要な二次変分側については、同じwitnessから追加の入力を生成できる。
$\widehat N_k=(N_k-N_k(0))^{\tau_r}$ とすると、一般の逆Davis評価により
\[
 \mathbb E_Q\sqrt{[\widehat N_k]_{r+1}}
 \le7\mathbb E_Q\sup_{t\le r+1}|\widehat N_k(t)|
 \le14c_{r,k}.
\]
最後の不等式は中心化後の最大値が元の停止マルチンゲールの最大値の2倍以下であることによる。
通常の条件の下で内在的二次変分を生成してこの評価を適用するため、
有界ジャンプや別途与えられたquadratic scheduleは不要である。
従って同じ選択でroot期待値の総和も有限となる。
両costの比較は次の構成で得られる。零始点性は確率1で仮定すれば十分である。
元のwitnessを $N,A$ とし、
\[
 M=((N-N_0)^\tau)^T,\qquad
 B=A^{\tau-}-A_0-\Delta N_\tau\mathbf1_{[\tau,\infty)}
\]
とする。$B$ はadaptedな零始点の大域有限変動過程であり、$M+B$ は $t<\tau$ 上で
元の過程に一致する。$\tau=0$ ではこの一致条件は空であり、witnessの初期値に余分な条件は不要である。
strict-prefix FVと境界jumpの変動評価から $\mathbb E\operatorname{Var}(B)\le2c$、
上の逆Davis評価から $M$ の全時間root期待値は $14c$ 以下となる。
従って実際のJ1分解を下限へ代入して
\[
 J_\tau(X)\le16c
\]
を得る。exact representationの区別不能性は、同じ零集合上で初期値とprefix一致に移送できる。
この比較を同時局所化へ適用すると、同じ行の $\sum_k J_{\tau_r}(Z_k)$ も有限となる。
さらに積分表示が与えられれば、補題\ref{lem:canonical-cost}と式\eqref{eq:actual-closed-cost}により
\[
 \mathbb E\sup_{t\le T}|M^{\mathrm{actual}}_t-M^{\mathrm{actual}}_0|
 +\mathbb E\operatorname{Var}_{[0,T]}(A^{\mathrm{actual}})
 \le30\bigl(32c+\mathbb E|\Delta X_\tau|\bigr)
\]
を得る。これは積分表示そのものの存在を主張しない。
補助測度上のspecial certificateを元の測度へ無条件に移送することは用いない。

\subsubsection{同じ選択からactual成分行を生成する}\label{sec:actual-series}
元の有界source $S$ とelementary代表列から、上の $Q,f,\tau_r$ を選ぶ。
価格の有界性は確率1で全時刻に共通の定数上界があれば十分である。
境界jumpの総和評価で初めの有限個の項を支配する際にも、共通零集合の外で上界を使い、
積分のa.e.単調性を適用できる。この上界は $Q\ll\mu$ により $Q$ へ移る。

$Q$ を固定してから元の $S$ のsource分解とcompleted scheduleを構成する。
連続差 $Z_k$ は二つのelementary strategyの差のgainであり、係数絶対値和は二つの行の
上界の和で支配される。このelementary strategyを $\tau_r$ で停止した積分graphを
前節の直接構成で作り、さらに $T_r=r+1$ で再停止する。
$\tau_r\le T_r$ なのでgainは $Z_k^{\tau_r}$ のままであり、係数は
\[
 \mathbf1_{(0,\tau_r]}
   \bigl(H^{f(k+1)}-H^{f(k)}\bigr)
\]
である。二成分は $T_r$ 後にそれぞれ一定で、FV成分は全時間上でBVである。
これらは元の $S$ に対する積分表示を与える。

成分評価は、停止した利得と積分表示された利得との区別不能性を
入力にすれば、独立に構成したこのcertificateへ適用できる。
各成分を $M^{r,k},A^{r,k}$ とすると、同じprelocal costとboundary jumpから
\[
 \sum_k\left(
 \mathbb E_Q\sup_{t\le T_r}|M^{r,k}_t-M^{r,k}_0|
 +\mathbb E_Q\operatorname{Var}_{[0,T_r]}(A^{r,k})\right)<\infty
\]
を得る。測度のtilt、密度上界、growing-horizon誤差包絡 $\eta\in L^2(Q)$ と
元の $L^2(\mu)$ 包絡からの移送も、同じ選択の情報として保持する。

\paragraph{全時間の成分級数}
行 $r$ を固定し、初期値を引いた成分を $m_k=M^{r,k}-M^{r,k}_0$、
$a_k=A^{r,k}-A^{r,k}_0$ とする。
actual行は具体的なelementary representativeを再停止して構成し、M成分の左極限も保持する。
$m_k$ はcàdlàgで $T_r$ 後に一定なので、可測な有限horizon包絡 $g_k$ が全時間で
$|m_k(t)|\le g_k$ を満たす。前段のcost評価から $\sum_k\mathbb E_Qg_k<\infty$ である。
各 $m_k$ はlocal martingaleであり、この可積分包絡によって真のmartingaleとなる。
可積分包絡付きmartingale級数定理により
\[
 m(t)=\sum_k m_k(t)
\]
は真のmartingaleである。部分和は確率1で全時間上一様に $m$ へ収束する。

FV成分については、初期値の差引きは変動を変えず、$T_r$ 後の定数性から
\[
 \operatorname{Var}_{[0,\infty)}(a_k)
 \le\operatorname{Var}_{[0,T_r]}(A^{r,k})
\]
となる。Tonelliと有限な期待値総和から、変動の総和は確率1で有限である。
変動級数定理を適用すると、$a(t)=\sum_k a_k(t)$ は確率1でglobal BVであり、
部分和は全時間上一様に、かつ全時間の変動距離で $a$ へ収束する。
各 $a_k$ は予測可能なので、点ごとの級数として定義した $a$ も予測可能である。
これらは同じ $Q,f,\tau_r$ の選択から実際に生成した成分の結論である。

\paragraph{初項を含むactual近似列と共通包絡}
段階 $r$ では $G_k=H^{f(k)}\cdot S$ と置き、$G_r^{\tau_r}$ の直接構成済みactual
certificateを初項 $P_{r,0}$ とする。同じactual増分行 $R_{r,k}$ を使ってraw層で
\[
 P_{r,n+1}=P_{r,n}+R_{r,r+n}
\]
と定義する。成分はこの加法で足される。係数とgainの有限和がtelescopingするため、
$P_{r,n}$ の係数は $1_{(0,\tau_r]}H^{f(r+n)}$、gainは $G_{r+n}^{\tau_r}$ である。
後者を直接構成したelementary certificateから、係数の等号とgainの区別不能性によって
積分としての実現を移す。一般actual加法calculusは必要ない。初項と増分のM・FV成分はいずれも零始点である。

各elementary gainの定数上界を $c_k$ とする。誤差包絡は
$\sup_{t\le k+1}|G_k(t)-X(t)|\le\eta$ を全ての $k$ について同時に満たす。
$\tau_r\le r+1\le r+n+1$ なので、全時間で
\[
 |G_{r+n}^{\tau_r}(t)|
 \le |G_r^{\tau_r}(t)|
   +|(G_{r+n}-X)^{\tau_r}(t)|+|(G_r-X)^{\tau_r}(t)|
 \le c_r+2\eta=:\zeta_r .
\]
従って $\zeta_r\in L^2(Q)$ は同じactual近似列の共通包絡である。
測度の選び直しも、指定された利得 $X$ の追加の $L^2$ 仮定も要らない。
この有限和近似列と包絡は、元のsourceからの同時局所化の出力へ接続済みである。

\paragraph{正則な成分極限とM\'emin入力}
部分和の右連続性と左極限は全時間一様極限へ移る。
通常の条件を使い、悪い経路を含む時刻0で可測な零集合上で増分級数を修正する。
二成分のversionを取った後、その区別不能性を共通の確率1の集合にまとめる。
Mの級数のmartingale性、FVのpredictability、全時間一様収束と変動距離収束を保持できる。
この正則版に初項を加えた極限を $M_r^\infty,A_r^\infty$ とする。
前者は全経路でcàdlàgなlocal martingale、後者は全経路で右連続かつglobal BVな予測可能過程である。
これらは固定した $r$ の結論であり、停止を外した全過程のglobal BVは主張しない。

隣接するraw有限和の成分差は、同じ増分行の成分に等しい。
FV経路は $T_r$ 後に一定なので、そのcanonical signed Stieltjes measureの全変動は
$[0,T_r]$ 上の経路変動に等しい。これは停止経路の全変動定理を使った等式であり、
raw recordの補助measure欄を使うものではない。
M成分についても任意のhorizon上の差分包絡は、$T_r$ 上の包絡以下である。
従って元の期待costの総和可能性とTonelliにより、両種類の隣接差ノルムは確率1で総和可能となる。

元の指定された利得との同定には、同じ測度選択で既に得た、より強い誤差の情報を保持する。
\[
 \sum_k\mathbb E_Q\sup_{t\le k+1}|G_k(t)-X(t)|\le1.
\]
Tonelliから包絡は経路ごとに総和可能であり、零へ収束する。
同じ確率1の集合上で全時刻 $t$ について $G_k(t)\to X(t)$ となるので、
$t\wedge\tau_r$ でも評価できる。$P_{r,n}$ のgainとの区別不能性と成分収束により
\[
 M_r^\infty+A_r^\infty\ \overset{\mathrm{indist}}=\ X^{\tau_r}
\]
を得る。新しい部分列や測度を選ぶ必要はない。

以上を共通 $L^2(Q)$ 包絡 $\zeta_r$ と合わせ、M\'emin component recordの
全フィールドとstored stopped gainの同定を、元のsourceと代表列から供給済みである。
極限のactual realizationは、以下の直接joint-passage構成を定理へ渡して得る。


\paragraph{停止則からの成分同定と有限加法}
右連続local martingale $M$ の局所化列を $\rho_n$ とし、
$B_n=\{\rho_n>0\}$ と置く。$1_{B_n}M$ を $\rho_n$ で停止した過程はmartingaleである。
これをさらに停止時刻 $\sigma$ で停止し、indicatorと停止の可換性、
$(M^{\rho_n})^{\sigma}=(M^{\sigma})^{\rho_n}$ を使うと、
同じ局所化列が $M^{\sigma}$ のlocal martingale性を与える。
この証明は $M_0=0$ を要求しない。

有限停止はpredictable FV成分のpredictabilityを保ち、local BVをglobal BVへ移す。
したがってraw strategyのgainと二成分を直接停止でき、各初期値も保持される。
積分graphの一座標では、このraw停止とcompleted strategyが同じgainを持つので、
中心化したcanonical成分の一意性からmartingale成分とFV成分を同定できる。
さらに二つのgraphの停止列の最小値で同じ直接構成を行い、energy kernel、
variation bridge、completed integralの停止同定を適用する。
これにより共通schedule上の加法を実供給できる。
以上と局所から全有限horizonへのStieltjes同定を合わせたFV calculusは、
一般のpredictable restriction calculusを仮定せず構成済みである。
同じ有限加法を使うintrinsic終端市場モデルも停止calculusを入力に取らない。
これはspecial graphからの市場構成であり、一般の市場との同定を意味しない。
joint-passage停止側も以下の構成で接続できる。


\subsubsection{Joint-passageの直接構成と停止極限の実現}\label{sec:joint-realization}
固定した外側の停止段階 $r$ で、近似列を $P_{r,k}$、その共通gain包絡を $\zeta_r$ とする。
sourceのcompleted scheduleと全近似列のmartingale passageの共通部分から、
有限horizon以下のexhaustiveな停止列 $\rho_{r,j}$ を作る。
source prefixは前段の直接停止で構成する。
各近似martingaleのpassage前の上界と、Corollary 2.4によるjump包絡を合わせると、
停止時のovershootを含めて
\[
 \|M_{r,k}^{\rho_{r,j}}(T_j)\|_{L^2(Q)}
 \le c_j+6\|\zeta_r\|_{L^2(Q)}
\]
を得る。この停止過程はlocal martingaleであり、積分可能な全時間包絡で支配されるため、
支配収束による定理から真のmartingaleとなる。

固定graph witnessのcompleted coordinateと直接停止したraw成分は同じgainを持つ。
中心化した成分の一意性から、各近似成分の積分表示を同定する。
上で得たFV隣接差の変動総和可能性と、共通座標でのM終端の $L^2$ 有界性を
補題\ref{lem:joint-coefficient}へ渡す。同補題ではFV変動によるcontrolで元の係数列全体の
極限を先に確保し、同じ係数のforward凸化をenergy座標で収束させる。
その結果、一つの予測可能係数が両積分を同時に表す。
energy測度が零となる部分のFV積分も、同補題の変動controlとの一致で同定される。
exhaustiveな全座標を合わせたgraph witnessは、成分和の極限をactual strategyとして実現する。

元のsourceと代表列から既に構成した同じ $Q$ と $\tau_r$ のcomponent dataをこの定理に渡すと、
元の価格過程に対するactual strategy $V_r$ が得られ、
\[
 (V_r\cdot S)\sim_Q X^{\tau_r},\qquad
 \tau_r\le r+1,\qquad \tau_r\longrightarrow\infty
 \quad Q\text{-a.s.}
\]
が成り立つ。ここで左辺は $V_r$ が保持するgainを表す。
新たな停止calculus、測度、近似列は入力に加えていない。
これは $Q$ 上の各停止段階でのspecial certificateであり、
元の測度上の一般graphへの復帰、停止区間の貼り合わせ、終端同定は、
\ref{sec:integral-calculus}節で示した別の段階である。

\subsection{二つの極限成分を表す一つの係数}
\label{sec:joint-coefficient}
直前の実現で用いた係数選択を詳述する。この節の分解は全て、固定した補助確率 $Q$ の下で、
元の価格 $S$ に対して行う。零始点の正準分解を $S-S_0=M^S+A^S$ とし、
実現済み近似積分の係数と成分を
\[
 M_n=H_n\cdot M^S,\qquad A_n=H_n\cdot A^S
\]
と書く。前述の成分級数の構成により、正則な極限 $M,A$、$M_n$ の局所一様収束、
$A_n$ の変動収束、および各有限 $T$ について概確実な
\[
 \sum_n\Var_{[0,T]}(A_{n+1}-A_n)<\infty
\]
が得られている。さらに共通の joint 停止列が得られ、各座標で停止したM成分の終端は、
近似列の添字について一様に $L^2(Q)$ で有界である。

\begin{lemma}[二つの制御測度に対する同時係数選択]\label{lem:joint-coefficient}
この条件の下で、一つの実数値予測可能係数 $H$ が存在し、各 joint 座標上の完成 M 積分と
FV 積分は、それぞれ $M,A$ の対応する停止と一致する。
energy measure と variation measure の間の絶対連続性は仮定しない。
\end{lemma}
\begin{proof}
まず決定論的な $T>0$ を固定し、元価格の経路的変動測度を
$v_T(\omega)=|d(A^S)^T|(\omega)$ と書く。
\[
 v_T(\omega)([0,\infty))+\Var_{[0,T]}A_0+
 \sum_n\Var_{[0,T]}(A_{n+1}-A_n)
\]
を支配する有限可測確率変数 $R_T\geq0$ を選ぶ。既に示した経路的総和可能性から選べる。
構成では、総和可能な M 成分の包絡も $R_T$ に加えてよく、その場合は同じ重みを全成分の評価に使える。
$a_T=(1+R_T)^{-2}$ と置き、零例外集合では $1$ とする。これは概確実に正である。
予測可能 sigma-field 上の有限測度を
\[
 \nu_T(B)=E_Q\left[a_T\int1_B(t,\omega)\,dv_T(\omega)(t)\right]
\]
で定める。変更するのは FV の制御測度だけであり、M 成分と energy の評価は $Q$ の下に保つ。

各近似積分の経路的密度等式より、
\[
 \begin{aligned}
 \int|H_0|\,d\nu_T&=E_Q[a_T\Var_{[0,T]}A_0]<\infty,\\
 \sum_n\int|H_{n+1}-H_n|\,d\nu_T
 &=E_Q\left[a_T\sum_n\Var_{[0,T]}(A_{n+1}-A_n)\right]<\infty.
 \end{aligned}
\]
実際、$d(A_{n+1}-A_n)=(H_{n+1}-H_n)dA^S$ であり、密度 $f$ を持つ符号付き測度の
全変動は、元の全変動測度に密度 $|f|$ を掛けたものである。
Tonelli と $L^1$ の完備性から、係数の列全体は $h_T$ に
$\nu_T$-a.e. かつ $L^1(\nu_T)$ で収束する。
経路ごとにも、最初の項は $v_T(\omega)$ 可積分で、隣接項の $L^1$ 距離の和は有限である。
従って一つの $Q$-零集合の外で、$v_T(\omega)$-a.e. および
経路的な $L^1(v_T(\omega))$ 収束が成り立つ。

可算個のjoint終端座標について、$\sup_n\|M_n(\rho_r)\|_2\leq b_r$ となる
$b_r>0$ を取る。ベクトル $(2^{-(r+1)/2}M_n(\rho_r)/b_r)_r$ は終端 $L^2(Q)$ 空間の
Hilbert直和で有界である。弱部分列コンパクト性とMazurの補題により、
この直和で強収束するforward凸平均を選べる。部分列の重みを元の添字に戻しても台は前進するので、
一つの重み族 $w_{n,k}$ によって全ての終端座標がCauchyとなる。
$\bar H_n=\sum_kw_{n,k}H_k$ と置き、$\bar M_n$ にも同じ重みを使う。
$r$ 番目の座標のpredictable energy measureを $q_r$ とすると、完成積分の等長性より
\[
 \|\bar H_n-\bar H_m\|_{L^2(q_r)}
 =\|\bar M_n(\rho_r)-\bar M_m(\rho_r)\|_{L^2(Q)}
\]
である。ここで $\rho_r$ は決定論的 horizon も含めた有限な座標停止時刻であり、
$\bar M_n$ は同じ重みの M 成分である。
従って係数は可算個の全ての $L^2(q_r)$ で Cauchy である。
一つの狭義増加な $f$ を
\[
 \|\bar H_{f(j+1)}-\bar H_{f(j)}\|_{L^2(q_r)}\leq2^{-j}
 \quad(r\leq j)
\]
となるように選ぶ。$q_r$ は有限なので、Cauchy--Schwarz と Tonelli により
$q_r$-a.e. に点ごとの収束が成り立つ。この列が $\mathbb R$ 内で収束する点ではその極限を、
それ以外では零を $H$ と定める。収束集合は予測可能なので $H$ も予測可能である。
$L^2$ の完備性より、各 $r$ で $\bar H_{f(j)}\to H$ は $L^2(q_r)$ 収束でもある。

この同じ定義は、$q_r$ が零となる部分も含めて FV 側の係数を決める。
$H_n\to h_T$ となる点では、forward 凸平均の性質から
\[
 |\bar H_n-h_T|\leq\sup_{k\geq n}|H_k-h_T|\longrightarrow0
\]
である。従って $H=h_T$ が $\nu_T$-a.e. に成立し、ほとんど全ての経路で
$v_T(\omega)$-a.e. にも成立する。特に $H\in L^1(\nu_T)$ であり、
経路ごとに $v_T$ に関して可積分である。
joint source prefix は決定論的 horizon より前での追加停止なので、その変動測度は $v_T$ の制限である。
同じ可積分性がその座標でも成立する。これで一つの係数の $L^2(q_r)$ と $L^1$ の所属がそろう。

完成 M 積分の等長性と Doob 不等式により、$\bar H_{f(j)}$ の M 積分は、
座標 $r$ 上で $H$ の M 積分へ収束する。一方、forward 凸平均は既存の極限 $M$ を保つ。
確率収束の極限の一意性から、この積分は $M^{\rho_r}$ に一致する。
FV 側では、ほとんど全ての経路で
\[
 \Var_{[0,T]}\bigl(H_n\cdot A^S-H\cdot A^S\bigr)
 =\int|H_n-H|\,dv_T\longrightarrow0
\]
である。同じ積分列は変動距離で $A^T$ に収束するので、零始点性と極限の一意性から
$H\cdot(A^S)^T=A^T$ を得る。
この符号付き測度の等式を $(0,\rho_r]$ に制限すれば、停止時の jump も含めた
joint 座標での FV 同定を得る。可算個の座標と右連続性により、過程等式の零集合を一つにできる。
従って同じ $H$ が両成分を実現する。
\end{proof}
異なる外側の停止段階で選んだ係数同士の一致は要求しない。
\ref{sec:integral-calculus}節では、元測度へ戻した後、互いに素な停止区間上でそれらを貼り合わせる。
補題\ref{lem:joint-coefficient}はその入力となる停止利得の実現を与えるものであり、
係数選択の段階で積分領域の閉性を仮定していない。


\clearpage
\section{本文の七つの解析入力の証明}
\label{sec:input-proofs}
本文との依存境界は第\ref{sec:graph}節の七つの命題と二つのデータ定義である。本付録はそれらの証明を集約する。先行する三付録の内部補題を本文で直接用いる必要はない。
\subsection{正則極限と可測包絡}
\label{sec:path-inputs}
\begin{proof}[命題\ref{input:paths}の証明]
\contractref{P1}について、全時間Cauchy性から厳密増加する $f$ を選び、
\[
 \mu\{\sup_t|Y_{f(n+1)}(t)-Y_{f(n)}(t)|>2^{-n}\}\leq2^{-n}
\]
とする。右連続性により上限は非負有理時刻の上限であり、事象は可測である。
Borel--Cantelliから、ほとんど全ての経路で、十分先の添字から隣接差の全時間上限が $2^{-n}$ 以下となる。
その尾部級数は有限なので一様Cauchy列となり、実数の完備性により全時間一様極限が存在する。
一様極限は右連続性と各正時刻の左極限を保つ。
共通の例外零集合は通常の条件により $\F_0$ に属するので、そこでは全時刻で零と定める。
各固定時刻で適合な実数値過程の極限であり、零集合上の修正も適合性を保つ。
これが強適合càdlàgな代表元を与える。

\contractref{P2}には同じ速い部分列を使う。各 $Y_n$ の経路は、コンパクト区間上のcàdlàg性と
無限時刻での有限極限により全時間有界である。
従って一様Cauchy性が成立する共通の確率1の集合上で
$\sup_{n,t}|Y_{f(n)}(t)|<\infty$ である。
右稠密な可算時刻集合 $D$ による拡張実数値可測関数
$\sup_{n,t\in D}|Y_{f(n)}(t)|$ を取り、有限な場合はその値、無限の場合は零とする。
この有限非負可測関数 $q$ は共通の確率1の集合上で全ての利得を支配する。

\contractref{P3}では $Z=E_\mu e^{-q}$、$dQ/d\mu=e^{-q}/Z$ と置く。
$q$ は有限なので $0<Z\leq1$、密度は正であり $Q\sim\mu$ となる。
$x^2e^{-x}$ は $[0,\infty)$ で有界なので $q\in L^2(Q)$ である。
同値測度は同じ零集合を持つため完備性を保ち、filtrationの右連続性は測度に依存しない。
従って通常の条件も保たれる。
\end{proof}

\subsection{積分構成で使う内部の定理}
\label{sec:integral-background}
\subsubsection{分解と積分の定義}
\begin{theorem}[包絡を持つ過程の分解]\label{thm:decomposition-input}
通常の条件の下で、適合càdlàg good integrator $Y$ が
$\sup_t|Y_t|\leq q$、$q\geq0$、$q\in L^2(Q)$ を満たすとする。
このとき $Y-Y_0=M+A$ と分解できる。$M$ は零始点càdlàg local martingale、
$A$ は零始点の予測可能càdlàg局所有限変動過程であり、両成分は区別不能性を除いて一意である。
$M$ には真のL²マルチンゲールとなる有限停止座標のexhaustiveな列がある。
有界価格 $S$ には $q$ を定数として適用する。同値測度の下でgood-integrator性は保たれるが、
分解はその測度の下で作り直す。
\end{theorem}
\begin{proof}[証明の所在]
付録の\ref{sec:interface-decomposition}節で、分解と補償の結果から導く。
市場のNFLVRや積分領域の閉性はこの入力の仮定ではない。
\end{proof}

予測可能elementary test $J$ は、有限個の停止区間と左端で可測な有界係数からなる。
任意のcàdlàg過程に対して $J\cdot Y$ は増分の有限和で定義できる。
\[
 d_T^Q(Y,Z)=\sup_{J:\,|J|\leq1}
 E_Q\left[1\wedge\sup_{t\leq T}|(J\cdot(Y-Z))_t|\right].
\]
各有限 $T$ で $d_T^Q(Y_n,Y)\to0$ となることをÉmery収束という\cite{Emery79}。
上限の可測性には可算な右稠密時刻集合を用いる。
Émery Cauchy性の量化順序は $\forall T,\varepsilon>0\ \exists N\ \forall n,k\geq N$ であり、
この $N$ は全ての単位有界testに共通である。

特殊分解を持つ過程の\emph{局所special積分}は、停止した座標上で
マルチンゲール成分のenergy測度に関する $L^2$ 積分と、FV成分の変動測度に関する
$L^1$ 積分を同じ予測可能係数に適用し、停止区間で貼り合わせたものをいう。
FV積分は経路ごとのStieltjes積分である。必要な局所化・完成・一意性は次の定理の一部である。

\begin{definition}[一般積分graph]\label{def:general-graph}
$H^{(n)}=H1_{\{|H|\leq n+1\}}$ とする。予測可能係数 $H$ と零始点・強適合・càdlàg過程 $X$ に対し、
$H^{(n)}$ の局所special積分が $X$ にÉmery収束するとき、$(H,X)$ を元価格 $S$ の一般積分graphに入れ、
$X=H\cdot S$ と書く。過程は区別不能性まで同一視する。
一般の $X$ 自体にspecial分解や無限時刻の終端を要求する定義ではない。
\end{definition}

\begin{theorem}[積分の基本則と評価]\label{thm:integral-input}
有界価格 $S$ と上の一般積分graphについて、次が成り立つ。
\begin{enumerate}
\item 有界予測可能係数の局所special積分が存在する。同じ係数の利得は区別不能性まで一意である。
局所special積分は、係数と利得を保って一般積分graphに属する。
各一般積分利得は元価格のelementary積分列でÉmery近似でき、good integratorである。
\item 局所special積分では線形性、閉停止、予測可能制限が成り立つ。
$Y=M+A$ に対する有界予測可能係数 $k$ の積分は
$k\cdot Y=k\cdot M+k\cdot A$ と分解され、
$\Delta(k\cdot Y)_t=k_t\Delta Y_t$ である。
同じ有限和の重みで両成分を取れる。有限時刻 $T$ での停止は局所FV成分を全時間BVにする。
\item 同値確率測度間で、Émery収束と一般積分graphは係数・利得を保って移送できる。
従って同じ許容下界と、存在する場合の同じ固有終端を保つ。
\item 固定した有限horizon上で、単位有界予測可能係数を、単位有界elementary係数で
M積分とFV積分の両方について同時に近似できる。M側は平方可積分な停止座標での終端 $L^2$、FV側は変動での収束である。
零始点 $L^2$ マルチンゲール $N$ と $|J|\leq1$ に対して
\[
 E_Q[1\wedge\sup_{t\leq T}|(J\cdot N)_t|]\leq2\|N_T\|_2.
\]
そのようなL²停止座標がhorizonを尽くす局所マルチンゲール $M$ について、
全単位testの同じhorizon上の期待値上界が $b$ なら、
$|k|\leq1$ に対する $k\cdot M$ の全単位testにも上界 $b$ が成り立つ。
\item 予測可能FV過程 $A$ には予測可能な正のHahn集合 $B$ があり、
$C=1_B\cdot A$ は非減少で、$C_c-C_a\geq A_c-A_a$ である。
また $|J|\leq1$ なら $\sup_{t\leq T}|(J\cdot A)_t|\leq\Var_{[0,T]}A$ である。
これらの区間は右端を含み、閉停止時のjumpも保持する。
\end{enumerate}
\end{theorem}
\begin{proof}[証明の所在]
付録の\ref{sec:integral-calculus}節で各項を示す。
一般利得の極限実現を仮定せず、有界係数の完成積分、切断、一意性から導く。
\end{proof}

\subsubsection{極限実現と元市場での操作}
次の定理は\cite[Theorem 4.2の証明]{DS}で用いるM\'eminの閉性定理\cite{Memin80}に対応する。
\begin{theorem}[一般積分graphの閉性]\label{thm:graph-closure}
同じ有界価格 $S$ の一般積分利得 $Y_n$ が、零始点・強適合・càdlàg過程 $X$ に
Émery収束するなら、予測可能係数 $H$ が存在し、$(H,X)$ も同じ一般積分graphに属する。
\end{theorem}
\begin{proof}[証明の所在]
付録の\ref{sec:closure-proof}節で、位相比較、共通局所化、同一係数による両成分の実現を接続して証明する。
以後はこの結論だけを用い、そこで選ぶ補助測度や停止列は用いない。
\end{proof}

\begin{corollary}[元市場での取引操作]\label{cor:market-operations}
一般積分利得の有限線形結合、閉停止、停止区間への制限、有限回の予測可能な貼り合わせは
再び元価格の一般積分利得である。
さらに、元価格に実現された零始点利得 $Y=M+A$ が $L^2(Q)$ 包絡を持つなら、
その局所special積分として実現された利得 $V=k\cdot Y$ も、元価格の一般積分利得として実現できる。
成分は同じ $V=k\cdot M+k\cdot A$ のまま保持する。
$Q$ が元測度と同値なら、元測度にも同じ利得を戻せる。
$V\geq-a$ で $T$ 以後一定なら、その固有終端は $V_T$ であり、同じ許容水準の終端集合に入る。
\end{corollary}
\begin{proof}
定理\ref{thm:integral-input}のelementary近似に、有限和・停止・制限を施す。
有限和の結合則と三角不等式により、その利得は指定した操作後の過程へÉmery収束する。
貼り合わせの係数は、停止時刻で既知の事象とその後の停止区間の指示関数の積であり、予測可能である。
定理\ref{thm:graph-closure}により、これらの正則な零始点極限を実現する。

$k$ が有界の場合、同じelementary係数で $M,A$ の積分を近似する。
各elementary変換は $Y$ の有限個の停止増分の和なので、さらに $Y$ の元価格上の近似を代入できる。
各有限和での近似誤差を先に小さくし、整数horizonについて対角選択すると、
元価格のelementary積分列が同じ $V$ へÉmery収束する。
再び閉性定理を使う。非有界な実現済み $k$ は、定理\ref{thm:integral-input}の切断収束で
有界な場合へ帰着し、同じ閉性定理を適用する。
測度移送は一般graphに適用する。下界・停止後の定数性・固有終端は同じ利得の性質なので保たれる。
\end{proof}


\subsection{包絡を持つ過程の分解の導出}
\label{sec:interface-decomposition}
\begin{proof}[定理\ref{thm:decomposition-input}の証明]
まず $Y-Y_0$ に定理\ref{thm:regular-gi-decomposition}を適用し、零始点の
$Y-Y_0=N+D$ を得る。ここで $N$ は局所マルチンゲール、$D$ は適合的な局所FV過程である。
この定理は、有界sourceの分解、大ジャンプの除去、停止後の貼り合わせから得ており、積分値域の閉性を使わない。
$|Y-Y_0|\leq2q$ であり、$q\in L^2(Q)\subseteq L^1(Q)$ である。

有限horizonを固定し、$N$ が真のマルチンゲールとなる局所化を先に施す。
さらに $|N|$ と $\Var(D)$ の水準 $r$ への同時passageで停止する。
閉停止を $\sigma_r$ とすると、停止直前の上界と
$\Delta D=\Delta Y-\Delta N$ から
\[
 \Var(D^{\sigma_r})\leq2r+2q+|N_{\sigma_r}|
\]
となる。右辺は任意抽出と $q\in L^1$ により可積分である。
付録の\ref{sec:canonical-cost}節の補償定理を、この有限horizon後一定な過程に適用する。
\[
 M^{(r)}=N^{\sigma_r}+D^{\sigma_r}-(D^{\sigma_r})^p,
 \qquad A^{(r)}=(D^{\sigma_r})^p
\]
はspecial分解を与える。予測可能FV局所マルチンゲールは零始点なら零なので、
成分は重なる停止区間で一致する。経路の有界性と局所有限変動性から停止列は各horizonを尽くす。
局所化を除き、整数horizonで貼り合わせると所要の分解と一意性を得る。
補題\ref{lem:explicit-jump-envelope}を、この存在が得られた分解と包絡 $2q$ に適用する。
固定した $T$ で $M$ のjumpは一つのL²確率変数 $\xi_T$ で抑えられる。
$|M|$ の水準 $n$ へのpassageと $T$ で停止すると、閉停止した最大値は $n+\xi_T$ 以下である。
従って支配された局所マルチンゲールは真のL²マルチンゲールである。
$T=n$ として水準も増加させれば、経路の局所有界性から停止列は全時間を尽くす。
このjump補題は、既に与えられたspecial分解への条件付き期待値とDoob評価から証明され、
本定理による分解の存在を逆に仮定するものではない。
最後に、good-integrator性は有限時刻の確率収束の条件であるため、同値測度間で保存される。
ここで移送するのはこの条件であり、得たspecial成分ではない。
\end{proof}

\subsection{積分則の構成と移送}
\label{sec:integral-calculus}
ここでは定理\ref{thm:integral-input}を証明する。局所special積分と一般graphは定義\ref{def:general-graph}のものとする。

\begin{lemma}[分解を仮定しない elementary-test 収束の測度不変性]
過程列 $X_n$ と $Y$ は強発展可測とする。全ての有限 $T$ と $\epsilon>0$ に対し、
ある $N$ が存在して、$n\ge N$ と全てのunit-bounded予測可能elementary test $J$ について
\[
 E_\mu\left[1\wedge\sup_{t\le T}
   |(J\mathbin{\cdot}X_n)_t-(J\mathbin{\cdot}Y)_t|\right]\le\epsilon
\]
となる性質は、同値確率測度間で不変である。
経路が右連続でない場合には、上式のsupを可算horizon skeleton上で取る。
\end{lemma}
\begin{proof}
各testのcapped envelopeは可測で $[0,1]$ に値を取る。
移送先で一様零収束が失敗すれば、ある $\epsilon>0$ と $n_k\ge k$、test列 $J_k$ を取り、
移送先の積分を常に $\epsilon$ より大きくできる。
元の測度ではtailがtestに一様なので、この選択後のenvelope列の積分も零へ収束する。
有界な非負関数列の積分零収束を確率収束へ移し、絶対連続性で移送し、
再び有界性から積分零収束を得ると矛盾する。逆方向も同じである。
この議論にspecial分解は不要である。
\end{proof}
elementary代表列の一様誤差評価は、有限和の結合則によりこの過程の収束を与える。
ただしtestは増分しか測らないので、零始点性は別途必要である。
この位相的事実だけで極限integrandの存在や一般integral graph membershipを結論しない。

\begin{lemma}[初期値を固定した極限の一意性]
強発展可測な過程列が上の意味で右連続な二過程 $Y,Z$ に収束し、
$Y_0=Z_0$ が概至る所成り立つなら、$Y,Z$ は区別不能である。
\end{lemma}
\begin{proof}
一つのelementary testを固定すると、capped envelopeの積分零収束は確率収束を与える。
コンパクト一様極限の一意性より、testを積分した二過程は区別不能となる。
固定時刻 $T$ に対し $(0,T]$ のunit testを選び、時刻 $T$ で評価すると、
$Y_T-Y_0=Z_T-Z_0$ を得る。初期値の一致と、可算な右稠密集合上の一致を右連続性で拡張して結論する。
\end{proof}

\begin{proposition}[切断graphの一意性と移送の還元]
固定sourceに対し、同じ $H$ に対応する切断graphのgainは区別不能である。
また、同値測度間で有界係数のspecial graphを移送できれば、切断graphも移送できる。
\end{proposition}
\begin{proof}
二つのcertificate列について、各行の係数は同じ $H^{(n)}$ なので、固定sourceの
graph一意性から行ごとのgainが区別不能となる。これで一つの共通近似列に揃え、
前補題を零始点の正則極限へ適用する。
測度移送では $|H^{(n)}|\le n+1$ により有界係数の移送を各行に適用する。
選び直した実現された利得は元の行のgainと区別不能であり、過程の収束の測度不変性と
区別不能な行への置換により、同じ正則極限と指定された利得を得る。
極限のspecial性を移送する操作はない。
\end{proof}

$|H|\le1$ の積分表示なら全ての切断が $H$ と等しいため、
零始点性と左極限を確認すれば定数近似列からこのgraphに入る。

\begin{proposition}[非有界係数の切断列と座標ごとの誤差制御]
大域積分graphの係数を $H$ とする。各 $n$ について
$H1_{\{|H|\le n+1\}}$ の局所積分表示を同じsourceに対して構成できる。
元のcompleted scheduleの座標 $k$ の係数を $h_k$、energy measureを $q_k$、
variation measureを $v_k$ とし、$r_{k,n}=h_k1_{\{|h_k|>n+1\}}$ とおくと、
\[
 \|r_{k,n}\|_{L^1(v_k)}\longrightarrow0,\qquad
 \|r_{k,n}\|_{L^2(q_k)}\longrightarrow0.
\]
同じ相補切断のcompletedマルチンゲール積分の終端 $L^2$ normも零へ収束する。
\end{proposition}
\begin{proof}
各座標の係数を同じ予測可能集合に制限し、整合したcompleted座標をgluingする。
座標の係数はもともと両control空間に属するので、可積分関数の大きな値のtail評価により
上記の二つのnormが消失する。control measureの全質量が有限であるという追加仮定は不要である。
completedマルチンゲール積分の終端normはhorizonで制限した係数のenergy normに等しく、
元の係数tailのenergy norm以下である。
\end{proof}
\begin{proposition}[非有界係数の一様切断収束と包含]
固定sourceのマルチンゲール成分が左極限を持つとする。
非有界係数を含む全ての大域・局所special積分表示は、
元の係数と指定された利得を保って、全時間の切断graphに入る。
\end{proposition}
\begin{proof}
まず一つのcompleted座標を固定する。零始点の平方可積分マルチンゲール $M$ と、
表されたintegrandが $|J|\le1$ を満たす任意の予測可能elementary testに対して、
離散積分の $L^2$ 縮小評価、chronological grid極限、Doob評価から
\[
 \mathbb E\!\left[1\wedge\sup_{t\le T}|(J\mathbin{\cdot}M)_t|\right]
 \le 2\|M_T\|_2
\]
を得る。testの項数や係数絶対値和への上界は不要である。
マルチンゲール積分はcompleted horizon後で定数なので、同じ評価を任意のtest horizonに使える。

有限変動部分では、variation可積分な係数 $r$ のcompleted積分 $B$ を
経路ごとのStieltjes積分と区別不能に同定する。符号付きStieltjes測度の密度表示より、
\[
 |dB|(\mathbb R_+)\le \int |r_t|\,d|A|_t,\qquad
 \sup_t |(J\mathbin{\cdot}B)_t|\le |dB|(\mathbb R_+).
\]
右辺の $L^1$ normはcanonical variation normで抑えられる。
reference measureから元の確率への確率収束の移送と上界1の支配により、
切断した右辺の期待値も零へ収束する。これは全testに共通の支配である。

両成分の相補係数tailにこれらを適用し、加法性でbounded cutとtailの和を元の積分と同定する。
従って各座標でbounded cutの利得が元の利得へ一様test収束する。
停止列の座標を先に選んで $\mathbb P(\tau_k\le T)$ を小さくし、次にその座標の行閾値を選ぶ。
停止前では全testの利得が一致し、停止後のcapped誤差は1以下であるため、局所化を外せる。
これで大域graphの包含を得る。

局所graphでは、各正の整数時刻での停止が大域積分graphを持つ。
一つのsource schedule上で有界切断係数を積分し、固定sourceの一意性により、
その各停止を停止graphのbounded cutと同定する。前段の停止ごとの収束から再び局所化を外す。
停止ごとの正則代表元を可算共通零集合の外で元の局所gainへ移し、通常の条件で零集合を補正する。
これにより全経路で正則な代表元と零始点性も得られる。
非有界な局所係数自体の共通schedule表示は用いない。
\end{proof}

\begin{proposition}[有界な大域certificateの包含]
固定sourceのマルチンゲール成分が左極限を持つとする。
任意の有限定数 $b$ に対し、$|H|\le b$ を満たす大域積分表示のgraphは切断graphに入る。
また、有界係数を持つ局所graphを任意の正の有限時刻 $T$ で停止したgraphも切断graphに入る。
\end{proposition}
\begin{proof}
大域certificateの各completed座標の係数を予測可能集合 $\{|H|\le n+1\}$ に制限する。
座標全体で同じ切断係数を持つため、整合性とgluingから各切断の局所積分表示を得る。
その全ての有限時刻停止の積分としての実現も、同じschedule上のcompleted停止則から従う。
$n+1\ge b$ となるtailでは係数が元の $H$ と等しいため、固定sourceの一意性から利得も区別不能である。
元のwitnessを自己refinementして、sourceの左極限を持つ座標を選ぶ。
completedマルチンゲール座標のcàdlàg gluingと局所有限変動座標のgluingを加えると、
固定sourceの一意性により元の利得と区別不能な正則代表元を得る。
その零始点性はcompleted graphから得られる。
tail上ではこの代表元と各切断利得が区別不能なので、全test共通の誤差積分は零となり、必要な収束を得る。
保存された利得とは区別不能性で結ぶ。

局所graphは時刻 $T$ で停止すると大域積分表示を持つ。
停止係数 $1_{(0,T]}H$ の上界は $\max(b,0)$ であり、利得は $X_{t\wedge T}$ と区別不能である。
前半をこの停止certificateへ適用する。停止前の利得に大域有限変動の分解を要求する必要はない。
\end{proof}

\begin{proposition}[有界予測可能係数の生成と局所graph全体の包含]
元の有界sourceと任意の有界な予測可能係数 $H$ から、同じsourceに対する
局所積分graphと切断graphを生成できる。さらに、有界係数の局所graphは
保存された利得を保ったまま、全時間の切断graphに入る。
\end{proposition}
\begin{proof}
元のsourceから直接構成したcompleted scheduleを使う。各座標のcanonical variation
measureとquadratic energy measureは有限である。従って $|H|\le b$ から、
前者に関する $L^1$ 可積分性と後者に関する $L^2$ 可積分性を得る。
全座標の係数を同じ $H$ としてgluingへ渡せば、局所積分graphが得られる。
各座標のマルチンゲールが左極限を持つようscheduleを自己refinementし、
càdlàgマルチンゲールのgluingと局所有限変動成分を加えて正則代表元を得る。
零始点性はcompleted graphの時刻0での同定から従う。

各 $H^{(n)}$ にも同じ構成を適用する。$n+1\ge b$ のtailでは $H^{(n)}=H$ なので、
固定sourceのgraph一意性により全利得が区別不能である。従って全elementary testに
共通のtailで誤差は零となり、切断graphの要求する収束が従う。
局所graphの利得とは、同じ $H$ のgraph一意性をもう一度使って同定する。

\end{proof}

\begin{proposition}[同値確率測度間の積分graphの移送]
同じ有界価格過程 $S$ のsourceを同値確率測度 $\mu,\nu$ の下に持つとする。
有界予測可能係数の局所積分graphは、係数と利得を保って移送できる。
従って一般切断graphも移送でき、Mémínの停止付き実現は元の $\mu$ 上のgraphへ戻る。
\end{proposition}
\begin{proof}
まず両測度に有界係数の大域certificateがあるとする。
二つのschedule座標のenergy controlの和とvariation controlの和は有限測度であり、
時刻0のsliceに質量を持たない。有界予測可能係数はそれぞれの $L^2,L^1$ に属する。
共通のelementary密度により、両controlで同時に収束する同じ列 $J_n$ を取る。
completed積分の差は、係数のenergy $L^2$ とvariation $L^1$ の消失から
全test一様に零へ収束する。各座標のelementary積分は、元の $S$ に対する
$J_n\mathbin{\cdot}S$ をその座標の停止時刻で止めたものと一致する。

二つの座標の停止時刻の最小値でさらに停止する。
有限値の停止はtest積分のcapped envelopeを増加させず、収束を保つ。
従って両測度の積分は同じ停止付きelementary列の極限となる。
収束の同値測度不変性と零始点の極限一意性により、停止した利得は区別不能である。
最小値の停止列もexhaustiveなので停止を除ける。
各決定論的停止にこの議論を適用すれば、局所certificateの利得も同定できる。

移送先では有界予測可能係数のgraphをsourceから直接構成する。
先の一意性で指定された利得を同定すれば有界graphの移送が得られる。
一般graphは各有界切断へこの移送を適用し、正則な極限への収束を移送する。
停止付きMémín実現にもこれを適用し、停止列のexhaustivenessは絶対連続性で移す。
非有界利得のspecial分解を移す操作は含まない。
\end{proof}

\begin{proposition}[停止区間の貼り合わせによる全時間の実現]
元の有界sourceの零始点で正則な選択利得 $X$ が、有限値のexhaustiveな停止列の各段階で
一般切断graphを持つなら、一つの予測可能係数 $H$ の全時間graphとして実現できる。
\end{proposition}
\begin{proof}
まず、有界係数のgraphをsourceの直接構成した共通scheduleに載せる。
各completed座標の停止則により、任意の停止時刻で止めた利得の積分graphを生成できる。
有限停止はelementary-test収束を保つので、一般切断graphにも停止則が渡る。
さらに $\sigma\le\tau$ なら、二つの停止を同じ切断行に適用し、その差を取ることで
$(\sigma,\tau]$ への係数制限と利得 $X^\tau-X^\sigma$ を得る。
この係数制限と大きさによる切断は可換である。

停止列の単調性と発散が失敗する共通零集合は時刻0で可測である。
その上で停止時刻を決定論的な $n+1$ に置き換え、全経路で単調な列 $\upsilon_n$ を得る。
元の停止graphの利得はこの修正後も区別不能である。
最初の $(0,\upsilon_0]$ と後続の $(\upsilon_n,\upsilon_{n+1}]$ 上に
各段階の係数を制限する。これらの台は交わらないので、係数の大きさによる切断は
有限和と可換となり、積分近似列の加法性と収束の加法性で有限和のgraphを構成できる。
利得はtelescopingにより $X^{\upsilon_n}$ と一致する。

係数の有限和は $(0,\upsilon_n]$ 上でそれ以降変化しない。
点ごとの極限 $H$ は予測可能であり、この区間への制限は対応する有限和と全点で等しい。
$H$ の各有界切断の実現された積分を元のsourceから直接生成する。
その停止利得を有限和graphの切断行と一意性で同定し、局所化誤差を消して全時間収束を得る。
これが $(H,X)$ の一般切断graphである。係数のoverlap一致は使っていない。

\end{proof}




\paragraph{単位有界変換とHahn成分}
定理\ref{thm:integral-input}の第4・5項の証明を与える。以下の局所化座標は有限horizonで真のL²成分を持つ。

\begin{proposition}[単位有界completed変換と有限horizonのHahn利得]
通常の条件を満たす確率空間上で、$Y=M+A$ とし、$Y,M,A$ は強適合・右連続であり、
$A$ は予測可能な全時間有界変動過程とする。
$M$ は真のマルチンゲールで、固定した $T$ に対し $M_T\in L^2$ とする。
任意の単位有界elementary test $J$ に対して
\[
 \mathbb E\left[1\wedge\sup_{t\leq T}|(J\cdot M)_t|\right]\leq b
\]
が成り立つとする。単位有界な予測可能係数 $k$ の有限horizonのcompleted積分を
$I^k=(k\mathbf 1_{(0,T]})\cdot M$ とすると、同じ $b$ で全elementary testに対する
$I^k$ の評価が成り立つ。

さらに、$A$ の正のHahn集合 $B$ と有限horizonの積分利得
$R=\mathbf 1_{B\cap(0,T]}\cdot Y$ を構成できる。
同じ $R=N+C$ について、ほとんど確実に全ての $0\leq a\leq c\leq T$ で
\[
 0\leq C_c-C_a,\qquad A_c-A_a\leq C_c-C_a
\]
となり、$N$ の全elementary testの上界と
$\mathbb E[1\wedge\sup_{t\leq T}|N_t|]\leq b$ が成り立つ。
\end{proposition}
\begin{proof}
正規化したvariation measureとmartingale energy measureはいずれも有限である。
単位有界な予測可能集合環密度により、$|J_n|\leq1$ を保ち、
同じ係数 $k\mathbf 1_{(0,T]}$ へL² energyとL¹ variationの両方で収束する
一つのelementary列を取る。
energy連続性とelementary積分との一致から、$J_n\cdot M$ は $I^k$ へ
elementary-Émery収束する。gain側も同じ列でcompleted積分に収束するが、
M側の収束の評価にvariation誤差は不要である。

別の単位test $J$ に対し、有限horizon内では
$J\cdot(J_n\cdot M)=(JJ_n)\cdot M$ であり、$|JJ_n|\leq1$ である。
従って近似積分のtest上界は全て $b$ 以下となる。
極限とのtest誤差に三角不等式を適用し、任意の正の誤差を零へ送ると、
completed積分も同じ上界を満たす。$b$ は係数やtestを選ぶ前の上界のままである。

正規化したFV bridgeからHahn集合を生成し、horizonの単位実現された積分をこの集合で制限する。
completed FV成分とStieltjes積分の一致により、$C_c-C_a$ は
$B$ の時間切片と $(a,c]$ の交わりのsigned massとなる。
Hahn分解の正部分の性質が二つの増分評価を与える。区間は右端を含む。
Mのtest上界は既に示したcompleted変換の上界から従う。
completed Mの零始点性を使ってhorizonの単位testを選べば、capped最大関数の上界も得られる。
\end{proof}

局所版では、次のように停止を除く。

積分graph自身の局所化列を $\sigma_r$、その有限completion horizonを $U_r$ とする。
正規化した積分表示されたM成分の停止は、被積分過程 $k$ のcompleted martingale積分と区別不能である。
係数を $(0,U_r]$ に切ってもenergy measure上の同値類は変わらず、単位有界性も保たれる。
従って前命題の単位近似を各座標へ適用できる。

ここでtest horizonは $T$ に固定したままとする。
elementary積分を $U_r$ で停止してから $J$ でtestしても、そのcapped最大関数は
停止しない $(JJ_n)\cdot M^{\sigma_r}$ のもの以下である。
また $M^{\sigma_r}$ の各testは、右連続性と停止したelementary gainの恒等式により、
同じ $T$ 上の $M$ のtest以下となる。$U_r$ と $T$ の大小関係は不要である。
これで全座標の上界を同じ $b$ にできる。

$\{\sigma_r>T\}$ では停止したtestと未停止のtestが一致する。
補集合での誤差は1以下なので、未停止testの期待値は
$b+\mathbb P(\sigma_r\leq T)$ 以下である。$r\to\infty$ として上界 $b$ を得る。
中心化はelementary gainを変えない。局所FVの積分graphは $T+1$ で停止してから適用し、
$[0,T]$ での一致を使う。



一般graphのelementary密度は、各有界切断行の密度と整数horizonでの対角選択から得る。
その利得は命題\ref{prop:gi-limit}によりgood integratorである。

\subsection{一般積分graphの閉性の証明}
\label{sec:closure-proof}
\begin{proof}[定理\ref{thm:graph-closure}の証明]
\emph{一つの元価格上の近似列を選ぶ。}
元価格の一様上界が失敗する共通零集合は通常条件により時刻0で可測である。
その上で価格を零に置き換えれば、利得の区別不能性を保ったまま全点での有界性を満たせる。
定理\ref{thm:integral-input}の有界切断とelementary密度を各 $Y_n$ に適用する。
各整数horizonの全test一様誤差を幾何的に小さくする対角選択により、
同じ $X$ にÉmery収束する元価格のelementary積分列 $P_n$ を得る。
$X$ は仮定により零始点・強適合・càdlàgである。
命題\ref{prop:gi-limit}が $X$ のgood-integrator性を与え、
定理\ref{thm:regular-gi-decomposition}が補助位相の領域に属する分解を与える。
ここでは $X$ の積分表示を仮定しない。

\emph{同じ列に共通する局所化を選ぶ。}
命題\ref{prop:topology-localization}を $P_n-X$ に適用し、補助 $r^1$ costが零へ収束することを得る。
同命題の位相比較は、分解可能過程の空間の完備性と閉グラフ定理に基づき、積分値域閉性から独立である。
さらに速い部分列を取ると、隣接差のcostが和可能となる。
命題\ref{prop:joint-costs}は、この部分列と一つの同値測度 $Q$、一つの停止列 $\tau_r$ を選び、
prelocal成分costと境界jumpの期待値を同時に和可能にする。
$\tau_r\leq r+1$、$\tau_r\to\infty$ を保持する。

\emph{停止した実際の積分成分へ評価を渡す。}
各 $P_n^{\tau_r}$ は元価格の実現済み積分である。
補題\ref{lem:canonical-cost}と式\eqref{eq:actual-closed-cost}により、
任意の分解から得たcostの上界を、この同じ利得の正準成分の上界へ移す。
\ref{sec:actual-series}項で、この評価を同じ $r$ の全隣接差に適用して成分級数を構成する。
その結果、M成分の局所一様極限とFV成分の変動極限、その和が $X^{\tau_r}$ であることを得る。
新たな列や独立の分解を選んで同一視する操作はない。

\emph{一つの係数で両成分を実現する。}
\ref{sec:joint-realization}項のjoint停止は、近似M成分の終端を全添字で一様に $L^2(Q)$ 有界にする。
FV隣接差の変動総和可能性とともに補題\ref{lem:joint-coefficient}の仮定を満たす。
同補題は、FVの変動controlでの係数極限と、同じ係数の凸化によるenergy極限を一致させる。
従って一つの予測可能係数が停止した両成分を表し、その和は同じ $X^{\tau_r}$ である。
energy測度が零となる部分にもFV側から係数が定まり、両測度間の絶対連続性は必要ない。

\emph{元測度と全時間へ戻す。}
定理\ref{thm:integral-input}の切断包含と測度移送により、各停止利得の一般graphを元測度へ戻す。
付録の\ref{sec:integral-calculus}節で証明した停止区間の貼り合わせをこの停止列に適用する。
互いに素な区間上で係数を継ぎ、停止利得の一致と $\tau_r\to\infty$ により
元価格の全時間の一般graph $(H,X)$ を得る。
貼り合わせが必要とするのは利得の区別不能性であり、別々に選んだ係数の重なりでの一致ではない。
補助測度上のspecial分解自体は移送しない。
\end{proof}

\paragraph{利得領域と終端集合}
elementary-test実現とは、各行の係数絶対値和が決定論的定数で抑えられた
予測可能初等戦略列と、その積分の強適合・càdlàgなucp極限の組であって、
積分列が全elementary testに一様なCauchy性を持つものをいう。
\begin{proposition}[利得領域と終端集合の同定]
同じ有界sourceの一般切断graphの利得領域は、elementary-test実現の利得領域と
区別不能性まで一致する。各候補の有限終端極限を明示して定めた
$a$-admissible終端集合も、両領域で一致する。
\end{proposition}
\begin{proof}
初等戦略列の各行の係数は有界なので、その積分は一般切断graphに属する。
そのÉmery Cauchy性とucp極限から、同じ極限へのÉmery収束を得る。
定理\ref{thm:graph-closure}を適用すると、elementary-test実現の一般graph帰属が従う。
逆に一般graphからは、各有界切断行を元のsourceのelementary積分で近似する。
停止を外す誤差は、全testに共通して停止時刻がhorizon以前となる確率で抑えられる。
切断行の収束と三角不等式により、正則代表元は任意の有限horizonでelementary近似を持つ。
整数horizon $n+1$ で誤差を $1/(n+1)$ 以下に選ぶ。
horizonの単調性から全horizonで収束し、全test共通のCauchy評価と
零始点性を用いたucp収束により、実現の全条件を満たす。
区別不能な利得の間で、各時刻の下界と候補自身の終端極限を移す。
有限horizon収束から終端極限の存在を推論する必要はない。
\end{proof}

\subsection{分解・市場操作の入力への接続}
\label{sec:auxiliary-transforms}
\begin{proof}[命題\ref{input:decomposition}の証明]
元市場の利得 $Y$ は定理\ref{thm:integral-input}のelementary近似によりgood integratorである。
その性質を $Q$ に移し、同じ利得と包絡 $|q|$ に定理\ref{thm:decomposition-input}を適用する。
$Y_0=0$ a.s.なので、得た零始点成分は同じ $Y$ のspecial分解となる。
\end{proof}

\begin{proof}[命題\ref{input:market}の証明]
\contractref{M1}--\contractref{M3}は\ref{sec:closure-proof}節の利得領域・終端集合の同定である。
$a=1$ と全ての正の $a$ に関する合併も同じ候補の終端証拠を保つ。
\contractref{M4}の可測性は、適合な利得の整数時刻での値の極限として終端を表すことによる。
同じ可算時刻での下界を共通の確率1の集合に集めて極限に渡せば、終端下界を得る。

\contractref{M5}の正則性・零始点性・固有終端は実現モデルのデータであり、
graph帰属は上の同定が与える。終端をa.e.等しい代表元に置き換えても、同じ利得がその終端へ収束する。
切替えの係数は、元の係数 $H,K$ を用いて
$H+1_E1_{(\tau,\infty)}(K-H)$ と書ける。
$E\in\F_\tau$ と有界停止時刻の条件により予測可能であり、
系\ref{cor:market-operations}の停止・制限・加法で指定した利得を得る。
$t\to\infty$ とすると、両利得の固有終端と有限の停止値から指定した終端公式が従う。

\contractref{M6}は定理\ref{thm:integral-input}の閉停止則を利得と両成分に同時に適用したものである。
局所BVの経路を有限 $T$ で停止すると全時間BVとなる。
\contractref{M7}では一般graphの線形性を有限和に適用する。
同じ係数でM/FV成分も加えれば、予測可能性、局所BV性、局所マルチンゲール性、零始点性を保つ。
従って成分を選び直すことなく、指定した重みの三つの有限和を同時に実現できる。
\end{proof}

以下の有限構成では、補助分解上に実現した積分を元価格へ戻す。
これは系\ref{cor:market-operations}による。同系は非有界な実現済み係数も切断で扱い、
下界と有限時刻後の定数性を保つ。同値測度間で移すのは一般graphであり、special性ではない。

\subsection{補助利得に対する定量評価}
\label{sec:ds-estimates}
第\ref{sec:components}節で使う数値評価を示す。
以下の成分は一つの確率 $Q$ の下での零始点正準成分である。
取引の帰還先の価格は常に $S$ とする。構成で使う操作は有界予測可能制限、加法、
正のスカラー倍、閉停止だけである。系\ref{cor:market-operations}によって、構成した許容利得は
全て元市場に実現される。議論は\cite[Lemmas 4.7--4.10]{DS}に従うが、
実現済み利得の族として述べるため、補助分解は添字ごとに異なってよい。

\begin{lemma}[具体的なjump包絡]\label{lem:explicit-jump-envelope}
$Y=M+A$ を予測可能FV成分を持つ正則special過程とし、
$0\leq q\in L^2(Q)$ により全時刻で $|Y_t|\leq q$ とする。
各有限 $T$ に対して非負の $\xi_T\in L^2(Q)$ が存在し、
\[
 |\Delta M_t|\leq\xi_T\quad(0<t\leq T),\qquad
 \|\xi_T\|_2\leq6\|q\|_2
\]
となる。
\end{lemma}
\begin{proof}
$Z_t=E_Q[q\mid\mathcal F_t]$ のcàdlàgマルチンゲール版を取り、
$Z_T^*=\sup_{t\leq T}|Z_t|$ と置く。
予測可能停止時刻 $\sigma\leq T$ では、正準分解のjumpの条件付き期待値から
\[
 \Delta A_\sigma=E_Q[\Delta Y_\sigma\mid\mathcal F_{\sigma-}],\qquad
 |\Delta A_\sigma|\leq2Z_{\sigma-}
\]
を得る。局所成分については先に局所化し、$\sigma$ を下から予告する。
$Y$ の可積分包絡により、条件付き期待値をjumpの極限へ移せる。
予測可能過程のjumpの列挙によって、$A$ の非零jumpは可算個のそのような時刻で覆われる。
共通の確率1の集合で全ての $0<t\leq T$ に対して $|\Delta A_t|\leq2Z_T^*$ となる。
従って
\[
 \xi_T=2q+2Z_T^*,\qquad
 \|\xi_T\|_2\leq2\|q\|_2+4\|q\|_2=6\|q\|_2
\]
と取れる。最後はDoobの $L^2$ 不等式である。
これは\cite[Corollary 2.4]{DS}のjump評価について、包絡とhorizonを明示したものである。
\end{proof}

\begin{lemma}[有限horizonの消失リスク戦略]\label{lem:quantitative-47}
元市場に実現された利得 $Y=M+A$ が $Y\geq-1$、$|Y|\leq q$、
$\|q\|_2\leq B$ を満たすとする。
$0<\alpha\leq1/4$、正整数 $n$、有限 $T>0$ に対し
\[
 Q\left(\sup_{t\leq T}|M_t|>n^3\right)>8\alpha
\]
とする。$k=\lfloor\alpha n/4\rfloor$ と置き、$k>0$、
$6B\leq n^2$、$(B/n)^2\leq\alpha$、
$m_n=\alpha^2-(4B/(\alpha n^2))^2\geq0$ と仮定する。
元市場の利得 $V_n$ で、$T$ 以後一定であり、
\[
 \begin{aligned}
 V_n(t)&\geq-n^{-1/4}-\frac{2}{nk}\quad(\forall t),\\
 Q\left(V_n(T)\geq\frac{\alpha m_n}{4}-n^{-1/4}\right)
 &\geq\frac{m_n}{2}-\frac{16\sqrt n}{k}
 \end{aligned}
\]
を満たすものが存在する。
\end{lemma}
\begin{proof}
最初の同時passageと再尺度化した利得を
\[
 \theta=T\wedge\inf\{t:|M_t|>n^3\text{ または }|Y_t|>n\},\qquad
 L=n^{-2}Y^\theta=N+B^L
\]
とする。停止前の利得上界と $Y\geq-1$ から、$\theta$ でのjumpを含めて
$\Delta L\geq-(n+1)/n^2\geq-2/n$ である。また $L^*\leq q/n^2$ である。
前補題により $N$ のjump包絡の $L^2$ ノルムは $6B/n^2\leq1$ 以下となる。
停止前には $|N|\leq n$ であるから、全時間の包絡は $n$ とこのjump包絡の和で抑えられる。
従って $N$ は真の $L^2$ マルチンゲールである。
Chebyshev不等式で利得上限に達する事象を除き、初回passageでの右連続性を使うと
\[
 Q(N_T^*>n/2)>8\alpha-(B/n)^2\geq7\alpha
\]
を得る。半径を $n/2$ とすることで、passageでの厳密不等号の問題を避ける。

$\kappa_0=0$ から始め、$\kappa_{j-1}$ 以後の $N$ の変位の絶対値が
初めて $1$ を厳密に超える時刻と $T$ の最小を $\kappa_j$ とする。
$1\leq j\leq k$ に対し $F_j=N_{\kappa_j}-N_{\kappa_{j-1}}$ と置く。
jump包絡より $\|F_j\|_2\leq2$、任意抽出定理より $E_QF_j=0$ である。
最初の $k$ 個のexcursionが完了していなければ、最大関数はこれらの停止excursionの
最大関数の和で抑えられる。Doob不等式はそれぞれの $L^2$ ノルムを $4$ で抑えるので、
\[
 Q(N_T^*>n/2,\ k\text{ 個のexcursionが未完了})
 \leq\left(\frac{4k}{n/2}\right)^2\leq\alpha.
\]
特に全ての $j\leq k$ に対し $Q(|F_j|\geq1)>6\alpha$ である。
従って $EF_j^-=E|F_j|/2>3\alpha$ となる。
$E_j=\{F_j^->\alpha\}$ とすると $E[F_j^-1_{E_j}]>2\alpha$ であり、
Cauchy--Schwarz不等式で同じ期待値は $2\sqrt{Q(E_j)}$ 以下である。
従って $Q(E_j)>\alpha^2$ となる。
対応する利得増分の $L^2$ ノルムは $2B/n^2$ 以下である。
水準 $\alpha/2$ のChebyshev不等式と $B^L=L-N$ から
\[
 Q(B^L_{\kappa_j}-B^L_{\kappa_{j-1}}\geq\alpha/2)\geq m_n
\]
を得る。

$dB^L$ の予測可能な正のHahn集合 $P$ を一つ選び、実際の利得を
\[
 R=(1_P1_{(0,\kappa_k]})\cdot L=N^R+C
\]
とする。$C$ は非減少で、各excursionでの増分は $B^L$ の対応する増分以上である。
上で得た $k$ 個の事象の成立個数を $Z$ とすると、$0\leq Z\leq k$、$EZ\geq km_n$ である。
$EZ\leq km_n/2+kQ(Z\geq km_n/2)$ より
\[
 Q(C_T\geq k\alpha m_n/4)\geq m_n/2
\]
となる。excursion事象の独立性は使わない。
予測可能制限は各マルチンゲールexcursionの $L^2$ ノルムを縮小する。
互いに交わらない停止区間の増分の直交性から $E|N^R_T|^2\leq4k$、従って
\[
 Q\left(\sup_{t\leq T}|N^R_t|>d\right)\leq16k/d^2
\]
を得る。ここでは全て固定horizonの真の $L^2$ マルチンゲールであり、
任意抽出と等長性がノルム縮小と直交性を正当化する。

最後に
\[
 d=kn^{-1/4},\qquad
 \lambda=T\wedge\inf\{t:R_t<-d\},\qquad V_n=k^{-1}R^\lambda
\]
と置く。正のHahn集合への制限は利得のjumpまたは零を選ぶから、$\Delta R\geq-2/n$ である。
従って停止時の飛び越しも含めて $R^\lambda\geq-d-2/n$ となる。
上のマルチンゲール最大事象の外では、$T$ まで $R=N^R+C\geq-d$、
かつ $R_T\geq C_T-d$ である。$C_T$ の事象と交差を取り $k$ で割れば、
$16k/d^2=16\sqrt n/k$ より両評価を得る。
$L,R,V_n$ を定義する全操作は本小節冒頭の元市場への実現を持ち、
最後の停止によって終端値は文字どおり $V_n(T)$ となる。
\end{proof}
この有限尺度の構成を用いるNFLVRの矛盾は、命題\ref{prop:indexed-maximal}で行う。

\begin{lemma}[凸尾部の許容性と具体的な包絡]\label{lem:quantitative-tails}
各利得が有限時刻以後一定である実現済み族 $Y_i=M_i+A_i\geq-1$ を考える。
共通包絡 $q\geq0$ は $L^2(Q)$ に属し、$\{M_i^*\}_i$ は確率有界と仮定する。
$\tau_i(c)=\inf\{t:|M_i(t)|>c\}$、$E_i(c)=\{\tau_i(c)<\infty\}$ と置く。
有限凸重み $w$ に対する尾部利得とそのM成分を
\[
 Z^{w,c}=\sum_iw_i(Y_i-Y_i^{\tau_i(c)}),\qquad
 L^{w,c}=\sum_iw_i(M_i-M_i^{\tau_i(c)})
\]
とする。利得の包絡は
\[
 G^{w,c}=2q\sum_iw_i1_{E_i(c)},\qquad
 \sup_w\|G^{w,c}\|_2\leq2\sup_i\|q1_{E_i(c)}\|_2\longrightarrow0
\]
で与えられる。各有限horizonでjump包絡 $\xi_T^{w,c}$ を持ち、
$\sup_w\|\xi_T^{w,c}\|_2\leq6\sup_w\|G^{w,c}\|_2\to0$ となる。
このノルム上界のcutoffは $T$ に依存しない。
さらに任意の $\varepsilon>0$ に対し $N\geq0$ を選べる。
各 $0<\delta\leq\varepsilon/4$ に対して、重みを選ぶ前に $c_0$ を定められ、
$c\geq c_0$ かつ $L^{w,c}$ の水準 $\varepsilon$ への有限passageの確率が
$\varepsilon$ を超えるなら、元市場の利得 $V=M^V+A^V$ と包絡 $g$ が存在して
\[
 V\geq-1,\qquad |V|\leq g,\qquad \|g\|_2\leq1,\qquad
 Q\!\left((M^V)^*>\frac{\varepsilon}{2(N+1)\delta}\right)>\frac\varepsilon2
\]
となる。この主張は一つの有限凸尾部の構成であり、NFLVRを仮定しない。

\end{lemma}
\begin{proof}
仮定した最大関数の確率有界性から $\sup_iQ(E_i(c))\to0$ である
（必要ならpassage水準を少し下げて適用する）。
$q^2$ の積分の絶対連続性により $\sup_i\|q1_{E_i(c)}\|_2\to0$ となる。
各尾部はpassage以前およびpassage時に零で、それ以後の絶対値は $2q$ 以下である。
包絡とそのノルム評価は三角不等式から従う。

有限尺度の構成を示す。$\varepsilon>0$ を固定する。
$Q(q\geq N)<\varepsilon/4$ となる $N$ を選び、全利得の絶対値のいずれかが
初めて $N$ を厳密に超える時刻を $\zeta$ とする。
$0<\delta\leq\varepsilon/4$ に対し、全ての $i$ で $Q(E_i(c))\leq\delta^2$、
全ての重みで $\|G^{w,c}\|_2\leq\delta$ となるcutoffを選ぶ。
これらは結論を破る重みを選ぶ前の選択である。
active mass
\[
 F_t=\sum_iw_i1_{\{\tau_i(c)<t\}},\qquad
 \gamma=\inf\{t:F_t>\delta\}
\]
は非減少、適合、左連続であるから $F_\gamma\leq\delta$ となる。
一方 $EF_\infty\leq\delta^2$ より $Q(\gamma<\infty)\leq\delta$ である。
$t\leq\zeta\wedge\gamma$ なら全てのactiveな開始時刻は $\zeta$ より厳密に前であるため、
$Y_i(\tau_i(c))\leq N$ である。また閉停止時を含めて $Y_i(t)\geq-1$ なので
\[
 Z^{w,c}_t\geq-(N+1)F_t\geq-(N+1)\delta
\]
となる。$L^{w,c}$ の水準 $\varepsilon$ へのpassageも停止に加える。
このpassageが有限である確率が $\varepsilon$ を超えるとする。
利得上限とactive massによる二つの停止が先に起こる事象を除いても、残る確率は
$\varepsilon-\varepsilon/4-\delta\geq\varepsilon/2$ より大きい。
その事象では閉停止時のM成分の絶対値が $\varepsilon$ 以上であり、
特に $\varepsilon/2$ を厳密に超える。有限和に現れる各利得は有限時刻後一定なので、
この停止利得も有限時刻後一定である。
$(N+1)\delta$ で正規化すれば、元市場の $1$ 許容利得を得る。
包絡ノルムは $1/(N+1)\leq1$ であり、表示した増幅後の確率下界が従う。
最後に補題\ref{lem:explicit-jump-envelope}を $q=G^{w,c}$ として適用する。
具体的には
\[
 \xi_T^{w,c}=2G^{w,c}+2\sup_{t\leq T}|E_Q[G^{w,c}\mid\mathcal F_t]|
\]
と取れる。ノルム評価はDoob不等式から従い、horizonに依存しない。
\end{proof}

\contractref{F1}は補題\ref{lem:quantitative-47}の終端 $f=V_n(T)$ を取り、
$K_0\subseteq K_0-L^0_+$ を使ったものである。\contractref{F2}は補題\ref{lem:quantitative-tails}の結論そのものである。

\subsection{停止前energyとマルチンゲールの近似}
\label{sec:approximation-inputs}
\begin{proof}[\contractref{F3}--\contractref{F5}の証明]
\contractref{F3}には、補題\ref{lem:quantitative-tails}の証明の包絡部分だけを使う。
この部分は利得の下界を使わない。凸尾部の利得包絡を $G^{w,c}$ とすると、
$\sup_w\|G^{w,c}\|_2\to0$ であり、補題\ref{lem:explicit-jump-envelope}から
そのM成分のjump包絡 $\xi_T^{w,c}$ を $\sup_w\|\xi_T^{w,c}\|_2\leq6\sup_w\|G^{w,c}\|_2$ と取れる。
$\eta>0$ に対して、この上界が $\eta$ 以下となるcutoffを先に選ぶ。
凸尾部 $L^{w,c}$ を水準 $\eta$ へのpassageと $T$ で停止すると、停止前の大きさと
jump包絡から停止過程 $R$ の全時間上限は $\eta+\xi_T^{w,c}$ 以下となる。
局所マルチンゲールの停止列をこの $L^2$ 包絡で外せるので、$R$ は真のL²マルチンゲールで、
$\|R_T\|_2\leq2\eta$ である。
定理\ref{thm:integral-input}のtest評価から $E_Q e_T^J(R,0)\leq4\eta$ を得る。
passageが有限でない経路では元のtailと同じなので、停止を外すcapped誤差は
$Q(B(w,c,\eta))$ 以下である。これが\contractref{F3}である。

\contractref{F4}でも同じ議論を元の $M$ の水準 $c$ で使う。
補題\ref{lem:explicit-jump-envelope}の $\|\xi_T\|_2\leq6\|q\|_2$ と
$\sup_t|M_{t\wedge\tau(c)\wedge T}|\leq c+\xi_T$ から、真のマルチンゲール性と
$\|P_T\|_2\leq c+6\|q\|_2$ を得る。
\contractref{F5}はcàdlàg経路の閉停止による右連続性と零始点性から従い、包絡を必要としない。
\end{proof}

\begin{proof}[命題\ref{input:approximation}の証明]
\contractref{A1}について、所定の誤差 $\varepsilon>0$ に対し、近似誤差を $\varepsilon/3$ とする列を選ぶ。
そのCauchy誤差も $\varepsilon/3$ とし、二つの近似誤差と中央の差の三角不等式を使う。
全ての上界がtestに一様なので、得られるcutoffもtestに依存しない。

\contractref{A2}では $P_n-P_k$ を $U$ で停止する。
定理\ref{thm:integral-input}のL² test評価と任意抽出・条件付き期待値の縮小性より、
全ての $T\leq U$ と $J$ に対して
\[
 E_Qe_T^J(P_n,P_k)\leq2\|P_n(U)-P_k(U)\|_2.
\]
終端のL² Cauchy性から、右辺を $\varepsilon$ 以下にする共通のcutoffを選ぶ。
\end{proof}

\subsection{有限Hahn取引の構成と閉停止評価}
\label{sec:finite-hahn-proof}
\begin{proposition}[単位有界completed変換と有限horizonのHahn利得]
通常の条件を満たす確率空間上で、$Y=M+A$ とし、$Y,M,A$ は強適合・右連続であり、
$A$ は予測可能な全時間有界変動過程とする。
$M$ は真のマルチンゲールで、固定した $T$ に対し $M_T\in L^2$ とする。
任意の単位有界elementary test $J$ に対して
\[
 \mathbb E\left[1\wedge\sup_{t\leq T}|(J\cdot M)_t|\right]\leq b
\]
が成り立つとする。単位有界な予測可能係数 $k$ の有限horizonのcompleted積分を
$I^k=(k\mathbf 1_{(0,T]})\cdot M$ とすると、同じ $b$ で全elementary testに対する
$I^k$ の評価が成り立つ。

さらに、$A$ の正のHahn集合 $B$ と有限horizonの積分利得
$R=\mathbf 1_{B\cap(0,T]}\cdot Y$ を構成できる。
同じ $R=N+C$ について、ほとんど確実に全ての $0\leq a\leq c\leq T$ で
\[
 0\leq C_c-C_a,\qquad A_c-A_a\leq C_c-C_a
\]
となり、$N$ の全elementary testの上界と
$\mathbb E[1\wedge\sup_{t\leq T}|N_t|]\leq b$ が成り立つ。
\end{proposition}
\begin{proof}
定理\ref{thm:integral-input}の第4項により、単位有界変換は同じ全test上界 $b$ を保つ。第5項のHahn集合を $[0,T]$ に制限し、第2項で同じ利得を $N+C$ と分解する。右端を含むStieltjes増分が二つの増分不等式を与える。零始点 $N$ にhorizonの単位testを適用すれば、capped最大関数にも上界 $b$ を得る。
\end{proof}

\begin{proposition}[局所積分の変換への評価移送と元価格への帰還]
通常の条件を満たす確率空間で、$Y=M+A$ を正則な局所special sourceとする。
積分表示の停止座標ではM成分は真のL²マルチンゲールであり、その停止列は全時間を尽くすとする。
固定した $T$ で $M$ の全単位elementary testのcapped期待値上界が $b$ であるとする。
$|k|\leq1$ のintrinsic 実現された積分 $k\cdot Y=N+C$ は、同じ $T$ で
$N$ の全単位test上界 $b$ を持つ。$N_0=0$ なら、そのcapped最大関数の期待値も $b$ 以下である。
$A,C$ には局所有限変動で十分であり、未停止 $M,N$ の大域L²性を仮定しない。

さらに、$Y$ が元の有界価格 $S$ に実現され、零始点の正規化された成分を持ち、
$|Y_t|\leq\xi$、$\xi\in L^2$ とする。
$Y^T$ の正のHahn集合を生成して作る積分利得 $V=N+C$ について、同時に
\[
 0\leq C_c-C_a,\qquad A_{c\wedge T}-A_{a\wedge T}\leq C_c-C_a
 \quad(0\leq a\leq c),
 \qquad
 \mathbb E[1\wedge\sup_{t\leq T}|N_t|]\leq b
\]
と全単位test上界が成り立ち、$V$ の一般積分graphを元の $S$ へ戻せる。
成分評価の測度が元測度と同値なら、一般積分graphのみを元測度へ移送する。
\end{proposition}
\begin{proof}
定理\ref{thm:integral-input}の局所版の変換評価を固定した $T$ で使う。局所化を除く際の誤差は停止が $T$ 以前となる確率であり、上界 $b$ はそのまま残る。Hahn成分は同定した同じFV積分を用い、前命題の増分評価を保つ。L²包絡を持つ元市場の利得の場合には、系\ref{cor:market-operations}を同じHahn利得へ適用する。移送するのは一般積分graphである。
\end{proof}

\paragraph{二つの指定利得への適用}
$R_n=\widetilde M_n+\widetilde A_n$、$R_k=\widetilde M_k+\widetilde A_k$ を、二つの指定されたspecial利得とする。共通包絡 $q\in L^2(Q)$、
下界 $-1$、固定 $T$ での両順序のM差の全test上界 $b$ を仮定する。
差の利得は $2q$ に支配され、指定した成分差をそのまま分解として持つ。
定理\ref{thm:decomposition-input}のL²停止座標と前命題をこの差に適用する。
得られる $V_{n,k}=N+C$ は零始点で、元価格の一般graphを持ち、
$C$ は非減少かつFV差の停止増分を支配し、$N$ の全testとcapped最大関数の上界は $b$ である。
ここで添字 $n,k$ は二利得を区別する記号にすぎず、列の選択、NFLVR、最大性は使わない。
以下の三命題をこの同じ取引に適用する。

\begin{proposition}[Hahn改善利得の閉じた有限停止と許容性]
上で構成した同じ利得を $V_{n,k}=N+C$、二つの入力を $R_n,R_k$ とする。
$\delta\geq0$ に対し
\[
 D_t=N_t-\max\{\widetilde M_n(t)-\widetilde M_k(t),0\},\qquad
 \tau=\inf\{t\geq0:D_t<-\delta\},\qquad \sigma=\tau\wedge T
\]
とおく。$\sigma$ は有限停止時刻であり、
$B=(R_k+V_{n,k})^\sigma$ は全時刻で $B\geq-(1+\delta)$ を満たす。
その固有終端 $B_\infty=(R_k+V_{n,k})_\sigma$ は、元市場の $K_{1+\delta}$ に属する。
\end{proposition}
\begin{proof}
$D$ は強適合かつ右連続なので、下方passageとその有限切断は停止時刻である。
成分は零始点であり、$t\leq T$ では
$C_t\geq0$、$C_t\geq\widetilde A_n(t)-\widetilde A_k(t)$ である。
$t<\sigma$ なら $D_t\geq-\delta$ なので
\[
 (R_k+V_{n,k})_t\geq\max\{R_n(t),R_k(t)\}-\delta.
\]
従って停止直前にも、左極限を取った同じ不等式が成り立つ。
実際のHahn制限のジャンプは、正のHahn集合では $\Delta(R_n-R_k)$、
その補集合では零である。この恒等式は $t=T$ を含む。
従って $R_k+V_{n,k}$ の停止時のジャンプは $\Delta R_n$ または $\Delta R_k$ に一致する。
左極限の不等式に選ばれた入力のジャンプを加え、その入力の下界 $-1$ を使えば
$B_\sigma\geq-(1+\delta)$ となる。$\sigma=0$ では零始点性を使う。

二つの一般積分graphを元価格の実現へ戻して加え、元市場の一般graphの有限停止則を適用する。
$\sigma\leq T$ により $T$ 後は一定であり、上の下界と合わせて、同じ有限終端の
$K_{1+\delta}$ への帰属を得る。成分評価は $Q$ に保ち、一般graphと下界を元測度へ戻す。
\end{proof}

\begin{proposition}[成功事象での元利得の継続と失敗確率]
前命題の $\tau,\sigma$ と $E=\{T<\tau\}$ を使い、任意の有限 $U\geq T$ に対して
\[
 Z_t=(R_k+V_{n,k})_{t\wedge\sigma}
       +1_E\bigl(R_k(t\wedge U)-R_k(t\wedge T)\bigr)
\]
とおく。$Z$ は元市場の一般積分graphを持ち、全時刻で $Z\geq-(1+\delta)$ であり、
$Z_U\in K_{1+\delta}$ となる。$E$ 上では
\[
 Z_U=R_k(U)+V_{n,k}(T).
\]
また、$N$ と $\widetilde M_n-\widetilde M_k$ の $[0,T]$ でのcapped最大関数の期待値が
ともに $b$ 以下なら、$0<\delta\leq1$ に対し
\[
 Q(\tau\leq T)\leq\frac{8b}{\delta}.
\]
この確率評価は $U$ に依存しない。
\end{proposition}
\begin{proof}
$E\in\mathcal F_T$ なので、元市場の実現のevent-tail制限を $T$ から適用できる。
$T$ までの停止利得と、この取引を $U$ で停止した利得を加え、同じ $Z$ の一般graphを得る。
$E^c$ 上では閉じた停止の下界を使う。$E$ 上では $\sigma=T$ であり、$T$ 以前も同じ下界を持つ。
$T$ 以後は $Z_t=R_k(t\wedge U)+V_{n,k}(T)$ である。
$D_T\geq-\delta$ と $C_T\geq0$ から $V_{n,k}(T)\geq-\delta$ なので下界を保つ。
$U$ 後は一定であり、固有終端の帰属も従う。

二つのcapped最大関数がともに $\delta/4$ 未満なら、$\delta/4<1$ により
$|N_t|,|\widetilde M_n(t)-\widetilde M_k(t)|<\delta/4$ が全 $t\leq T$ で成り立つ。
従って $D_t\geq-\delta/2$ である。$T$ での右連続性はこの厳密な余裕を
$T$ より少し先まで保つので、$\tau>T$ となる。よって $\{\tau\leq T\}$ は、二つのcapped最大関数の
いずれかが $\delta/4$ 以上となる事象の和に含まれる。Markov不等式と和の評価から結論を得る。
元のM差の期待値上界は、零始点性とhorizonの単位testにより、仮定した有限尺度のtest評価から得る。
\end{proof}

\begin{proposition}[両向きのHahn増分による変動支配と向きの選択]
指定した二つの利得の成分差 $A=\widetilde A_n-\widetilde A_k$ と逆向きの差 $-A$ に、
上の積分構成をそれぞれ適用する。そのFV成分を $C^+,C^-$ とすると、ほとんど確実に
\[
 \operatorname{Var}_{[0,T]}(A)\leq C^+_T+C^-_T.
\]
従って $Q(\operatorname{Var}_{[0,T]}(A)>\alpha)>\alpha$ なら、
少なくとも一方の符号について
\[
 Q(C^{\pm}_T>\alpha/2)>\alpha/2
\]
となる。各向きの利得は同じ構成から得たM上界、downside失敗確率、
任意の有限 $U\geq T$ での元市場の終端帰属を保持する。
\end{proposition}
\begin{proof}
$C^+,C^-$ は零始点の単調過程である。$0\leq a\leq c\leq T$ に対して
\[
 A_c-A_a\leq C^+_c-C^+_a,\qquad
 -(A_c-A_a)\leq C^-_c-C^-_a.
\]
よって $|A_c-A_a|\leq(C^+_c-C^+_a)+(C^-_c-C^-_a)$ である。
任意の有限分割に沿って和を取り、上限を取ると変動支配を得る。
局所BVの経路を $T$ で停止したStieltjes測度の全変動は、
右連続性によりこの区間の経路変動と一致するため、成分Cauchy消費側と同じ変動量を評価している。
大きな変動の事象は $\{C^+_T>\alpha/2\}\cup\{C^-_T>\alpha/2\}$ に零集合を除いて含まれる。
確率の劣加法性から向きの選択が従う。
両順序のtest上界 $b$ を仮定しているので、両向きの構成は同時に使える。
\end{proof}

これらの構成は\contractref{H1}の全結論を同じ二つの利得について供給する。

\subsection{成分評価からの終端実現}
\label{sec:component-realization}
\begin{lemma}[成分評価からÉmery Cauchy性へ]
有限 $T$ について
\[
 U^J_{n,k}=1\wedge\sup_{t\le T}|(J\cdot(\widetilde M_n-\widetilde M_k))_t|,
 \qquad V_{n,k}=\Var_{[0,T]}(\widetilde A_n-\widetilde A_k)
\]
と置く。任意の $\delta>0$ について
\[
 \begin{aligned}
 \sup_J Q_*(U^J_{n,k}>\delta)&\longrightarrow0,\\
 Q_*(V_{n,k}>\delta)&\longrightarrow0
 \end{aligned}
 \qquad(n,k\longrightarrow\infty).
\]
が成り立つなら、$R_n$ はÉmery Cauchyである。ここで二つの添字は独立に無限大へ送る。
すなわち、各確率誤差に対する閾値は $n,k$ と test $J$ を選ぶ前に固定する。
\end{lemma}
\begin{proof}
Stieltjes積分の変動評価より、FV差のtest積分の最大値は $V_{n,k}$ 以下である。
利得差のcapped test誤差を $Z^J_{n,k}$ と書けば、$0\leq Z^J_{n,k}\leq1$ かつ
$Z^J_{n,k}\leq U^J_{n,k}+V_{n,k}$ である。従って $0<\delta<1/2$ について
\[
 E_{Q_*}Z^J_{n,k}\leq2\delta+Q_*(U^J_{n,k}>\delta)+Q_*(V_{n,k}>\delta).
\]
$\delta$ と二つの確率上界を先に小さく選ぶと、全ての $J$ に共通のCauchy閾値が得られる。
M成分の期待値評価から確率評価への移行はMarkov不等式である。
\end{proof}

\begin{lemma}[同じ一様極限へのÉmery収束]
$R_n$ がÉmery Cauchyで、正則過程 $X$ へucp収束するなら、$R_n\to X$ はÉmery収束である。
\end{lemma}
\begin{proof}
一つのtestを固定すると有限和の増分積分はucp極限と両立し、capped誤差の期待値も零へ収束する。
$\varepsilon$ に対するCauchy閾値 $N$ を全test共通に取る。
$n,k\geq N$ に対し三角不等式を適用し、固定した $n,J$ について $k\to\infty$ とすれば、
$R_n$ と $X$ の期待値誤差は $\varepsilon$ 以下となる。最後に全testにわたる上限を取る。
\end{proof}


\begin{proof}[命題\ref{input:realization}の証明]
同じ $Y_n=M_n+A_n$ に上の二補題を適用する。各確率誤差を小さくするcutoffは全testに共通である。
従って $Y_n\to X$ は $Q$ の下でÉmery収束する。
定理\ref{thm:integral-input}の同値測度移送で収束を元測度へ戻し、
定理\ref{thm:graph-closure}で同じ $X$ を元価格の一般積分に実現する。
指定された $X\geq-1$ と $X_\infty=h$ により $h\in K_1$ である。
\end{proof}

